% Géométrie analytique du plan — Mathématiques 1ère E — Cours signé OnBuch+
% Programme officiel MINESEC. Compiler avec : tectonic N14-geometrie-analytique-plan.tex
\newcommand{\DOCMATIERE}{Mathématiques}
\newcommand{\DOCNIVEAU}{1ère E}
\newcommand{\DOCMODULE}{Configurations et transformations élémentaires du plan}
\newcommand{\DOCLECON}{N14}
\newcommand{\DOCTITRE}{Géométrie analytique du plan}
% OnBuch+ — préambule « cours signés » ENRICHI (compilé avec tectonic / XeLaTeX)
% Système visuel « Pop » d'OnBuch+ : fond crème (jamais blanc), bordures encre
% épaisses, ombres « dures » décalées, titres Archivo Black en capitales.
% Version MATHÉMATIQUES : graphes (pgfplots/tikz), boîtes pédagogiques
% (définition, propriété, méthode, attention, le savais-tu, exercice résolu,
% exercice, TP, bilan), page de garde et signature OnBuch+ ancrée en bas.
%
% Macros attendues dans main.tex AVANT \input{preamble} :
%   \DOCMATIERE  \DOCNIVEAU  \DOCMODULE  \DOCLECON (numéro)  \DOCTITRE
\documentclass[11pt]{article}

\usepackage{fontspec}
\usepackage[french]{babel}
\usepackage[a4paper,margin=2cm,top=3.4cm,bottom=2.8cm,headheight=1.4cm,footskip=1.3cm]{geometry}
\usepackage{amsmath,amssymb,amsfonts}
\usepackage{mathtools} % \xmapsto, \coloneqq... souvent utilisés par les modèles
\usepackage{cancel} % simplifications barrées
% \abs{z} (module, valeur absolue) et \norm{u} (norme), utilisés par les modèles
\providecommand{\abs}{}\let\abs\relax\DeclarePairedDelimiter{\abs}{\lvert}{\rvert}
\providecommand{\norm}{}\let\norm\relax\DeclarePairedDelimiter{\norm}{\lVert}{\rVert}
\usepackage[table]{xcolor} % [table] : fournit aussi \rowcolors (alternance de lignes)
\usepackage{pagecolor}
\usepackage{tikz}
\usetikzlibrary{arrows.meta,positioning,calc,shapes.geometric,shapes.misc,shapes.arrows,decorations.pathmorphing,decorations.pathreplacing,decorations.markings,patterns,shadows,mindmap,trees,fit,backgrounds,matrix,intersections,angles,quotes}
\usepackage{pgfplots}
\usepackage{pgf-pie} % diagrammes circulaires \pie (statistiques)
\pgfplotsset{compat=1.18}
\usepgfplotslibrary{fillbetween,groupplots} % aires entre courbes (calcul intégral)
\usepackage{enumitem}
\usepackage{fancyhdr}
% [most] charge toutes les bibliothèques tcolorbox (listings, minted, raster,
% documentation...) et gonfle inutilement la mémoire TeX sur un document long
% (~300 boîtes) : on ne charge que ce qui est réellement utilisé.
\usepackage[skins,breakable]{tcolorbox}
\usepackage{graphicx}
\usepackage{colortbl}
\usepackage{booktabs}
\usepackage{tabularx}
\usepackage{diagbox} % case d'en-tête diagonale des tableaux à double entrée
% Colonnes extensibles centrée / gauche / droite (C, L, R), souvent utilisées par les modèles
\newcolumntype{C}{>{\centering\arraybackslash}X}
\newcolumntype{L}{>{\raggedright\arraybackslash}X}
\newcolumntype{R}{>{\raggedleft\arraybackslash}X}
\usepackage{array}
\usepackage{multicol}
\usepackage{siunitx}
% parse-numbers=false : accepte aussi \SI{2,5 \times 10^{-3}}{...}
\sisetup{locale=FR,output-decimal-marker={,},group-minimum-digits=4,parse-numbers=false}
\DeclareSIUnit\ohm{\ensuremath{\symup{\Omega}}} % U+2126 absent de la police
\usepackage{qrcode}
% \degree : commande fréquemment utilisée par les modèles pour un angle ou une
% température en dehors de \SI{}{} (non fournie par les paquets chargés).
\providecommand{\degree}{^{\circ}}
% \qed (fin de démonstration), utilisé par les modèles sans amsthm
\providecommand{\qed}{\hfill$\square$}
% \barre : double barre des valeurs interdites dans un tableau de variations (tkz-tab)
\providecommand{\barre}{\ensuremath{\|}}
% environnement proof (amsthm non chargé) utilisé par les modèles
\ifdefined\proof\else\newenvironment{proof}[1][Démonstration]{\par\noindent\textit{#1.}\ }{\qed\par}\fi
\usepackage{lastpage}
\usepackage{hyperref}
\hypersetup{hidelinks}

% ============================================================
% Palette « Pop » OnBuch+ — mêmes hex que la classe P (plus_ui.dart)
% ============================================================
\definecolor{poporange}{HTML}{F59321}
\definecolor{popdark}{HTML}{E07A0C}
\definecolor{popcream}{HTML}{FFF6E8}
\definecolor{popink}{HTML}{1C1714}
\definecolor{poppurple}{HTML}{7A5AE0}
\definecolor{popgreen}{HTML}{2E9E6A}
\definecolor{poppink}{HTML}{E0457B}
\definecolor{popblue}{HTML}{2F7DE1}
\definecolor{popgold}{HTML}{C98A12}
\definecolor{popmuted}{HTML}{8A7F76}
\definecolor{popline}{HTML}{EADFD2}
\definecolor{popgray}{HTML}{8A7F76}
\colorlet{popgrey}{popgray}
% Alias tolérants (le modèle nomme parfois les couleurs différemment)
\colorlet{popwhite}{white}
\colorlet{popblack}{popink}
\colorlet{popred}{poppink}
\colorlet{popyellow}{popgold}
% Teintes claires pour fonds de tableaux / zones
\colorlet{popblueL}{popblue!10!white}
\colorlet{poporangeL}{poporange!14!white}
\colorlet{popgreenL}{popgreen!12!white}
\colorlet{poppurpleL}{poppurple!12!white}
\colorlet{poppinkL}{poppink!12!white}
\colorlet{popgoldL}{popgold!14!white}

\pagecolor{popcream}

% ============================================================
% Typographie
% ============================================================
\defaultfontfeatures{Ligatures=TeX}
\setmainfont{PlusJakartaSans-Medium.ttf}[Path=../../../fonts/,BoldFont=PlusJakartaSans-Bold.ttf]
% Maths en sans-serif (Fira Math) avec chiffres et lettres droites de Plus
% Jakarta, pour que les formules se fondent dans le texte.
\usepackage[math-style=ISO,bold-style=ISO]{unicode-math}
\setmathfont{FiraMath-Regular.otf}[Path=../../../fonts/]
\setmathfont{PlusJakartaSans-Medium.ttf}[Path=../../../fonts/,range={up/num,up/{Latin,latin},\mathup}]
\usepackage{icomma} % virgule décimale sans espace en mode math
% Caractères mathématiques Unicode absents des polices de texte (Plus Jakarta,
% Archivo Black) : rendus par leur équivalent mathématique, en texte comme en
% formule — sinon ils s'affichent en carré vide (ex. « Arithmétique dans ℤ »).
\usepackage{newunicodechar}
\newunicodechar{ℝ}{\ensuremath{\mathbb{R}}}
\newunicodechar{ℤ}{\ensuremath{\mathbb{Z}}}
\newunicodechar{ℂ}{\ensuremath{\mathbb{C}}}
\newunicodechar{ℕ}{\ensuremath{\mathbb{N}}}
\newunicodechar{ℚ}{\ensuremath{\mathbb{Q}}}
\newunicodechar{𝔻}{\ensuremath{\mathbb{D}}}
\newunicodechar{α}{\ensuremath{\alpha}}
\newunicodechar{β}{\ensuremath{\beta}}
\newunicodechar{γ}{\ensuremath{\gamma}}
\newunicodechar{δ}{\ensuremath{\delta}}
\newunicodechar{ε}{\ensuremath{\varepsilon}}
\newunicodechar{θ}{\ensuremath{\theta}}
\newunicodechar{λ}{\ensuremath{\lambda}}
\newunicodechar{μ}{\ensuremath{\mu}}
\newunicodechar{σ}{\ensuremath{\sigma}}
\newunicodechar{τ}{\ensuremath{\tau}}
\newunicodechar{φ}{\ensuremath{\varphi}}
\newunicodechar{ψ}{\ensuremath{\psi}}
\newunicodechar{ω}{\ensuremath{\omega}}
\newunicodechar{Σ}{\ensuremath{\Sigma}}
% majuscules produites par \MakeUppercase dans les titres (α-aminés -> Α)
\newunicodechar{Α}{\ensuremath{\alpha}}
\newunicodechar{Β}{\ensuremath{\beta}}
\newunicodechar{Γ}{\ensuremath{\Gamma}}
\newunicodechar{Δ}{\ensuremath{\Delta}}
\newunicodechar{Ε}{\ensuremath{\varepsilon}}
\newunicodechar{Θ}{\ensuremath{\Theta}}
\newunicodechar{Λ}{\ensuremath{\Lambda}}
\newunicodechar{Μ}{\ensuremath{\mu}}
\newunicodechar{Τ}{\ensuremath{\tau}}
\newunicodechar{Φ}{\ensuremath{\Phi}}
\newunicodechar{Ψ}{\ensuremath{\Psi}}
\newunicodechar{Ω}{\ensuremath{\Omega}}
\newunicodechar{↦}{\ensuremath{\mapsto}}
\newunicodechar{∈}{\ensuremath{\in}}
\newunicodechar{∉}{\ensuremath{\notin}}
\newunicodechar{∧}{\ensuremath{\wedge}}
\newunicodechar{⇔}{\ensuremath{\Leftrightarrow}}
\newunicodechar{⟺}{\ensuremath{\Longleftrightarrow}}
\newunicodechar{⇒}{\ensuremath{\Rightarrow}}
\newunicodechar{⟹}{\ensuremath{\Longrightarrow}}
\newunicodechar{∘}{\ensuremath{\circ}}
\newunicodechar{∪}{\ensuremath{\cup}}
\newunicodechar{∩}{\ensuremath{\cap}}
\newunicodechar{≡}{\ensuremath{\equiv}}
\newunicodechar{⊂}{\ensuremath{\subset}}
\newunicodechar{⊥}{\ensuremath{\perp}}
\newunicodechar{∃}{\ensuremath{\exists}}
\newunicodechar{∀}{\ensuremath{\forall}}
\newunicodechar{∛}{\ensuremath{\sqrt[3]{\,}}}
\newunicodechar{∜}{\ensuremath{\sqrt[4]{\,}}}
\newunicodechar{⃗}{\ensuremath{{}^{\to}}}
\newunicodechar{ⁿ}{\ifmmode^{n}\else\textsuperscript{\ensuremath{n}}\fi}
\newunicodechar{ˣ}{\ifmmode^{x}\else\textsuperscript{\ensuremath{x}}\fi}
\newunicodechar{ᵇ}{\ifmmode^{b}\else\textsuperscript{\ensuremath{b}}\fi}
\newunicodechar{ʳ}{\ifmmode^{r}\else\textsuperscript{\ensuremath{r}}\fi}
\newunicodechar{ᵘ}{\ifmmode^{u}\else\textsuperscript{\ensuremath{u}}\fi}
\newunicodechar{ʸ}{\ifmmode^{y}\else\textsuperscript{\ensuremath{y}}\fi}
\newunicodechar{ᵃ}{\ifmmode^{a}\else\textsuperscript{\ensuremath{a}}\fi}
\newunicodechar{ᵉ}{\ifmmode^{e}\else\textsuperscript{\ensuremath{e}}\fi}
\newunicodechar{ᵏ}{\ifmmode^{k}\else\textsuperscript{\ensuremath{k}}\fi}
\newunicodechar{ᵖ}{\ifmmode^{p}\else\textsuperscript{\ensuremath{p}}\fi}
\newunicodechar{ᴺ}{\ifmmode^{N}\else\textsuperscript{\ensuremath{N}}\fi}
\newunicodechar{ᶜ}{\ifmmode^{c}\else\textsuperscript{\ensuremath{c}}\fi}
\newunicodechar{ʰ}{\ifmmode^{h}\else\textsuperscript{\ensuremath{h}}\fi}
\newunicodechar{ᵠ}{\ifmmode^{q}\else\textsuperscript{\ensuremath{q}}\fi}
\newunicodechar{ₙ}{\ifmmode_{n}\else\textsubscript{\ensuremath{n}}\fi}
\newunicodechar{ₐ}{\ifmmode_{a}\else\textsubscript{\ensuremath{a}}\fi}
\newunicodechar{ᵢ}{\ifmmode_{i}\else\textsubscript{\ensuremath{i}}\fi}
\newunicodechar{ᵣ}{\ifmmode_{r}\else\textsubscript{\ensuremath{r}}\fi}
\newunicodechar{ₖ}{\ifmmode_{k}\else\textsubscript{\ensuremath{k}}\fi}
\newunicodechar{ᵦ}{\ifmmode_{\beta}\else\textsubscript{\ensuremath{\beta}}\fi}
\newunicodechar{ₓ}{\ifmmode_{x}\else\textsubscript{\ensuremath{x}}\fi}
\newfontfamily\popdisplay{ArchivoBlack-Regular.ttf}[Path=../../../fonts/]
\newfontfamily\poplogo{SpaceGrotesk-Bold.ttf}[Path=../../../fonts/]
\newfontfamily\popsemi{PlusJakartaSans-SemiBold.ttf}[Path=../../../fonts/]
\newfontfamily\popxbold{PlusJakartaSans-ExtraBold.ttf}[Path=../../../fonts/]
\newfontfamily\popxboldls{PlusJakartaSans-ExtraBold.ttf}[Path=../../../fonts/,LetterSpace=3.0]

\setlist[enumerate]{leftmargin=1.7em,itemsep=5pt,topsep=4pt}
\setlist[itemize]{leftmargin=1.7em,itemsep=4pt,topsep=4pt,label={\color{poporange}\textbullet}}
\setlength{\parindent}{0pt}
\setlength{\parskip}{5pt}
\renewcommand{\arraystretch}{1.25}

% pgfplots : style Pop par défaut
\pgfplotsset{
  every axis/.append style={axis line style={popink,line width=1pt},tick style={popink},
    grid style={popline,line width=0.5pt},label style={font=\small},tick label style={font=\footnotesize}},
  popcurve/.style={poporange,line width=1.8pt,no markers,smooth},
}
% Même style disponible dans un \draw TikZ simple (hors pgfplots)
\tikzset{popcurve/.style={poporange,line width=1.8pt}}

% ============================================================
% Pill / squiggle
% ============================================================
\newcommand{\poppill}[3]{%
  \tikz[baseline=-0.55ex]\node[draw=popink,line width=1.2pt,rounded corners=2.6pt,%
    fill=#2,inner xsep=6.5pt,inner ysep=3pt]{%
    \popxboldls\fontsize{7}{7}\selectfont\color{#3}\MakeUppercase{#1}};%
}
\newcommand{\popsquiggle}[1]{%
  \tikz[baseline=-0.4ex]{
    \draw[#1,line width=2.6pt,line cap=round]
      plot [smooth] coordinates {(0,0.09) (0.34,0.01) (0.68,0.21) (1.02,0.01) (1.36,0.10)};
  }%
}
% Mot-clé surligné (marqueur orange sous le texte)
\newcommand{\cle}[1]{{\popxbold\color{popdark}#1}}

% ============================================================
% En-tête / pied
% ============================================================
\pagestyle{fancy}
\fancyhf{}
\renewcommand{\headrulewidth}{0pt}
\renewcommand{\footrulewidth}{0pt}
\lhead{\raisebox{-0.15ex}{\poplogo\fontsize{13}{13}\selectfont\color{popink}OnBuch\color{poporange}+}}
\rhead{\poppill{\DOCMATIERE\ \textbullet\ \DOCNIVEAU\ \textbullet\ Leçon \DOCLECON}{white}{popink}}
\lfoot{\popxbold\fontsize{7.6}{7.6}\selectfont\color{popmuted}\MakeUppercase{Cours signé OnBuch+}\ \textbullet\ {\popsemi\DOCTITRE}}
\rfoot{\tikz[baseline=-0.6ex]\node[rounded corners=8.5pt,fill=popink,minimum height=17pt,minimum width=17pt,inner xsep=5pt,inner sep=0]{\popxbold\fontsize{8}{8}\selectfont\color{popcream}\thepage\,/\,\pageref*{LastPage}};}

% ============================================================
% Titres
% ============================================================
\newcommand{\chaptitre}[2]{%
  \begin{tcolorbox}[enhanced,colback=poporange,colframe=popink,boxrule=2.6pt,arc=11pt,
      drop shadow=popink,left=15pt,right=15pt,top=12pt,bottom=11pt,before skip=3pt,after skip=18pt]
    \poppill{#2}{popink}{popcream}\par\vspace{9pt}
    {\raggedright\hyphenpenalty=10000\exhyphenpenalty=10000\popdisplay\fontsize{22}{24}\selectfont\color{popcream}\MakeUppercase{#1}\par}
  \end{tcolorbox}
}
% Section numérotée : pastille orange avec le numéro + titre Archivo
\newcounter{popsec}
\newcommand{\coursec}[1]{%
  \stepcounter{popsec}\setcounter{popsub}{0}%
  \par\vspace{16pt}\noindent
  \begin{minipage}{\linewidth}
  \tikz[baseline=-0.7ex]\node[draw=popink,line width=1.6pt,fill=poporange,rounded corners=4pt,minimum size=22pt,inner sep=0]{\popdisplay\fontsize{12}{12}\selectfont\color{popcream}\thepopsec};%
  \hspace{8pt}{\popdisplay\fontsize{15}{17}\selectfont\color{popink}\MakeUppercase{#1}}\par
  \vspace{4pt}\noindent{\color{poporange}\rule{36pt}{3.6pt}}
  \end{minipage}\par\nopagebreak\vspace{8pt}
}
\newcounter{popsub}
\newcommand{\courssub}[1]{%
  \stepcounter{popsub}%
  \par\vspace{9pt}\noindent{\popxbold\fontsize{11.5}{13}\selectfont\color{popink}{\color{poporange}\thepopsec.\thepopsub}\hspace{6pt}#1}\par\nopagebreak\vspace{4pt}
}

% ============================================================
% Boîtes pédagogiques — même recette : fond blanc, bordure épaisse colorée,
% ombre dure encre, badge plein. Titre optionnel [#1].
% ============================================================
\tcbset{popbase/.style={breakable,enhanced,colback=white,boxrule=2.3pt,arc=9pt,
  shadow={3pt}{-3pt}{0pt}{popink},left=14pt,right=14pt,top=11pt,bottom=11pt,before skip=11pt,after skip=15pt,
  fontupper=\popsemi\fontsize{10.6}{14.5}\selectfont\color{popink},
  fontlower=\popsemi\fontsize{10.6}{14.5}\selectfont\color{popink}}}
% Coche dessinée (le glyphe ✓ n'existe pas dans les polices du document)
\newcommand{\popcheck}[1]{\tikz[baseline=-0.5ex]\draw[#1,line width=1.6pt,line cap=round,line join=round](-0.13,0.0)--(-0.04,-0.09)--(0.13,0.1);}
\newcommand{\popboxhead}[3]{\poppill{#1}{#2}{white}\ifx\relax#3\relax\else\hspace{6pt}{\popxbold\fontsize{10.6}{12}\selectfont\color{popink}#3}\fi\par\vspace{7pt}}

\newtcolorbox{aretenir}[1][]{popbase,colframe=popgreen,before upper={\popboxhead{À retenir}{popgreen}{#1}}}
\newtcolorbox{exemplebox}[1][]{popbase,colframe=poporange,before upper={\popboxhead{Exemple}{poporange}{#1}}}
\newtcolorbox{definition}[1][]{popbase,colframe=popblue,before upper={\popboxhead{Définition}{popblue}{#1}}}
\newtcolorbox{propriete}[1][]{popbase,colframe=poppurple,before upper={\popboxhead{Propriété}{poppurple}{#1}}}
\newtcolorbox{methode}[1][]{popbase,colframe=popgold,before upper={\popboxhead{Méthode}{popgold}{#1}}}
\newtcolorbox{attention}[1][]{popbase,colframe=poppink,before upper={\popboxhead{Attention}{poppink}{#1}}}
\newtcolorbox{experience}[1][]{popbase,colframe=popblue,colback=popblueL,before upper={\popboxhead{Expérience}{popblue}{#1}}}
% « Le savais-tu ? » : carte violette pleine (contexte, culture, Cameroun)
\newtcolorbox{savaistu}[1][]{popbase,colback=poppurple,colframe=popink,
  fontupper=\popsemi\fontsize{10.4}{14.2}\selectfont\color{white},
  before upper={\poppill{Le savais-tu\,?}{white}{poppurple}\ifx\relax#1\relax\else\hspace{6pt}{\popxbold\color{white}#1}\fi\par\vspace{7pt}}}
% Exercice résolu : énoncé en haut, \tcblower puis la solution sur fond crème
\newcounter{popexo}\newcounter{popexr}
\newtcolorbox{exoresolu}[1][]{popbase,colframe=poporange,colbacklower=popcream,
  segmentation style={popink,line width=1.2pt,dashed},
  before upper={\stepcounter{popexr}\popboxhead{Exercice résolu \thepopexr}{poporange}{#1}},
  before lower={\poppill{Solution}{popink}{popcream}\par\vspace{6pt}}}
% Exercice (non résolu en place) — la correction est dans la partie Corrigés
\newtcolorbox{exercice}[1][]{popbase,colframe=popink,
  before upper={\stepcounter{popexo}\popboxhead{Exercice \thepopexo}{popink}{#1}}}
% Corrigé
\newtcolorbox{corrige}[1][]{popbase,colframe=popgreen,colback=popgreenL,
  before upper={\popboxhead{Corrigé}{popgreen}{#1}}}
% Objectifs / prérequis / bilan
\newtcolorbox{objectifs}{popbase,colframe=popink,colback=poporangeL,
  before upper={\popboxhead{Objectifs de la leçon}{popink}{}}}
\newtcolorbox{prerequis}{popbase,colframe=popink,colback=popblueL,
  before upper={\popboxhead{Prérequis}{popblue}{}}}
\newtcolorbox{bilan}{popbase,colframe=popink,colback=white,boxrule=2.8pt,
  before upper={\popboxhead{Fiche bilan}{poporange}{L'essentiel à connaître pour l'examen}}}
% Figure encadrée avec légende
\newtcolorbox{popfigure}[1][]{enhanced,breakable=false,colback=white,colframe=popline,boxrule=1.4pt,arc=8pt,
  left=10pt,right=10pt,top=10pt,bottom=8pt,before skip=10pt,after skip=12pt,halign=center,
  lower separated=false,halign lower=center,
  fontlower=\popsemi\fontsize{9}{11.5}\selectfont\color{popmuted},
  #1}
\newcommand{\legende}[1]{\par\vspace{6pt}{\popsemi\fontsize{9}{11.5}\selectfont\color{popmuted}#1}}

% ============================================================
% Signature OnBuch+
% ============================================================
\newcommand{\poplogoinline}[1]{{\poplogo\fontsize{#1}{#1}\selectfont\color{popink}OnBuch\color{poporange}+}}
\newcommand{\signatureonbuch}{%
  \begin{tcolorbox}[enhanced,colback=popink,colframe=popink,boxrule=2.2pt,arc=10pt,
      drop shadow=poporange,left=14pt,right=14pt,top=10pt,bottom=10pt,before skip=0pt,after skip=0pt]
    \tikz[baseline=-0.7ex]\node[circle,fill=popgreen,draw=popcream,line width=1.3pt,minimum size=22pt,inner sep=0]{\popcheck{white}};%
    \hspace{9pt}\begin{minipage}[c]{0.62\linewidth}
      {\popxbold\fontsize{12}{13}\selectfont\color{popcream}Signé\ \textbullet\ {\poplogo OnBuch\color{poporange}+}}\par\vspace{2pt}
      {\raggedright\popsemi\fontsize{8}{10.5}\selectfont\color{popline}\MakeUppercase{Équipe pédagogique OnBuch+}\\\MakeUppercase{Conforme au programme officiel MINESEC}\par}
    \end{minipage}\hfill
    {\poplogo\fontsize{22}{22}\selectfont\color{popcream}OnBuch\color{poporange}+}
  \end{tcolorbox}%
}

% ============================================================
% Page de garde
% #1 titre · #2 matière · #3 niveau · #4 module · #5 n° leçon · #6 durée ·
% #7 description / objectifs (texte ou itemize)
% ============================================================
\newcommand{\pagegarde}[7]{%
  \thispagestyle{empty}
  \newgeometry{a4paper,margin=1.7cm,top=1.6cm,bottom=1.4cm}%
  \begin{tikzpicture}[remember picture,overlay]
    \fill[poporange!18] ([xshift=-3.2cm,yshift=-3.6cm]current page.north east) circle (4.6cm);
    \fill[poppurple!14] ([xshift=2.4cm,yshift=7.6cm]current page.south west) circle (3.8cm);
  \end{tikzpicture}%
  {\poplogo\fontsize{30}{30}\selectfont\color{popink}OnBuch\color{poporange}+}\hfill
  \raisebox{0.6ex}{\poppill{Cours signé \textbullet\ Premium}{popink}{popcream}}\par
  \vspace{1pt}\popsquiggle{poppurple}\par
  \vspace{8pt}
  \begin{tcolorbox}[enhanced,colback=poporange,colframe=popink,boxrule=2.8pt,arc=13pt,
      drop shadow=popink,left=18pt,right=18pt,top=12pt,bottom=13pt,after skip=10pt]
    \poppill{#2 \textbullet\ #3}{popink}{popcream}\hspace{5pt}\poppill{#4}{white}{popink}\par\vspace{8pt}
    {\popdisplay\fontsize{14}{14}\selectfont\color{popink}LEÇON #5}\par\vspace{4pt}
    {\raggedright\hyphenpenalty=10000\exhyphenpenalty=10000\popdisplay\fontsize{25}{27}\selectfont\color{popcream}\MakeUppercase{#1}\par}
    \vspace{8pt}
    {\popxbold\fontsize{9.5}{9.5}\selectfont\color{popink}\MakeUppercase{Durée conseillée : #6}}
  \end{tcolorbox}
  \begin{tcolorbox}[enhanced,colback=white,colframe=popink,boxrule=2.2pt,arc=10pt,
      drop shadow=popink,left=16pt,right=16pt,top=10pt,bottom=10pt,after skip=8pt]
    \poppill{Dans cette leçon}{poppurple}{white}\par\vspace{5pt}
    \popsemi\fontsize{9.3}{12.6}\selectfont\color{popink}#7
  \end{tcolorbox}
  \noindent\begin{minipage}[t]{0.487\linewidth}
    \begin{tcolorbox}[enhanced,colback=popgreen,colframe=popink,boxrule=2.2pt,arc=10pt,
        drop shadow=popink,left=11pt,right=11pt,top=10pt,bottom=10pt,height=3.1cm,valign=center]
      \tcbox[colback=white,colframe=popink,boxrule=1.3pt,arc=5pt,boxsep=0pt,nobeforeafter,
        left=3pt,right=3pt,top=3pt,bottom=3pt,valign=center]{\qrcode[height=1.9cm]{https://play.google.com/store/apps/details?id=com.luvvix.onbuch&hl=fr}}%
      \hspace{8pt}\begin{minipage}[c]{0.5\linewidth}
        {\popdisplay\fontsize{11}{12.5}\selectfont\color{white}\MakeUppercase{Télécharge OnBuch}}\par\vspace{4pt}
        {\raggedright\popsemi\fontsize{8.4}{11}\selectfont\color{white}Gratuit \textbullet\ Google Play\\Tuteur IA, annales, concours, résultats\par}
      \end{minipage}
    \end{tcolorbox}
  \end{minipage}\hfill
  \begin{minipage}[t]{0.487\linewidth}
    \begin{tcolorbox}[enhanced,colback=poppurple,colframe=popink,boxrule=2.2pt,arc=10pt,
        drop shadow=popink,left=11pt,right=11pt,top=10pt,bottom=10pt,height=3.1cm,valign=center]
      \tikz[baseline=-0.7ex]\node[circle,fill=white,draw=popink,line width=1.3pt,minimum size=1.6cm,inner sep=0]{\popdisplay\fontsize{24}{24}\selectfont\color{poppurple}+};%
      \hspace{8pt}\begin{minipage}[c]{0.6\linewidth}
        {\popdisplay\fontsize{11}{12.5}\selectfont\color{white}\MakeUppercase{Passe à OnBuch+}}\par\vspace{4pt}
        {\raggedright\popsemi\fontsize{8.4}{11}\selectfont\color{white}Tous les cours signés, Tuteur IA illimité, fiches PDF — onglet OnBuch+ de l'app\par}
      \end{minipage}
    \end{tcolorbox}
  \end{minipage}
  \vspace{8pt}
  \popcontenu
  \vfill
  \signatureonbuch
  \restoregeometry
  \newpage
}

% Rangée « Ce cours contient » (page de garde)
\newcommand{\poptile}[3]{% #1 couleur #2 gros texte #3 légende
  \begin{tcolorbox}[enhanced,colback=#1,colframe=popink,boxrule=2pt,arc=9pt,shadow={2.5pt}{-2.5pt}{0pt}{popink},
      left=6pt,right=6pt,top=9pt,bottom=9pt,halign=center,valign=center,height=1.75cm,width=\linewidth,nobeforeafter]
    {\popdisplay\fontsize{12.5}{13}\selectfont\color{white}\MakeUppercase{#2}}\par\vspace{3pt}
    {\centering\popsemi\fontsize{7.8}{9.5}\selectfont\color{white}#3\par}
  \end{tcolorbox}}
\newcommand{\popcontenu}{%
  \par\noindent{\popxboldls\fontsize{7.5}{7.5}\selectfont\color{popmuted}CE COURS SIGNÉ CONTIENT}\par\vspace{7pt}
  \noindent\begin{minipage}[t]{0.235\linewidth}\poptile{popblue}{Cours}{complet, illustré, pas à pas}\end{minipage}\hfill
  \begin{minipage}[t]{0.235\linewidth}\poptile{poporange}{Méthodes}{et exercices résolus}\end{minipage}\hfill
  \begin{minipage}[t]{0.235\linewidth}\poptile{poppink}{Exercices}{type examen + corrigés}\end{minipage}\hfill
  \begin{minipage}[t]{0.235\linewidth}\poptile{popgreen}{Fiche bilan}{l'essentiel à retenir}\end{minipage}\par}

% Sommaire
\newtcolorbox{sommaire}{enhanced,colback=white,colframe=popink,boxrule=2.2pt,arc=10pt,
  drop shadow=popink,left=18pt,right=18pt,top=13pt,bottom=13pt,before skip=6pt,after skip=20pt,
  before upper={\poppill{Sommaire}{poppurple}{white}\par\vspace{9pt}\popsemi\fontsize{10.6}{14.5}\selectfont\color{popink}}}

% Signature de fin de cours — poussée tout en bas de la dernière page
\newcommand{\findecours}{%
  \par\vfill\nopagebreak
  \begin{center}\popsquiggle{poporange}\end{center}\vspace{4pt}
  \signatureonbuch
}
\AtBeginDocument{\def\llbracket{\mathopen{[\mkern-3.5mu[}}\def\rrbracket{\mathclose{]\mkern-3.5mu]}}} % intervalles d'entiers (glyphe ⟦ absent de la police)
% Glyphes absents de Fira Math : redéfinis à partir de glyphes disponibles
\AtBeginDocument{%
  \def\setminus{\mathbin{\backslash}}\let\smallsetminus\setminus
  \def\vdots{\mathord{\vbox{\baselineskip3.2pt\lineskiplimit0pt\kern4pt\hbox{.}\hbox{.}\hbox{.}}}}%
  \def\ddots{\mathinner{\mkern1mu\raise6.5pt\hbox{.}\mkern2mu\raise3.6pt\hbox{.}\mkern2mu\raise0.7pt\hbox{.}\mkern1mu}}%
  \def\triangle{\mathord{\scalebox{0.9}{$\symup{\Delta}$}}}%
  \def\nabla{\mathord{\raisebox{\depth}{\scalebox{1}[-1]{$\symup{\Delta}$}}}}%
  \def\checkmark{\popcheck{popdark}}%
  \def\star{\mathbin{\ast}}\def\bigstar{\mathord{\ast}}%
  \def\diamond{\mathbin{\scalebox{0.6}{$\Diamond$}}}%
  \def\bigcup{\mathop{\mathchoice{\raisebox{-0.3em}{\scalebox{1.7}{$\cup$}}}{\scalebox{1.25}{$\cup$}}{\cup}{\cup}}\displaylimits}%
  \def\bigcap{\mathop{\mathchoice{\raisebox{-0.3em}{\scalebox{1.7}{$\cap$}}}{\scalebox{1.25}{$\cap$}}{\cap}{\cap}}\displaylimits}%
  \def\lesseqqgtr{\mathrel{\lessgtr}}\def\bigoplus{\mathop{\oplus}\displaylimits}%
}
\providecommand{\ssi}{si et seulement si} % abréviation fréquente
\AtBeginDocument{\providecommand{\wideparen}{\overparen}} % arc de cercle

%% FIGFIX-BEGIN v5 — correctifs globaux des figures (docs/onbuch-generation/tools/figfix)
\usepackage{adjustbox}\usepackage{etoolbox}\tcbuselibrary{raster}
% toute figure TikZ/pgfplots plus large que la ligne est réduite à la largeur disponible
\BeforeBeginEnvironment{tikzpicture}{\begin{adjustbox}{max width=\linewidth}}
\AfterEndEnvironment{tikzpicture}{\end{adjustbox}}
% légendes pgfplots lisibles par-dessus les courbes
\pgfplotsset{every axis/.append style={legend style={fill=white,fill opacity=0.92,text opacity=1,draw=black!15,inner sep=3pt}}}
% étiquettes d'arêtes (syntaxe edge["..."]) avec fond blanc
\tikzset{every edge quotes/.append style={fill=white,inner sep=1.2pt}}
%% FIGFIX-END
\begin{document}

\pagegarde{Géométrie analytique du plan}{Mathématiques}{1ère E}{Configurations et transformations élémentaires du plan}{N14}{10 h}{Cette leçon poursuit l'étude de la géométrie analytique en introduisant la distance d'un point à une droite, l'équation normale d'une droite et les équations paramétriques d'un cercle. Elle permet de résoudre des problèmes concrets de tangence et de distance, compétences clés pour le Probatoire C et E.
\vspace{4pt}
\begin{multicols}{2}\small
\begin{itemize}
\item À la fin de cette leçon, je sais déterminer un vecteur normal d'une droite à partir de son équation cartésienne ou d'un vecteur directeur
\item À la fin de cette leçon, je sais calculer la distance d'un point à une droite du plan
\item À la fin de cette leçon, je sais écrire une équation normale d'une droite et l'utiliser pour calculer des distances
\item À la fin de cette leçon, je sais déterminer le centre et le rayon d'un cercle à partir d'une de ses équations cartésiennes
\item À la fin de cette leçon, je sais écrire des équations paramétriques d'un cercle dont on connaît le centre et le rayon
\item À la fin de cette leçon, je sais déterminer une équation cartésienne de la tangente en un point d'un cercle
\end{itemize}\end{multicols}}

\begin{sommaire}
\begin{multicols}{2}
\begin{enumerate}
\item Vecteur normal et équation normale d'une droite
\item Distance d'un point à une droite
\item Cercle : centre, rayon et équations cartésiennes
\item Équations paramétriques d'un cercle
\item Tangente à un cercle en un de ses points
\item Tangentes à un cercle passant par un point extérieur
\item Activité d'intégration
\item Méthodes et exercices résolus
\item Exercices
\item Corrigés des exercices
\item Fiche bilan
\end{enumerate}
\end{multicols}
\end{sommaire}

\begin{objectifs}
\begin{itemize}
\item À la fin de cette leçon, je sais déterminer un vecteur normal d'une droite à partir de son équation cartésienne ou d'un vecteur directeur
\item À la fin de cette leçon, je sais calculer la distance d'un point à une droite du plan
\item À la fin de cette leçon, je sais écrire une équation normale d'une droite et l'utiliser pour calculer des distances
\item À la fin de cette leçon, je sais déterminer le centre et le rayon d'un cercle à partir d'une de ses équations cartésiennes
\item À la fin de cette leçon, je sais écrire des équations paramétriques d'un cercle dont on connaît le centre et le rayon
\item À la fin de cette leçon, je sais déterminer une équation cartésienne de la tangente en un point d'un cercle
\item À la fin de cette leçon, je sais déterminer les équations des tangentes à un cercle passant par un point extérieur
\item À la fin de cette leçon, je sais vérifier la position relative d'une droite et d'un cercle en comparant distance et rayon
\end{itemize}
\end{objectifs}

\begin{prerequis}
\begin{itemize}
\item Repère orthonormé, coordonnées d'un point et d'un vecteur
\item Produit scalaire et ses propriétés (orthogonalité, expression analytique)
\item Équation cartésienne d'une droite, vecteur directeur
\item Équation cartésienne d'un cercle : (x − a)² + (y − b)² = R²
\item Distance entre deux points, norme d'un vecteur
\end{itemize}
\end{prerequis}

\newpage

\coursec{Vecteur normal et équation normale d'une droite}

\textbf{Situation d'accroche.} Sur le site de la SONARA à Limbé, un technicien doit installer une canalisation rectiligne entre deux postes de pompage. La réglementation de sécurité impose que la canalisation passe à au moins \SI{50}{m} d'un réservoir cylindrique de stockage d'hydrocarbures. Sur le plan du site, muni d'un repère orthonormé, le réservoir est assimilé à un point $A$ et la canalisation à une droite $(D)$. Deux questions se posent alors au technicien :
\begin{itemize}
\item comment calculer \emph{exactement} la distance entre le point $A$ (le réservoir) et la droite $(D)$ (la canalisation), sans avoir à mesurer sur le terrain ?
\item comment déterminer la position de la tangente au réservoir menée depuis un point extérieur, par exemple pour poser une échelle de maintenance ?
\end{itemize}
Ces problèmes d'ingénierie, en apparence très concrets, se résolvent entièrement par le calcul, grâce à la \cle{géométrie analytique} : on traduit les objets géométriques (points, droites, cercles) par des équations, et les propriétés géométriques (perpendicularité, distance, tangence) par des relations algébriques. Toute cette leçon repose sur une idée simple mais puissante : une droite peut être décrite non seulement par sa direction, mais aussi par la direction qui lui est \emph{perpendiculaire}. C'est l'objet de cette première section.

\courssub{Équation cartésienne d'une droite et vecteur directeur : rappels}

Dans tout le chapitre, le plan est muni d'un repère orthonormé $(O\,;\vec{\imath},\vec{\jmath})$.

Tu as vu en classe de Seconde qu'une droite $(D)$ du plan admet une \cle{équation cartésienne} de la forme
\[ ax + by + c = 0, \quad \text{avec } (a\,;b) \neq (0\,;0), \]
ce qui signifie qu'un point $M(x\,;y)$ appartient à $(D)$ si et seulement si ses coordonnées vérifient cette relation. La condition $(a\,;b)\neq(0\,;0)$ est essentielle : sinon l'équation se réduirait à $c=0$, qui ne décrit pas une droite.

Tu as aussi rencontré la notion de \cle{vecteur directeur} : un vecteur non nul $\vec{u}$ est directeur de $(D)$ s'il a la même direction que $(D)$, c'est-à-dire s'il est colinéaire à tout vecteur $\overrightarrow{AB}$ joignant deux points de $(D)$.

\begin{propriete}[Vecteur directeur lu sur l'équation cartésienne]
Si la droite $(D)$ a pour équation cartésienne $ax+by+c=0$, alors le vecteur
\[ \vec{u}\begin{pmatrix} -b \\ a \end{pmatrix} \]
est un vecteur directeur de $(D)$. Tout vecteur directeur de $(D)$ est colinéaire à $\vec{u}$.
\end{propriete}

\begin{exemplebox}[Lecture directe]
La droite $(D) : 3x + 4y - 10 = 0$ admet pour vecteur directeur $\vec{u}\begin{pmatrix} -4 \\ 3 \end{pmatrix}$ : on échange les coefficients de $x$ et $y$, et on change le signe de l'un des deux. Le vecteur $\begin{pmatrix} 8 \\ -6 \end{pmatrix}$ est aussi directeur de $(D)$, car il est colinéaire à $\vec{u}$.
\end{exemplebox}

\courssub{Vecteur normal à une droite}

\textbf{Intuition.} Reprenons notre canalisation. Pour mesurer la distance du réservoir à la canalisation, le technicien ne marche pas \emph{le long} de la canalisation : il marche \emph{perpendiculairement} à celle-ci, en ligne droite vers le réservoir. La direction de son déplacement est une direction \emph{normale} à la droite. C'est exactement l'idée mathématique : à côté des vecteurs directeurs, qui « suivent » la droite, il existe des vecteurs qui lui sont orthogonaux, et ce sont eux qui pilotent les calculs de distances et d'angles.

\begin{definition}[Vecteur normal à une droite]
Soit $(D)$ une droite du plan. On appelle \cle{vecteur normal} à $(D)$ tout vecteur \textbf{non nul} $\vec{n}$ orthogonal à un vecteur directeur de $(D)$.
\end{definition}

Deux remarques importantes découlent immédiatement de cette définition :
\begin{itemize}
\item si $\vec{n}$ est normal à $(D)$, alors tout vecteur $k\vec{n}$ avec $k \neq 0$ est aussi normal à $(D)$ : le vecteur normal est donc \textbf{unique à un facteur multiplicatif non nul près} ;
\item dire que $\vec{n}$ est normal à $(D)$, c'est dire que la droite portée par $\vec{n}$ est perpendiculaire à $(D)$.
\end{itemize}

\begin{propriete}[Vecteur normal lu sur l'équation cartésienne]
Si la droite $(D)$ a pour équation cartésienne $ax+by+c=0$, alors le vecteur
\[ \vec{n}\begin{pmatrix} a \\ b \end{pmatrix} \]
est un vecteur normal à $(D)$. \emph{(Démonstration exigible au Probatoire.)}
\end{propriete}

\begin{experience}[Démonstration guidée : pourquoi $\vec{n}(a\,;b)$ est-il normal à $(D)$ ?]
On sait que $\vec{u}\begin{pmatrix} -b \\ a \end{pmatrix}$ dirige $(D)$. Pour prouver que $\vec{n}\begin{pmatrix} a \\ b \end{pmatrix}$ est normal à $(D)$, il suffit de montrer que $\vec{n}$ et $\vec{u}$ sont orthogonaux. Or, dans un repère orthonormé, deux vecteurs sont orthogonaux si et seulement si leur produit scalaire est nul. Calculons :
\begin{align*}
\vec{n}\cdot\vec{u} &= a\times(-b) + b\times a \\
&= -ab + ab \\
&= 0.
\end{align*}
Le produit scalaire est nul, donc $\vec{n}\perp\vec{u}$ : le vecteur $\vec{n}\begin{pmatrix} a \\ b \end{pmatrix}$ est bien orthogonal à un vecteur directeur de $(D)$, c'est donc un vecteur normal à $(D)$. De plus $\vec{n}\neq\vec{0}$ car $(a\,;b)\neq(0\,;0)$. \hfill$\square$
\end{experience}

\begin{aretenir}[À retenir absolument]
Pour une droite $(D) : ax+by+c=0$ :
\begin{itemize}
\item un \textbf{vecteur directeur} : $\vec{u}\begin{pmatrix} -b \\ a \end{pmatrix}$ (on échange et on change un signe) ;
\item un \textbf{vecteur normal} : $\vec{n}\begin{pmatrix} a \\ b \end{pmatrix}$ (on recopie les coefficients dans l'ordre).
\end{itemize}
Le coefficient $c$ n'intervient ni dans l'un ni dans l'autre : il fixe la \emph{position} de la droite, pas sa \emph{direction}.
\end{aretenir}

\begin{popfigure}
\begin{tikzpicture}[scale=1.1,>=Stealth]
% Droite (D) : direction u = (3,1)
\draw[popblue,thick] (-1.2,-0.4) -- (4.2,1.4) node[below right] {$(D)$};
% Point A sur la droite
\coordinate (A) at (1.5,0.5);
\fill[popdark] (A) circle (1.5pt) node[below left=1pt] {$A$};
% Vecteur directeur u = (1.5,0.5)
\draw[->,popgreen,very thick] (A) -- ++(1.5,0.5) node[midway,below right] {$\vec{u}$};
% Vecteur normal n = (-0.5,1.5) (orthogonal à u)
\draw[->,poporange,very thick] (A) -- ++(-0.5,1.5) node[midway,above left] {$\vec{n}$};
% Angle droit
\draw[popdark] ($(A)+(0.3,0.1)$) -- ++(-0.1,0.3) -- ++(-0.3,-0.1);
\end{tikzpicture}
\legende{Depuis un point $A$ de $(D)$ : un vecteur directeur $\vec{u}$ (vert) suit la droite, un vecteur normal $\vec{n}$ (orange) lui est perpendiculaire. L'angle droit est marqué.}
\end{popfigure}

\textbf{Passage dans les deux sens.} Le lien entre équation et vecteur normal fonctionne dans les deux directions :
\begin{itemize}
\item \emph{De l'équation vers le vecteur} : si $(D) : 2x - 5y + 7 = 0$, alors $\vec{n}\begin{pmatrix} 2 \\ -5 \end{pmatrix}$ est normal à $(D)$.
\item \emph{Du vecteur vers l'équation} : si $\vec{n}\begin{pmatrix} a \\ b \end{pmatrix}$ est normal à une droite $(D)$, alors $(D)$ admet une équation de la forme $ax+by+c=0$, où seul $c$ reste à déterminer, généralement à l'aide d'un point de $(D)$.
\end{itemize}

\begin{exoresolu}[Trouver l'équation d'une droite connaissant un point et un vecteur normal]
Déterminer une équation cartésienne de la droite $(D)$ passant par $A(1\,;-2)$ et de vecteur normal $\vec{n}\begin{pmatrix} 3 \\ 4 \end{pmatrix}$.
\tcblower
\textbf{Solution.} Puisque $\vec{n}\begin{pmatrix} 3 \\ 4 \end{pmatrix}$ est normal à $(D)$, une équation de $(D)$ est de la forme
\[ 3x + 4y + c = 0. \]
Le point $A(1\,;-2)$ appartient à $(D)$, donc ses coordonnées vérifient l'équation :
\[ 3\times 1 + 4\times(-2) + c = 0 \iff 3 - 8 + c = 0 \iff c = 5. \]
Une équation cartésienne de $(D)$ est donc $\boxed{3x + 4y + 5 = 0}$.
\end{exoresolu}

\courssub{Équation normale d'une droite}

\textbf{Intuition.} Parmi tous les vecteurs normaux à une droite $(D)$, il y en a deux qui sont particulièrement intéressants : ceux de \emph{norme $1$} (vecteurs unitaires). Pourquoi ? Parce qu'avec un vecteur normal unitaire, les coefficients de l'équation cessent d'être arbitraires et prennent une signification géométrique précise : ils deviennent le cosinus et le sinus d'un angle, et le second membre devient la \emph{distance de l'origine à la droite}. C'est toute la richesse de l'équation normale.

Soit $(D)$ une droite, et soit $H$ le projeté orthogonal de l'origine $O$ sur $(D)$. Notons $p = OH$ la distance de $O$ à $(D)$, et $\theta$ une mesure de l'angle $\left(\vec{\imath}\,;\overrightarrow{OH}\right)$ (si $H\neq O$). Le vecteur $\overrightarrow{OH}$ est normal à $(D)$, et le vecteur unitaire de même direction est
\[ \vec{n}_0 = \begin{pmatrix} \cos\theta \\ \sin\theta \end{pmatrix}. \]

\begin{definition}[Équation normale d'une droite]
On appelle \cle{équation normale} de la droite $(D)$ toute équation de la forme
\[ x\cos\theta + y\sin\theta = p, \]
où $p \geq 0$ est la \textbf{distance de l'origine $O$ à la droite $(D)$}, et $\theta$ est une mesure de l'angle que fait le vecteur normal unitaire $\vec{n}_0\begin{pmatrix} \cos\theta \\ \sin\theta \end{pmatrix}$ avec l'axe des abscisses.
\end{definition}

\begin{popfigure}
\begin{tikzpicture}[scale=1.3,>=Stealth]
% Repère
\draw[->,popmuted] (-0.6,0) -- (4.2,0) node[below] {$x$};
\draw[->,popmuted] (0,-0.6) -- (0,3.2) node[left] {$y$};
\fill[popdark] (0,0) circle (1.5pt) node[below left] {$O$};
% Droite (D) : normale à 40°, p = 2. H = (2cos40, 2sin40) = (1.532, 1.286)
\coordinate (H) at (1.532,1.286);
% Direction de la droite : (-sin40, cos40) = (-0.643, 0.766)
\draw[popblue,thick] ($(H)+2.2*(-0.643,0.766)$) -- ($(H)-2.6*(-0.643,0.766)$) node[below left] {$(D)$};
% Segment OH
\draw[->,poporange,very thick] (0,0) -- (H) node[midway,above right,inner sep=1pt] {$p = OH$};
\fill[popdark] (H) circle (1.5pt) node[above left=1pt] {$H$};
% Angle droit en H
\draw[popdark] ($(H)+0.18*(-0.643,0.766)$) -- ++($-0.18*(0.766,0.643)$) -- ++($-0.18*(-0.643,0.766)$);
% Angle theta
\draw[popgreen,thick] (0.9,0) arc (0:40:0.9) node[midway,right] {$\theta$};
% Vecteur normal unitaire
\draw[->,popgreen,thick] (H) -- ++(0.766,0.643) node[right] {$\vec{n}_0$};
\end{tikzpicture}
\legende{Le projeté orthogonal $H$ de $O$ sur $(D)$ : $p=OH$ est la distance de l'origine à la droite, et $\theta$ est l'angle entre l'axe des abscisses et la direction normale $\overrightarrow{OH}$.}
\end{popfigure}

\begin{aretenir}[Lecture géométrique de l'équation normale]
Dans l'équation $x\cos\theta + y\sin\theta = p$ :
\begin{itemize}
\item le vecteur $\vec{n}_0\begin{pmatrix} \cos\theta \\ \sin\theta \end{pmatrix}$ est un vecteur normal \textbf{unitaire} à $(D)$, car $\cos^2\theta+\sin^2\theta = 1$ ;
\item $\theta$ indique la \emph{direction} de la perpendiculaire à $(D)$ issue de $O$ ;
\item $p \geq 0$ est la \emph{distance} de $O$ à $(D)$ : $p = d(O\,;D) = OH$.
\end{itemize}
\end{aretenir}

\courssub{Passage d'une équation cartésienne à une équation normale}

Partons d'une équation cartésienne quelconque $ax+by+c=0$. Le vecteur $\vec{n}\begin{pmatrix} a \\ b \end{pmatrix}$ est normal à la droite, mais sa norme $\sqrt{a^2+b^2}$ n'a aucune raison de valoir $1$. Pour obtenir un vecteur normal \emph{unitaire}, il suffit de diviser $\vec{n}$ par sa norme, c'est-à-dire de diviser toute l'équation par $\sqrt{a^2+b^2}$. Reste un détail crucial : dans une équation normale, le second membre $p$ représente une distance, donc il doit être \textbf{positif ou nul}. On choisit donc le signe du diviseur en conséquence.

\begin{methode}[Normaliser une équation cartésienne]
Pour transformer $(D) : ax+by+c=0$ en équation normale :
\begin{enumerate}
\item calculer $\rho = \sqrt{a^2+b^2}$ ;
\item choisir le signe $\varepsilon = +1$ si $c<0$, et $\varepsilon = -1$ si $c>0$ (de sorte que $-\varepsilon c > 0$) ; si $c=0$, le signe est libre ;
\item diviser les deux membres par $\varepsilon\rho$ :
\[ \frac{a}{\varepsilon\rho}\,x + \frac{b}{\varepsilon\rho}\,y = \frac{-c}{\varepsilon\rho}\,; \]
\item identifier : $\cos\theta = \dfrac{a}{\varepsilon\rho}$, $\sin\theta = \dfrac{b}{\varepsilon\rho}$ et $p = \dfrac{-c}{\varepsilon\rho} = \dfrac{|c|}{\sqrt{a^2+b^2}} \geq 0$.
\end{enumerate}
On obtient alors directement la distance de l'origine à la droite : $d(O\,;D) = p$.
\end{methode}

\begin{exemplebox}[Exemple chiffré de référence]
Soit $(D) : 3x + 4y - 10 = 0$.
\begin{enumerate}
\item $\rho = \sqrt{3^2+4^2} = \sqrt{9+16} = \sqrt{25} = 5$.
\item Ici $c = -10 < 0$, donc on divise par $+5$ :
\[ \frac{3}{5}x + \frac{4}{5}y = 2. \]
\item On lit : $\cos\theta = \dfrac{3}{5}$, $\sin\theta = \dfrac{4}{5}$ (donc $\theta \approx 53{,}1^\circ$) et $p = 2$.
\end{enumerate}
Conclusion : la distance de l'origine $O$ à la droite $(D)$ vaut $d(O\,;D) = 2$, et le vecteur normal unitaire $\vec{n}_0\begin{pmatrix} 3/5 \\ 4/5 \end{pmatrix}$ pointe de $O$ vers la droite. Vérification : $\left(\dfrac{3}{5}\right)^2+\left(\dfrac{4}{5}\right)^2 = \dfrac{9+16}{25}=1$ : le vecteur est bien unitaire.
\end{exemplebox}

\begin{exoresolu}[Normaliser avec un second membre négatif à corriger]
Déterminer une équation normale de la droite $(D) : x - y + 3 = 0$ et en déduire la distance de $O$ à $(D)$.
\tcblower
\textbf{Solution.} On a $a=1$, $b=-1$, $c=3$.
\begin{enumerate}
\item $\rho = \sqrt{1^2+(-1)^2} = \sqrt{2}$.
\item Comme $c = 3 > 0$, on divise par $-\sqrt{2}$ pour rendre le second membre positif :
\[ -\frac{1}{\sqrt{2}}\,x + \frac{1}{\sqrt{2}}\,y = \frac{3}{\sqrt{2}}, \qquad\text{soit encore}\qquad -\frac{\sqrt{2}}{2}\,x + \frac{\sqrt{2}}{2}\,y = \frac{3\sqrt{2}}{2}. \]
\item On lit $\cos\theta = -\dfrac{\sqrt{2}}{2}$, $\sin\theta = \dfrac{\sqrt{2}}{2}$, donc $\theta = \dfrac{3\pi}{4}$, et $p = \dfrac{3\sqrt{2}}{2}$.
\end{enumerate}
La distance de l'origine à $(D)$ est donc $d(O\,;D) = \dfrac{3\sqrt{2}}{2} \approx 2{,}12$.
\end{exoresolu}

\begin{attention}[Les trois erreurs classiques]
\begin{itemize}
\item \textbf{Oublier le signe.} Si tu divises $x - y + 3 = 0$ par $+\sqrt{2}$, tu obtiens un second membre $-\dfrac{3\sqrt{2}}{2}$ \emph{négatif}, qui ne peut pas être une distance. Choisis toujours le signe du diviseur pour que $p \geq 0$.
\item \textbf{Diviser par $a^2+b^2$ au lieu de $\sqrt{a^2+b^2}$.} Le facteur de normalisation est la \emph{norme} du vecteur normal, donc la racine carrée de la somme des carrés. Diviser $3x+4y-10=0$ par $25$ au lieu de $5$ donnerait $p = 0{,}4$ au lieu de $2$ : résultat faux.
\item \textbf{Confondre vecteur normal et vecteur directeur.} Pour $ax+by+c=0$, le vecteur normal est $\begin{pmatrix} a \\ b \end{pmatrix}$ (coefficients dans l'ordre), le vecteur directeur est $\begin{pmatrix} -b \\ a \end{pmatrix}$ (échange et changement de signe). Une inversion ruine tout le calcul d'angle.
\end{itemize}
\end{attention}

\begin{savaistu}[Pourquoi « normale » ?]
Le mot « normale » vient du latin \emph{norma}, l'équerre du charpentier : une droite normale est une droite « à l'équerre », c'est-à-dire perpendiculaire. L'équation normale, introduite systématiquement au XIX\textsuperscript{e} siècle notamment par Ludwig Otto Hesse (forme dite « normale de Hesse »), est un outil de base en infographie et en robotique : pour savoir si un robot mobile risque de percuter un mur rectiligne, on calcule exactement la distance d'un point à une droite, comme notre technicien de la SONARA avec son réservoir. Dans la suite de la leçon, cette même idée nous donnera la formule générale de la distance d'un point quelconque à une droite, puis les équations des tangentes à un cercle.
\end{savaistu}

\begin{exercice}[À toi de jouer]
\begin{enumerate}
\item Déterminer un vecteur normal et un vecteur directeur de chacune des droites : $(D_1) : 5x - 12y + 8 = 0$ et $(D_2) : y = 2x - 3$.
\item Déterminer une équation cartésienne de la droite passant par $B(-1\,;3)$ et de vecteur normal $\vec{n}\begin{pmatrix} 2 \\ -1 \end{pmatrix}$.
\item Écrire une équation normale de $(D_1) : 5x - 12y + 8 = 0$ et en déduire $d(O\,;D_1)$.
\end{enumerate}
\end{exercice}

\begin{corrige}[Exercice n°1]
\begin{enumerate}
\item Pour $(D_1)$ : $\vec{n}_1\begin{pmatrix} 5 \\ -12 \end{pmatrix}$ et $\vec{u}_1\begin{pmatrix} 12 \\ 5 \end{pmatrix}$. Pour $(D_2)$, on réécrit l'équation sous forme cartésienne : $2x - y - 3 = 0$, d'où $\vec{n}_2\begin{pmatrix} 2 \\ -1 \end{pmatrix}$ et $\vec{u}_2\begin{pmatrix} 1 \\ 2 \end{pmatrix}$.
\item L'équation est de la forme $2x - y + c = 0$. Comme $B(-1\,;3)\in(D)$ : $2\times(-1) - 3 + c = 0$, soit $c = 5$. D'où $(D) : 2x - y + 5 = 0$.
\item $\rho = \sqrt{25+144} = 13$. Comme $c = 8 > 0$, on divise par $-13$ :
\[ -\frac{5}{13}x + \frac{12}{13}y = \frac{8}{13}. \]
On en déduit $d(O\,;D_1) = \dfrac{8}{13} \approx 0{,}62$.
\end{enumerate}
\end{corrige}

\coursec{Distance d'un point à une droite}

Dans la section précédente, nous avons associé à toute droite $(D)$ un \cle{vecteur normal} $\vec{n}(a\,;b)$ et une \cle{équation normale} $x\cos\theta + y\sin\theta - p = 0$, dans laquelle le nombre $p$ mesure la distance de l'origine du repère à la droite. Nous allons maintenant généraliser cette idée : au lieu de mesurer la distance de l'origine à une droite, nous voulons mesurer la distance de \emph{n'importe quel point} du plan à \emph{n'importe quelle droite}. C'est l'un des outils les plus utiles de la géométrie analytique : il permet de calculer des hauteurs de triangles, des aires, des distances entre routes parallèles, et il sera indispensable lorsque nous étudierons les tangentes à un cercle.

\courssub{Notion de distance d'un point à une droite}

\textbf{Intuition.}\quad Imagine que tu te trouves en un point $M_0$ au bord d'une route rectiligne $(D)$, et que tu veuilles traverser pour rejoindre la route le plus vite possible. Quel chemin choisir ? L'expérience (et le théorème de Pythagore) montre que le chemin le plus court est le segment \emph{perpendiculaire} à la route. Tout autre point $K$ de la route donne un trajet $M_0K$ strictement plus long, car dans le triangle $M_0HK$ rectangle en $H$, l'hypoténuse $M_0K$ est plus grande que le côté $M_0H$.

\begin{definition}[Distance d'un point à une droite]
Soit $(D)$ une droite du plan et $M_0$ un point. On appelle \cle{projeté orthogonal} de $M_0$ sur $(D)$ l'unique point $H$ de $(D)$ tel que $(M_0H) \perp (D)$.\\
On appelle \cle{distance du point $M_0$ à la droite $(D)$}, notée $d(M_0\,;D)$, le nombre :
\[ d(M_0\,;D) = M_0H. \]
C'est la plus petite des distances $M_0K$ lorsque $K$ décrit $(D)$. En particulier, $d(M_0\,;D) = 0$ si et seulement si $M_0 \in (D)$.
\end{definition}

\begin{popfigure}
\begin{tikzpicture}[scale=1.1,>=Stealth]
\coordinate (A) at (-0.5,-0.75);
\coordinate (B) at (6,3.15);
\coordinate (M) at (2.2,3.4);
\coordinate (H) at (2.62,1.86);
\draw[popline,thick] (A) -- (B) node[pos=0.95,above left,popdark]{$(D)$};
\draw[poporange,very thick,dashed] (M) -- (H);
\fill[popblue] (M) circle (2pt) node[above right,popdark]{$M_0$};
\fill[popink] (H) circle (2pt) node[below right,popdark]{$H$};
% Marque d'angle droit robuste avec calc
\coordinate (H1) at ($(H)!0.3!(M)$);
\coordinate (H2) at ($(H)!0.3!(B)$);
\draw[popmuted] (H1) -- ($(H1)!0.3!-90:(H2)$) -- ($(H2)!0.3!90:(H1)$);
\draw[->,popgreen,thick] (H) -- ++(0.45,-0.675) node[below,popdark]{$\vec{n}$};
\fill[popmuted] (4.6,2.55) circle (1.5pt) node[above,popdark]{$K$};
\draw[popmuted,dashed] (M) -- (4.6,2.55);
\end{tikzpicture}
\legende{Le projeté orthogonal $H$ de $M_0$ sur $(D)$ : la distance $d(M_0\,;D)=M_0H$ est le plus court chemin (pour tout autre point $K$ de $(D)$, $M_0K > M_0H$).}
\end{popfigure}

\courssub{Le théorème fondamental}

Le plan est muni d'un repère orthonormé $(O\,;\vec{i},\vec{j})$. La droite $(D)$ admet une équation cartésienne $ax + by + c = 0$ (avec $(a\,;b) \neq (0\,;0)$), et on note $M_0(x_0\,;y_0)$. Comment calculer $M_0H$ sans avoir à déterminer les coordonnées de $H$ ? L'idée géniale consiste à \emph{projeter} le vecteur $\overrightarrow{HM_0}$ sur la direction du vecteur normal $\vec{n}(a\,;b)$ : comme $\overrightarrow{HM_0}$ est précisément colinéaire à $\vec{n}$, la valeur absolue de cette projection est exactement $M_0H$.

\begin{propriete}[Formule de la distance d'un point à une droite]
Dans un repère orthonormé, soit $(D)$ la droite d'équation $ax + by + c = 0$ et $M_0(x_0\,;y_0)$ un point. Alors :
\[ \boxed{\,d(M_0\,;D) = \dfrac{|ax_0 + by_0 + c|}{\sqrt{a^2 + b^2}}\,} \]
\emph{Démonstration exigible au Probatoire : elle est guidée ci-dessous.}
\end{propriete}

\begin{experience}[Démonstration guidée du théorème]
On note $H$ le projeté orthogonal de $M_0$ sur $(D)$, et $\vec{n}(a\,;b)$ un vecteur normal à $(D)$.
\begin{enumerate}
\item \textbf{Introduire le vecteur normal unitaire.} Le vecteur $\vec{u} = \dfrac{\vec{n}}{\|\vec{n}\|} = \left(\dfrac{a}{\sqrt{a^2+b^2}}\,;\dfrac{b}{\sqrt{a^2+b^2}}\right)$ est un vecteur normal à $(D)$ de norme $1$.
\item \textbf{Projeter $\overrightarrow{HM_0}$ sur $\vec{u}$.} Comme $\overrightarrow{HM_0}$ est colinéaire à $\vec{n}$ (la droite $(M_0H)$ est perpendiculaire à $(D)$), la valeur absolue du produit scalaire donne :
\[ \left|\overrightarrow{HM_0}\cdot\vec{u}\right| = \|\overrightarrow{HM_0}\|\times\|\vec{u}\| = M_0H = d(M_0\,;D). \]
\item \textbf{Choisir un point auxiliaire de $(D)$.} Soit $A(x_A\,;y_A)$ un point quelconque de $(D)$. On a $ax_A + by_A + c = 0$, c'est-à-dire $c = -ax_A - by_A$.
\item \textbf{Calculer le produit scalaire avec les coordonnées.} En écrivant $\overrightarrow{HM_0}\cdot\vec{u} = \overrightarrow{AM_0}\cdot\vec{u}$ (car $\overrightarrow{HA}\cdot\vec{u}=0$, le vecteur $\overrightarrow{HA}$ étant directeur de $(D)$ et donc orthogonal à $\vec{u}$) :
\begin{align*}
\overrightarrow{AM_0}\cdot\vec{u} &= (x_0 - x_A)\dfrac{a}{\sqrt{a^2+b^2}} + (y_0 - y_A)\dfrac{b}{\sqrt{a^2+b^2}}\\
&= \dfrac{ax_0 + by_0 - (ax_A + by_A)}{\sqrt{a^2+b^2}} = \dfrac{ax_0 + by_0 + c}{\sqrt{a^2+b^2}}.
\end{align*}
\item \textbf{Conclure.} En prenant la valeur absolue : $d(M_0\,;D) = \dfrac{|ax_0 + by_0 + c|}{\sqrt{a^2+b^2}}$. $\blacksquare$
\end{enumerate}
\end{experience}

\begin{aretenir}[Lien avec l'équation normale]
Si l'on utilise l'\cle{équation normale} de $(D)$ : $x\cos\theta + y\sin\theta - p = 0$, alors $\cos^2\theta + \sin^2\theta = 1$, donc le dénominateur vaut $1$ et la formule devient simplement :
\[ d(M_0\,;D) = \left|x_0\cos\theta + y_0\sin\theta - p\right|. \]
C'est une raison supplémentaire d'aimer l'équation normale : elle « lit » directement les distances, sans division.
\end{aretenir}

\begin{exemplebox}[Un calcul direct]
Soit $A(1\,;2)$ et $(D)$ la droite d'équation $3x - 4y + 10 = 0$. On a $a=3$, $b=-4$, $c=10$, donc $\sqrt{a^2+b^2} = \sqrt{9+16} = \sqrt{25} = 5$ et :
\[ d(A\,;D) = \dfrac{|3(1) - 4(2) + 10|}{5} = \dfrac{|3 - 8 + 10|}{5} = \dfrac{|5|}{5} = 1. \]
Le point $A$ est donc à une distance $1$ de la droite $(D)$.
\end{exemplebox}

\begin{attention}[Les deux erreurs classiques]
\begin{itemize}
\item \textbf{Oublier la valeur absolue au numérateur.} Une distance est toujours positive ! Si tu trouves $\dfrac{-5}{5} = -1$, c'est que tu as oublié les barres de valeur absolue. Le signe de $ax_0+by_0+c$ indique seulement de quel \emph{côté} de la droite se trouve le point.
\item \textbf{Oublier la racine carrée au dénominateur.} On divise par $\sqrt{a^2+b^2}$ (norme du vecteur normal), \emph{pas} par $a^2+b^2$ ni par $a+b$. Vérifie aussi que l'équation est bien écrite sous la forme $ax+by+c=0$ (tout dans le membre de gauche) avant d'identifier $a$, $b$, $c$.
\end{itemize}
\end{attention}

\courssub{Applications}

\courssub{Distance entre deux droites parallèles}

Deux droites parallèles $(D)$ et $(D')$ sont partout à la même distance l'une de l'autre : pour la mesurer, il suffit de prendre \emph{n'importe quel} point de l'une et de calculer sa distance à l'autre.

\begin{methode}[Calculer la distance entre deux droites parallèles]
\begin{enumerate}
\item Vérifier que $(D)$ et $(D')$ sont bien parallèles (leurs vecteurs normaux sont colinéaires, ou leurs coefficients $a$ et $b$ sont proportionnels).
\item Choisir un point $M_0$ de $(D)$ : on fixe par exemple une coordonnée commode (souvent $0$) pour déterminer l'autre.
\item Calculer $d(M_0\,;D')$ avec la formule du théorème : c'est la distance cherchée $d(D\,;D')$.
\end{enumerate}
\end{methode}

\begin{exoresolu}[Distance entre deux droites parallèles]
Soient $(D) : 2x - y + 3 = 0$ et $(D') : 2x - y - 5 = 0$. Calculer $d(D\,;D')$.
\tcblower
\textbf{Solution.} Les deux droites ont le même vecteur normal $\vec{n}(2\,;-1)$ : elles sont parallèles.\\
\textbf{Étape 1 : un point de $(D)$.} Pour $x_0 = 0$, on obtient $-y_0 + 3 = 0$, soit $y_0 = 3$. Le point $M_0(0\,;3)$ appartient à $(D)$.\\
\textbf{Étape 2 : distance de $M_0$ à $(D')$.}
\[ d(D\,;D') = d(M_0\,;D') = \dfrac{|2(0) - 3 - 5|}{\sqrt{2^2 + (-1)^2}} = \dfrac{|-8|}{\sqrt{5}} = \dfrac{8}{\sqrt{5}} = \dfrac{8\sqrt{5}}{5} \approx 3{,}58. \]
\end{exoresolu}

\courssub{Hauteur et aire d'un triangle}

Dans un triangle $ABC$, la \cle{hauteur} issue de $A$ est précisément la distance du point $A$ à la droite $(BC)$. L'aire s'obtient alors par la formule classique $\mathcal{A} = \dfrac{1}{2}\times\text{base}\times\text{hauteur}$.

\begin{exoresolu}[Aire d'un triangle connaissant ses trois sommets]
Dans un repère orthonormé, on donne $A(1\,;5)$, $B(-2\,;0)$ et $C(4\,;0)$. Calculer l'aire du triangle $ABC$.
\tcblower
\textbf{Solution.} \textbf{Étape 1 : équation de $(BC)$.} Un vecteur directeur est $\overrightarrow{BC}(6\,;0)$ ; la droite $(BC)$ est l'axe des abscisses : $y = 0$, soit $y = 0$ avec $a=0$, $b=1$, $c=0$.\\
\textbf{Étape 2 : hauteur issue de $A$.}
\[ h_A = d(A\,;(BC)) = \dfrac{|0\times 1 + 1\times 5 + 0|}{\sqrt{0^2+1^2}} = 5. \]
\textbf{Étape 3 : aire.} La base $BC = \sqrt{6^2+0^2} = 6$, donc :
\[ \mathcal{A} = \dfrac{1}{2}\times BC \times h_A = \dfrac{1}{2}\times 6\times 5 = 15 \text{ unités d'aire.} \]
\end{exoresolu}

\begin{popfigure}
\begin{tikzpicture}[scale=0.85,>=Stealth]
\draw[->,popmuted] (-3,0) -- (5.2,0) node[below left]{$x$};
\draw[->,popmuted] (0,-0.8) -- (0,6) node[below left]{$y$};
\coordinate (B) at (-2,0);
\coordinate (C) at (4,0);
\coordinate (A) at (1,5);
\coordinate (H) at (1,0);
\draw[popblue,very thick] (B) -- (C) -- (A) -- cycle;
\fill[poppinkL,opacity=0.5] (B) -- (C) -- (A) -- cycle;
\draw[poporange,very thick,dashed] (A) -- (H);
\draw[popmuted] ($(H)+(0.25,0)$) -- ($(H)+(0.25,0.25)$) -- ($(H)+(0,0.25)$);
\fill[popink] (A) circle (2pt) node[above,popdark]{$A(1\,;5)$};
\fill[popink] (B) circle (2pt) node[below left,popdark]{$B(-2\,;0)$};
\fill[popink] (C) circle (2pt) node[below right,popdark]{$C(4\,;0)$};
\fill[popink] (H) circle (1.8pt) node[below,popdark]{$H$};
\node[poporange] at (1.35,2.6) {$h_A = 5$};
\node[popdark] at (1,-0.55) {base $BC = 6$};
\end{tikzpicture}
\legende{La hauteur issue de $A$ est la distance de $A$ à la droite $(BC)$ ; l'aire vaut $\frac12 \times 6 \times 5 = 15$.}
\end{popfigure}

\courssub{Complément : position d'un point par rapport à une bande}

\begin{aretenir}[Point entre deux droites parallèles (complément utile)]
Soient $(D) : ax+by+c=0$ et $(D') : ax+by+c'=0$ deux droites parallèles (mêmes coefficients $a$ et $b$). Un point $M(x_0\,;y_0)$ est situé \cle{strictement entre} les deux droites (dans la « bande » qu'elles délimitent) si et seulement si les nombres $ax_0+by_0+c$ et $ax_0+by_0+c'$ sont de \textbf{signes contraires}, c'est-à-dire :
\[ (ax_0+by_0+c)(ax_0+by_0+c') < 0. \]
Ce critère pratique permet par exemple de vérifier si un point se trouve dans un couloir délimité par deux frontières parallèles.
\end{aretenir}

\begin{savaistu}[La formule des distances au service de l'aménagement du territoire]
La formule $d(M_0\,;D) = \dfrac{|ax_0+by_0+c|}{\sqrt{a^2+b^2}}$ est utilisée chaque jour en \cle{géomatique}, dans les Systèmes d'Information Géographique (SIG). Au Cameroun, les services d'urbanisme de Douala ou de Yaoundé s'appuient sur de tels calculs pour vérifier les \emph{distances réglementaires} : une construction doit respecter un recul minimal par rapport à l'axe d'une route nationale ou d'une voie ferrée (comme la ligne Transcam). Le logiciel modélise la route par une droite d'équation connue, la parcelle par des points de coordonnées GPS, et calcule automatiquement les distances... avec exactement la formule que tu viens d'apprendre !
\end{savaistu}

\begin{exercice}[Pour t'entraîner]
\begin{enumerate}
\item Calculer la distance du point $M_0(3\,;-1)$ à la droite $(D) : x + 2y - 6 = 0$.
\item Soient $(D) : 3x + 4y - 1 = 0$ et $(D') : 3x + 4y + 9 = 0$. Montrer qu'elles sont parallèles et calculer $d(D\,;D')$.
\item On donne $A(0\,;4)$, $B(-3\,;-2)$, $C(5\,;-2)$. Calculer l'aire du triangle $ABC$.
\item Le point $P(1\,;1)$ est-il situé entre les droites $(D) : x + y = 0$ et $(D') : x + y - 4 = 0$ ?
\end{enumerate}
\end{exercice}

\begin{corrige}[Exercice n°1]
\begin{enumerate}
\item $d(M_0\,;D) = \dfrac{|3 + 2(-1) - 6|}{\sqrt{1^2+2^2}} = \dfrac{|-5|}{\sqrt{5}} = \dfrac{5}{\sqrt{5}} = \sqrt{5} \approx 2{,}24$.
\item Même vecteur normal $\vec{n}(3\,;4)$ : les droites sont parallèles. Un point de $(D)$ : $M_0(0\,;\tfrac14)$ car $4y = 1$. Alors $d(D\,;D') = \dfrac{|0 + 1 + 9|}{\sqrt{9+16}} = \dfrac{10}{5} = 2$.
\item $(BC)$ a pour équation $y = -2$, soit $y + 2 = 0$. Hauteur : $h_A = \dfrac{|4+2|}{\sqrt{0+1}} = 6$. Base $BC = 8$. Aire : $\mathcal{A} = \dfrac12 \times 8 \times 6 = 24$ unités d'aire.
\item On écrit $(D) : x + y = 0$ et $(D') : x + y - 4 = 0$. Pour $P(1\,;1)$ : $1+1 = 2 > 0$ et $1+1-4 = -2 < 0$. Les deux quantités sont de signes contraires : $P$ est bien situé \emph{entre} les deux droites.
\end{enumerate}
\end{corrige}

\begin{aretenir}[L'essentiel de la section]
\begin{itemize}
\item La distance de $M_0$ à $(D)$ est la longueur $M_0H$, où $H$ est le projeté orthogonal de $M_0$ sur $(D)$ : c'est le plus court chemin.
\item Pour $(D) : ax+by+c=0$ et $M_0(x_0\,;y_0)$ : $\;d(M_0\,;D) = \dfrac{|ax_0+by_0+c|}{\sqrt{a^2+b^2}}$.
\item Avec l'équation normale : $d(M_0\,;D) = |x_0\cos\theta + y_0\sin\theta - p|$.
\item Distance entre deux parallèles : un point de l'une + distance à l'autre. Hauteur d'un triangle : distance d'un sommet au côté opposé.
\end{itemize}
\end{aretenir}

\coursec{Cercle : centre, rayon et équations cartésiennes}

Dans la section précédente, nous avons appris à mesurer la distance d'un point à une droite : un outil précieux qui nous servira très bientôt, lorsqu'il s'agira de comparer la position d'une droite par rapport à un cercle. Mais avant cela, il faut d'abord apprendre à \emph{reconnaître} un cercle dans une équation, et à en extraire ses deux données essentielles : son \cle{centre} et son \cle{rayon}. C'est l'objet de cette section.

\courssub{Le cercle comme ensemble de points : l'équation réduite}

Reprenons l'intuition géométrique la plus simple. Un cercle, c'est l'ensemble des points situés à une distance fixe d'un point fixe. Le point fixe s'appelle le \cle{centre}, la distance fixe s'appelle le \cle{rayon}. Un élève qui trace un cercle au compas fait exactement cela : il pique la pointe sèche en $\Omega$ et fait tourner la mine à distance constante $R$.

\begin{definition}[Cercle]
Soit $\Omega$ un point du plan et $R$ un réel strictement positif. Le \cle{cercle} de centre $\Omega$ et de rayon $R$, noté $\mathcal{C}(\Omega, R)$, est l'ensemble des points $M$ du plan tels que :
\[ \Omega M = R. \]
\end{definition}

Plaçons-nous maintenant dans un repère orthonormé $(O ; \vec{i}, \vec{j})$, avec $\Omega(a ; b)$ et $M(x ; y)$. La distance $\Omega M$ se calcule grâce à la formule vue en géométrie analytique :
\[ \Omega M = \sqrt{(x-a)^2 + (y-b)^2}. \]
La condition $\Omega M = R$ équivaut, en élevant au carré (les deux membres sont positifs), à $\Omega M^2 = R^2$. Nous obtenons ainsi l'équation fondamentale du cercle.

\begin{propriete}[Équation réduite d'un cercle]
Dans un repère orthonormé, le cercle $\mathcal{C}(\Omega, R)$ de centre $\Omega(a ; b)$ et de rayon $R > 0$ admet pour équation, dite \cle{équation réduite} (ou forme canonique) :
\[ (x-a)^2 + (y-b)^2 = R^2. \]
Autrement dit : $M(x ; y) \in \mathcal{C}(\Omega, R) \iff (x-a)^2 + (y-b)^2 = R^2$.
\end{propriete}

\begin{popfigure}
\begin{center}
\begin{tikzpicture}[scale=1.1]
\draw[popline,->] (-1.2,0) -- (5.2,0) node[below left]{$x$};
\draw[popline,->] (0,-1.2) -- (0,4.6) node[below left]{$y$};
\coordinate (O) at (0,0);
\coordinate (Om) at (2.6,2.2);
\coordinate (M) at (4.0,3.47);
\draw[popblue,thick] (Om) circle (1.8);
\fill[popdark] (Om) circle (1.6pt) node[below left]{$\Omega(a\,;\,b)$};
\fill[poporange] (M) circle (1.6pt) node[above right]{$M(x\,;\,y)$};
\draw[poporange,thick] (Om) -- (M) node[midway, above, sloped]{$R$};
\draw[popmuted,dashed] (Om) -- (2.6,0) node[below]{$a$};
\draw[popmuted,dashed] (Om) -- (0,2.2) node[left]{$b$};
\draw[popmuted,dashed] (M) -- (4.0,0) node[below]{$x$};
\draw[popmuted,dashed] (M) -- (0,3.47) node[left]{$y$};
\fill[popdark] (O) circle (1.2pt) node[below left]{$O$};
\end{tikzpicture}
\end{center}
\legende{Le point $M(x\,;\,y)$ appartient au cercle de centre $\Omega(a\,;\,b)$ si et seulement si $\Omega M = R$, c'est-à-dire $(x-a)^2+(y-b)^2=R^2$.}
\end{popfigure}

\begin{exemplebox}[Lire centre et rayon dans une équation réduite]
L'équation $(x-2)^2 + (y+3)^2 = 16$ définit le cercle de centre $\Omega(2 ; -3)$ et de rayon $R = \sqrt{16} = 4$.

Attention à l'écriture : $(y+3)^2$ doit se lire $(y-(-3))^2$, donc l'ordonnée du centre est $-3$, et non $3$.
\end{exemplebox}

\begin{attention}[Le piège des signes dans les coordonnées du centre]
C'est l'erreur la plus fréquente au Probatoire. Dans l'écriture $(x-a)^2$, l'abscisse du centre est $\mathbf{+a}$, et non $-a$. Ainsi :
\begin{itemize}
\item $(x-5)^2$ correspond à $a = 5$ ;
\item $(x+5)^2 = (x-(-5))^2$ correspond à $a = -5$.
\end{itemize}
Retiens la phrase : \emph{« le centre, c'est ce qui annule chaque parenthèse »}. La parenthèse $(x-5)^2$ s'annule pour $x = 5$ ; la parenthèse $(y+3)^2$ s'annule pour $y = -3$. Le centre est donc le point où les deux carrés valent zéro simultanément.
\end{attention}

\courssub{Forme développée d'une équation de cercle}

Dans la pratique — et notamment dans les sujets du Probatoire —, l'équation d'un cercle n'est presque jamais donnée sous forme réduite. On la rencontre \emph{développée}. Développons l'équation réduite :
\begin{align*}
(x-a)^2 + (y-b)^2 &= R^2 \\
x^2 - 2ax + a^2 + y^2 - 2by + b^2 &= R^2 \\
x^2 + y^2 - 2ax - 2by + (a^2 + b^2 - R^2) &= 0.
\end{align*}
En posant $c = a^2 + b^2 - R^2$, on obtient la \cle{forme développée} :
\[ x^2 + y^2 - 2ax - 2by + c = 0. \]

Deux questions se posent alors naturellement :
\begin{enumerate}
\item Toute équation de cette forme définit-elle un cercle ?
\item Si oui, comment retrouver le centre et le rayon ?
\end{enumerate}

Pour y répondre, remontons le calcul dans l'autre sens. Partons de $x^2 + y^2 - 2ax - 2by + c = 0$ et reformons les carrés :
\begin{align*}
x^2 - 2ax + y^2 - 2by + c &= 0 \\
(x^2 - 2ax + a^2) + (y^2 - 2by + b^2) + c - a^2 - b^2 &= 0 \\
(x-a)^2 + (y-b)^2 &= a^2 + b^2 - c.
\end{align*}
Le membre de gauche est une somme de carrés, donc un réel \emph{positif ou nul}. Tout dépend alors du signe du membre de droite, la quantité $a^2 + b^2 - c$ :

\begin{itemize}
\item si $a^2 + b^2 - c > 0$ : on peut écrire $a^2 + b^2 - c = R^2$ avec $R = \sqrt{a^2+b^2-c}$, et l'équation est celle du cercle de centre $\Omega(a ; b)$ et de rayon $R$ ;
\item si $a^2 + b^2 - c = 0$ : l'équation devient $(x-a)^2 + (y-b)^2 = 0$. Une somme de carrés n'est nulle que si chaque carré est nul, donc $x = a$ et $y = b$ : l'ensemble des solutions est réduit au \emph{seul point} $\Omega(a ; b)$. On parle de \cle{cercle de rayon nul} ou cercle dégénéré ;
\item si $a^2 + b^2 - c < 0$ : le membre de gauche est positif ou nul, le membre de droite strictement négatif : l'égalité est impossible. L'ensemble des solutions est \emph{vide}.
\end{itemize}

\begin{propriete}[Condition d'existence d'un cercle]
Soit $(E)$ l'équation $x^2 + y^2 - 2ax - 2by + c = 0$, où $a$, $b$, $c$ sont des réels.
\begin{itemize}
\item Si $a^2 + b^2 - c > 0$, alors $(E)$ définit le cercle de centre $\Omega(a ; b)$ et de rayon $R = \sqrt{a^2 + b^2 - c}$.
\item Si $a^2 + b^2 - c = 0$, l'ensemble des points solutions est réduit au point $\Omega(a ; b)$.
\item Si $a^2 + b^2 - c < 0$, l'ensemble des points solutions est vide.
\end{itemize}
\end{propriete}

\begin{aretenir}[Les deux signatures d'une équation de cercle]
Une équation du second degré en $x$ et $y$ représente un cercle (éventuellement dégénéré) seulement si :
\begin{enumerate}
\item les coefficients de $x^2$ et $y^2$ sont \textbf{égaux} (et non nuls) — on peut alors diviser toute l'équation par ce coefficient pour se ramener à $x^2 + y^2 + \dots$ ;
\item il n'y a \textbf{pas de terme en $xy$}.
\end{enumerate}
Par exemple, $2x^2 + 2y^2 - 4x + 3 = 0$ peut représenter un cercle (diviser par $2$), mais $x^2 + 2y^2 - 4x = 0$ (coefficients différents) et $x^2 + y^2 + xy = 1$ (terme en $xy$) n'en représentent pas.
\end{aretenir}

\courssub{Méthode : retrouver centre et rayon par complétion de carrés}

Face à une équation développée, deux stratégies coexistent : appliquer directement la propriété précédente (lecture de $a$, $b$, $c$), ou refaire la mise sous forme canonique à la main. La seconde est plus sûre, car elle évite les erreurs de signe sur $a$ et $b$.

\begin{methode}[De la forme développée à la forme réduite]
Pour traiter l'équation $x^2 + y^2 + \alpha x + \beta y + \gamma = 0$ :
\begin{enumerate}
\item \textbf{Regrouper} les termes en $x$ d'un côté, les termes en $y$ de l'autre : $(x^2 + \alpha x) + (y^2 + \beta y) = -\gamma$.
\item \textbf{Compléter chaque carré} : $x^2 + \alpha x$ est le début de $\left(x + \dfrac{\alpha}{2}\right)^2$, auquel il manque $\left(\dfrac{\alpha}{2}\right)^2$. Faire de même en $y$.
\item \textbf{Ajouter les constantes manquantes des deux côtés} de l'égalité pour ne pas la rompre.
\item \textbf{Lire le résultat} : $(x-a)^2 + (y-b)^2 = k$. Si $k > 0$ : cercle de centre $(a ; b)$, rayon $\sqrt{k}$. Si $k = 0$ : un point. Si $k < 0$ : ensemble vide.
\end{enumerate}
\end{methode}

\begin{exoresolu}[Étude complète de l'équation $x^2 + y^2 - 4x + 6y - 3 = 0$]
Montrer que l'équation $x^2 + y^2 - 4x + 6y - 3 = 0$ définit un cercle $\mathcal{C}$, dont on déterminera le centre $\Omega$ et le rayon $R$.
\tcblower
\textbf{Méthode 1 : complétion de carrés.} Regroupons et complétons :
\begin{align*}
x^2 - 4x + y^2 + 6y - 3 &= 0 \\
(x^2 - 4x + 4) + (y^2 + 6y + 9) - 3 - 4 - 9 &= 0 \\
(x-2)^2 + (y+3)^2 &= 16.
\end{align*}
On reconnaît l'équation réduite d'un cercle de centre $\Omega(2 ; -3)$ et de rayon $R = \sqrt{16} = 4$.

\textbf{Méthode 2 : application directe de la propriété.} L'équation est de la forme $x^2 + y^2 - 2ax - 2by + c = 0$ avec $-2a = -4$, $-2b = 6$, $c = -3$, c'est-à-dire $a = 2$, $b = -3$, $c = -3$. On calcule :
\[ a^2 + b^2 - c = 4 + 9 - (-3) = 16 > 0. \]
L'équation définit donc bien un cercle, de centre $\Omega(2 ; -3)$ et de rayon $R = \sqrt{16} = 4$. Les deux méthodes concordent.

\textbf{Vérification (toujours profitable) :} le point $A(6 ; -3)$ vérifie $(6-2)^2 + (-3+3)^2 = 16$, donc $A \in \mathcal{C}$ ; et effectivement $\Omega A = |6-2| = 4 = R$.
\end{exoresolu}

\begin{exoresolu}[Cas dégénérés : un point, ou rien du tout]
Déterminer la nature de l'ensemble des points $M(x ; y)$ dans chacun des cas suivants :
\begin{enumerate}
\item $x^2 + y^2 - 2x + 4y + 5 = 0$ ;
\item $x^2 + y^2 + 6x - 2y + 15 = 0$.
\end{enumerate}
\tcblower
\textbf{Cas 1.} Complétons : $(x^2 - 2x + 1) + (y^2 + 4y + 4) + 5 - 1 - 4 = 0$, soit $(x-1)^2 + (y+2)^2 = 0$. Une somme de carrés est nulle si et seulement si chaque carré est nul : $x = 1$ et $y = -2$. L'ensemble cherché est réduit au \textbf{point unique} $\Omega(1 ; -2)$. (Vérification par la propriété : $a^2+b^2-c = 1 + 4 - 5 = 0$.)

\textbf{Cas 2.} Complétons : $(x^2 + 6x + 9) + (y^2 - 2y + 1) + 15 - 9 - 1 = 0$, soit $(x+3)^2 + (y-1)^2 = -5$. Le membre de gauche est positif ou nul, celui de droite vaut $-5 < 0$ : égalité impossible. L'ensemble est \textbf{vide}. (Vérification : $a^2+b^2-c = 9 + 1 - 15 = -5 < 0$.)
\end{exoresolu}

\begin{attention}[Deux réflexes de sécurité avant de conclure]
\begin{itemize}
\item Si l'équation est donnée sous la forme $2x^2 + 2y^2 - 8x + 12y - 6 = 0$, \textbf{divise d'abord toute l'équation par 2} avant d'identifier $a$, $b$, $c$. Sinon tes valeurs seront fausses.
\item Le rayon est $R = \sqrt{a^2+b^2-c}$, et non $a^2+b^2-c$. Dire « le rayon est 16 » dans l'exemple résolu ci-dessus est une erreur : le rayon est $4$.
\end{itemize}
\end{attention}

\courssub{Position relative d'une droite et d'un cercle}

Voici où la section précédente reprend du service. Imaginons un phare sur la côte de Kribi, dont le faisceau balaie la mer, et une route rectiligne au bord de l'eau. Trois situations peuvent se produire entre une droite $(D)$ et un cercle $\mathcal{C}(\Omega, R)$ :

\begin{itemize}
\item la droite \textbf{coupe} le cercle en deux points distincts : on dit que $(D)$ est \cle{sécante} au cercle ;
\item la droite \textbf{touche} le cercle en un seul point : on dit que $(D)$ est \cle{tangente} au cercle ;
\item la droite \textbf{ne rencontre pas} le cercle : on dit que $(D)$ est \cle{extérieure} au cercle.
\end{itemize}

Comment décider algébriquement ? L'idée est lumineuse : la distance $d(\Omega, D)$ du centre à la droite est la \emph{plus courte} distance entre $\Omega$ et les points de $(D)$. Si cette distance minimale est déjà supérieure au rayon, aucun point de $(D)$ ne peut être sur le cercle. Si elle vaut exactement $R$, le point de $(D)$ le plus proche de $\Omega$ est sur le cercle — et c'est le seul. Si elle est inférieure à $R$, la droite « traverse » le cercle en deux points.

\begin{propriete}[Position relative droite-cercle (admise)]
Soit $\mathcal{C}(\Omega, R)$ un cercle, $(D)$ une droite, et $d = d(\Omega, D)$ la distance du centre à la droite. Alors :
\begin{itemize}
\item $d < R \iff (D)$ est sécante à $\mathcal{C}$ (deux points d'intersection) ;
\item $d = R \iff (D)$ est tangente à $\mathcal{C}$ (un unique point de contact) ;
\item $d > R \iff (D)$ est extérieure à $\mathcal{C}$ (aucun point commun).
\end{itemize}
En particulier : $(D)$ tangente à $\mathcal{C}(\Omega, R) \iff d(\Omega, D) = R$. Ce critère sera l'outil central des sections 5 et 6.
\end{propriete}

\begin{popfigure}
\begin{center}
\begin{tikzpicture}[scale=0.86]
% Cas 1 : sécante
\begin{scope}[shift={(0,0)}]
\coordinate (O1) at (0,0);
\draw[popblue,thick] (O1) circle (1.2);
\draw[poporange,thick] (-1.7,-0.5) -- (1.7,-0.5);
\fill[popdark] (O1) circle (1.4pt) node[above]{$\Omega$};
\draw[popgreen,thick,<->] (0,0) -- (0,-0.5) node[midway,right]{$d$};
\draw[popgreen] (0,-0.5) ++(0,0.12) -- ++(0.12,0) -- ++(0,-0.12);
\fill[popdark] (-1.08,-0.5) circle (1.4pt);
\fill[popdark] (1.08,-0.5) circle (1.4pt);
\node[popdark] at (0,-1.9) {$d < R$ : sécante};
\end{scope}
% Cas 2 : tangente
\begin{scope}[shift={(4.6,0)}]
\coordinate (O2) at (0,0);
\draw[popblue,thick] (O2) circle (1.2);
\draw[poporange,thick] (-1.7,-1.2) -- (1.7,-1.2);
\fill[popdark] (O2) circle (1.4pt) node[above]{$\Omega$};
\draw[popgreen,thick,<->] (0,0) -- (0,-1.2) node[midway,right]{$d$};
\draw[popgreen] (0,-1.2) ++(0,0.12) -- ++(0.12,0) -- ++(0,-0.12);
\fill[poporange] (0,-1.2) circle (1.8pt) node[below left]{$T$};
\node[popdark] at (0,-1.9) {$d = R$ : tangente};
\end{scope}
% Cas 3 : extérieure
\begin{scope}[shift={(9.2,0)}]
\coordinate (O3) at (0,0);
\draw[popblue,thick] (O3) circle (1.2);
\draw[poporange,thick] (-1.7,-1.8) -- (1.7,-1.8);
\fill[popdark] (O3) circle (1.4pt) node[above]{$\Omega$};
\draw[popgreen,thick,<->] (0,0) -- (0,-1.8) node[midway,right]{$d$};
\draw[popgreen] (0,-1.8) ++(0,0.12) -- ++(0.12,0) -- ++(0,-0.12);
\node[popdark] at (0,-1.9) {$d > R$ : extérieure};
\end{scope}
\end{tikzpicture}
\end{center}
\legende{Les trois positions relatives d'une droite $(D)$ et d'un cercle $\mathcal{C}(\Omega, R)$, décidées par la comparaison de $d(\Omega, D)$ et de $R$.}
\end{popfigure}

\begin{methode}[Étudier la position d'une droite par rapport à un cercle]
\begin{enumerate}
\item Mettre l'équation du cercle sous forme réduite pour obtenir $\Omega(a ; b)$ et $R$.
\item Calculer $d(\Omega, D) = \dfrac{|\alpha a + \beta b + \gamma|}{\sqrt{\alpha^2 + \beta^2}}$ où $(D) : \alpha x + \beta y + \gamma = 0$ (formule de la section précédente).
\item Comparer $d$ et $R$, puis conclure : sécante, tangente ou extérieure.
\end{enumerate}
\end{methode}

\begin{exoresolu}[Position de la droite $(D) : 3x - 4y + 7 = 0$ par rapport au cercle $x^2 + y^2 - 4x + 6y - 3 = 0$]
On considère le cercle $\mathcal{C}$ d'équation $x^2 + y^2 - 4x + 6y - 3 = 0$ et la droite $(D) : 3x - 4y + 7 = 0$. Étudier la position relative de $(D)$ et $\mathcal{C}$.
\tcblower
\textbf{Étape 1.} D'après l'exercice résolu précédent, $\mathcal{C}$ a pour centre $\Omega(2 ; -3)$ et pour rayon $R = 4$.

\textbf{Étape 2.} Calculons la distance de $\Omega$ à $(D)$ :
\[ d(\Omega, D) = \frac{|3 \times 2 - 4 \times (-3) + 7|}{\sqrt{3^2 + (-4)^2}} = \frac{|6 + 12 + 7|}{\sqrt{25}} = \frac{25}{5} = 5. \]

\textbf{Étape 3.} On a $d(\Omega, D) = 5 > 4 = R$ : la droite $(D)$ est \textbf{extérieure} au cercle $\mathcal{C}$ ; elle ne le rencontre en aucun point.
\end{exoresolu}

\courssub{Intersection d'une droite et d'un cercle par résolution du système}

La comparaison de $d(\Omega, D)$ à $R$ donne le \emph{nombre} de points d'intersection, mais pas leurs coordonnées. Pour les obtenir, on résout le \cle{système} formé par les deux équations : un point d'intersection est un point dont les coordonnées vérifient \emph{à la fois} l'équation de la droite et celle du cercle.

\begin{methode}[Déterminer les points d'intersection d'une droite et d'un cercle]
\begin{enumerate}
\item Dans l'équation de la droite, exprimer une inconnue en fonction de l'autre (par exemple $y$ en fonction de $x$).
\item Substituer cette expression dans l'équation du cercle : on obtient une \textbf{équation du second degré} en une seule inconnue.
\item Résoudre cette équation. Le signe du discriminant $\Delta$ confirme la position relative : $\Delta > 0$ (deux points), $\Delta = 0$ (un point, tangence), $\Delta < 0$ (aucun point).
\item Revenir à l'équation de la droite pour calculer l'autre coordonnée de chaque point.
\end{enumerate}
\end{methode}

\begin{exoresolu}[Points d'intersection de $\mathcal{C} : x^2 + y^2 = 5$ et $(D) : y = x - 1$]
Déterminer les coordonnées des points d'intersection du cercle $\mathcal{C}$ d'équation $x^2 + y^2 = 5$ et de la droite $(D)$ d'équation $y = x - 1$.
\tcblower
\textbf{Étape 1.} Le cercle est centré en $O(0 ; 0)$, de rayon $R = \sqrt{5}$. La droite donne directement $y$ en fonction de $x$.

\textbf{Étape 2.} Substituons $y = x - 1$ dans l'équation du cercle :
\begin{align*}
x^2 + (x-1)^2 &= 5 \\
x^2 + x^2 - 2x + 1 &= 5 \\
2x^2 - 2x - 4 &= 0 \\
x^2 - x - 2 &= 0.
\end{align*}

\textbf{Étape 3.} Discriminant : $\Delta = (-1)^2 - 4 \times 1 \times (-2) = 1 + 8 = 9 > 0$. Deux solutions :
\[ x_1 = \frac{1 - 3}{2} = -1 \qquad \text{et} \qquad x_2 = \frac{1 + 3}{2} = 2. \]

\textbf{Étape 4.} Avec $y = x - 1$ : pour $x_1 = -1$, $y_1 = -2$ ; pour $x_2 = 2$, $y_2 = 1$.

\textbf{Conclusion.} La droite est sécante au cercle en deux points : $A(-1 ; -2)$ et $B(2 ; 1)$.

\textbf{Vérification.} $A$ : $(-1)^2 + (-2)^2 = 5$ et $-2 = -1 - 1$ : correct. $B$ : $2^2 + 1^2 = 5$ et $1 = 2 - 1$ : correct. On peut aussi vérifier la cohérence avec la distance : $d(O, D) = \dfrac{|0 - 0 - 1|}{\sqrt{1^2+(-1)^2}} = \dfrac{1}{\sqrt{2}} \approx \num{0,71} < \sqrt{5} \approx \num{2,24}$, ce qui confirme bien deux points d'intersection.
\end{exoresolu}

\begin{savaistu}[Des cercles partout autour de nous]
L'équation d'un cercle est l'un des outils préférés des ingénieurs en télécommunications. Pour estimer la zone de couverture d'une antenne relais MTN ou Orange installée en un point $\Omega$ et portant à $R$ kilomètres, on modélise la zone couverte par le disque de centre $\Omega$ et de rayon $R$ : un abonné situé en $M$ capte le signal si $(x_M - a)^2 + (y_M - b)^2 \leq R^2$. Les mêmes équations servent en cartographie pour le géopositionnement par trilatération : ta position est l'intersection de plusieurs cercles centrés sur des satellites ou des antennes !
\end{savaistu}

\begin{exercice}[Pour t'entraîner]
\begin{enumerate}
\item Déterminer la nature de l'ensemble des points $M(x ; y)$ vérifiant $x^2 + y^2 + 2x - 8y + 8 = 0$. S'il s'agit d'un cercle, préciser son centre et son rayon.
\item Soit $\mathcal{C}$ le cercle d'équation $(x-1)^2 + (y-2)^2 = 9$ et $(D)$ la droite d'équation $x + y - 6 = 0$. Étudier la position relative de $(D)$ et $\mathcal{C}$ par deux méthodes : comparaison de $d(\Omega, D)$ à $R$, puis résolution du système.
\item Déterminer les points d'intersection du cercle $x^2 + y^2 - 2x - 4y - 4 = 0$ avec la droite d'équation $x = 3$.
\end{enumerate}
\end{exercice}

\begin{corrige}[Exercice n°1]
\textbf{1.} Complétons les carrés : $(x^2 + 2x + 1) + (y^2 - 8y + 16) + 8 - 1 - 16 = 0$, soit $(x+1)^2 + (y-4)^2 = 9$. C'est le cercle de centre $\Omega(-1 ; 4)$ et de rayon $R = 3$.

\textbf{2.} Centre $\Omega(1 ; 2)$, rayon $R = 3$. Distance : $d(\Omega, D) = \dfrac{|1 + 2 - 6|}{\sqrt{2}} = \dfrac{3}{\sqrt{2}} \approx \num{2,12} < 3$ : la droite est sécante. Par le système : $y = 6 - x$, d'où $(x-1)^2 + (4-x)^2 = 9$, soit $2x^2 - 10x + 8 = 0$, c'est-à-dire $x^2 - 5x + 4 = 0$, d'où $x = 1$ ou $x = 4$. Points d'intersection : $(1 ; 5)$ et $(4 ; 2)$.

\textbf{3.} Avec $x = 3$ : $9 + y^2 - 6 - 4y - 4 = 0$, soit $y^2 - 4y - 1 = 0$. $\Delta = 16 + 4 = 20$, d'où $y = 2 \pm \sqrt{5}$. Points : $(3 ; 2 - \sqrt{5})$ et $(3 ; 2 + \sqrt{5})$.
\end{corrige}

\coursec{Équations paramétriques d'un cercle}

Dans la section précédente, nous avons caractérisé un cercle $\mathcal{C}(\Omega, R)$ par une \cle{équation cartésienne} $(x-a)^2+(y-b)^2=R^2$ : une relation que doivent vérifier les coordonnées $(x\,;y)$ de tout point $M$ du cercle. Cette relation est \emph{implicite} : elle ne donne pas directement $x$ et $y$, elle impose seulement une contrainte entre eux.

Imagine maintenant la grande roue de la fête foraine sur l'esplanade du stade Ahmadou Ahidjo : pour décrire la position d'une nacelle, le plus naturel n'est pas de donner son abscisse et son ordonnée séparément, mais de donner \emph{l'angle} dont la roue a tourné depuis sa position initiale. Un seul nombre — l'angle — suffit à déterminer entièrement la position. C'est exactement l'idée du \cle{paramétrage} d'un cercle : décrire chaque point du cercle à l'aide d'un unique paramètre $t$, qui sera un angle.

\courssub{Intuition et construction du paramétrage}

\textbf{Le cercle trigonométrique comme modèle}\par

Rappelons un fait fondamental vu en classe de Seconde et réinvesti en Première : sur le \cle{cercle trigonométrique} (cercle de centre $O$ et de rayon $1$), le point associé au réel $t$ a pour coordonnées $(\cos t\,;\sin t)$. Autrement dit, si $\vec{i}$ et $\vec{j}$ sont les vecteurs de base du repère orthonormé $(O\,;\vec{i},\vec{j})$, alors
\[
\overrightarrow{OM} = \cos t\,\vec{i} + \sin t\,\vec{j},
\]
où $t$ est une mesure en radians de l'angle orienté $(\vec{i},\overrightarrow{OM})$.

Pour un cercle de centre $\Omega(a\,;b)$ et de rayon $R$ quelconques, deux modifications interviennent :
\begin{itemize}
\item le rayon n'est plus $1$ mais $R$ : le vecteur $\overrightarrow{\Omega M}$ a pour norme $R$, donc $\overrightarrow{\Omega M} = R\cos t\,\vec{i} + R\sin t\,\vec{j}$ ;
\item le centre n'est plus $O$ mais $\Omega(a\,;b)$ : il faut translater, c'est-à-dire ajouter les coordonnées de $\Omega$.
\end{itemize}

\begin{popfigure}
\begin{center}
\begin{tikzpicture}[scale=1.5,>=Stealth]
% axes
\draw[->,popmuted] (-1.6,0) -- (4.6,0) node[below left]{$x$};
\draw[->,popmuted] (0,-2.2) -- (0,2.6) node[below left]{$y$};
% centre Omega
\coordinate (Om) at (1.5,0.6);
\fill[popblue] (Om) circle (1.6pt) node[below left]{$\Omega(a\,;b)$};
% cercle
\draw[popblue,thick] (Om) circle (1.5);
% point M
\coordinate (M) at ($(Om)+(50:1.5)$);
\fill[poporange] (M) circle (1.6pt) node[above right]{$M(x\,;y)$};
% rayon
\draw[poporange,thick,->] (Om) -- (M) node[midway,above left]{$R$};
% direction i depuis Omega
\draw[popdark,->] (Om) -- ($(Om)+(1.1,0)$) node[midway,below]{$\vec{i}$};
% angle t
\draw[popgreen,thick] ($(Om)+(0.55,0)$) arc (0:50:0.55);
\node[popgreen] at ($(Om)+(25:0.85)$){$t$};
% projections
\draw[dashed,popmuted] (M) -- (M |- Om);
\draw[dashed,popmuted] (M) -- (M -| Om);
% annotations projections
\draw[<->,popdark] ($(Om)+(0,-0.35)$) -- ($(M |- Om)+(0,-0.35)$) node[midway,below]{$R\cos t$};
\draw[<->,popdark] ($(M -| Om)+(-0.35,0)$) -- ($(Om)+(-0.35,0)$) node[midway,left]{$R\sin t$};
% points projections
\fill[popmuted] (M |- Om) circle (1pt);
\fill[popmuted] (M -| Om) circle (1pt);
\end{tikzpicture}
\end{center}
\legende{Le point $M$ du cercle $\mathcal{C}(\Omega,R)$ est repéré par l'angle $t$ : les projections du rayon $\overrightarrow{\Omega M}$ sur les directions de $\vec{i}$ et $\vec{j}$ donnent $R\cos t$ et $R\sin t$.}
\end{popfigure}

\courssub{Définition et justification rigoureuse}

\begin{definition}[Représentation paramétrique d'un cercle]
Le plan est muni d'un repère orthonormé $(O\,;\vec{i},\vec{j})$. Soit $\mathcal{C}$ le cercle de centre $\Omega(a\,;b)$ et de rayon $R>0$. Une \cle{représentation paramétrique} (ou \cle{équations paramétriques}) de $\mathcal{C}$ est :
\[
\begin{cases}
x = a + R\cos t \\
y = b + R\sin t
\end{cases}
\qquad t\in\mathbb{R}.
\]
Le réel $t$ est appelé le \cle{paramètre} ; il représente une mesure, en radians, de l'angle $(\vec{i},\overrightarrow{\Omega M})$.
\end{definition}

\begin{experience}[Démonstration guidée : pourquoi ces formules ?]
Soit $M(x\,;y)$ un point du plan. Notons $\mathcal{C}$ le cercle de centre $\Omega(a\,;b)$ et de rayon $R$.
\begin{enumerate}
\item \textbf{Supposons $M\in\mathcal{C}$.} Alors $\Omega M = R$, donc le vecteur $\overrightarrow{\Omega M}$ a pour norme $R$. Posons $\vec{u}=\dfrac{1}{R}\overrightarrow{\Omega M}$ : c'est un vecteur \emph{unitaire}. Or tout vecteur unitaire du plan s'écrit $\vec{u}=(\cos t\,;\sin t)$ pour un certain réel $t$ (c'est la définition même du cosinus et du sinus d'un nombre réel, lue sur le cercle trigonométrique). On en déduit :
\[
\overrightarrow{\Omega M} = R\vec{u} = \big(R\cos t\,;\,R\sin t\big).
\]
Comme $\overrightarrow{\Omega M}=(x-a\,;\,y-b)$, l'unicité des coordonnées dans la base $(\vec{i},\vec{j})$ donne $x-a=R\cos t$ et $y-b=R\sin t$.
\item \textbf{Réciproquement}, supposons qu'il existe un réel $t$ tel que $x=a+R\cos t$ et $y=b+R\sin t$. Alors
\[
\Omega M^2 = (x-a)^2+(y-b)^2 = R^2\cos^2 t + R^2\sin^2 t = R^2(\cos^2 t+\sin^2 t)=R^2,
\]
donc $\Omega M = R$, c'est-à-dire $M\in\mathcal{C}$.
\end{enumerate}
L'équivalence est établie : le système paramétrique décrit \emph{exactement} le cercle, ni plus ni moins.
\end{experience}

\begin{aretenir}[L'équivalence fondamentale]
\[
M(x\,;y)\in\mathcal{C}\big(\Omega(a\,;b),R\big) \iff \exists\, t\in\mathbb{R},\ 
\begin{cases}
x = a + R\cos t \\
y = b + R\sin t
\end{cases}
\]
Le paramètre $t$ est une mesure de l'angle $(\vec{i},\overrightarrow{\Omega M})$, exprimée en \cle{radians}. Quand $t$ parcourt $[0\,;2\pi[$, le point $M$ décrit le cercle entier, une seule fois, dans le sens trigonométrique.
\end{aretenir}

\begin{exemplebox}[Un premier paramétrage chiffré]
Soit $\mathcal{C}$ le cercle de centre $\Omega(1\,;-2)$ et de rayon $3$. Une représentation paramétrique de $\mathcal{C}$ est :
\[
\begin{cases}
x = 1 + 3\cos t \\
y = -2 + 3\sin t
\end{cases}
\qquad t\in\mathbb{R}.
\]
Vérifions sur quelques valeurs remarquables :
\begin{itemize}
\item pour $t=0$ : $x=1+3=4$, $y=-2$. Le point $A(4\,;-2)$ est bien sur le cercle : $(4-1)^2+(-2+2)^2=9=3^2$ ;
\item pour $t=\dfrac{\pi}{2}$ : $x=1+0=1$, $y=-2+3=1$. Le point $B(1\,;1)$ est le point « le plus haut » du cercle ;
\item pour $t=\pi$ : $x=1-3=-2$, $y=-2$. Le point $C(-2\,;-2)$ est le point « le plus à gauche ».
\end{itemize}
\end{exemplebox}

\begin{attention}[Le paramètre $t$ n'est pas une longueur !]
L'erreur la plus fréquente est de confondre le paramètre $t$ avec une distance ou une longueur d'arc. \textbf{$t$ est un angle, exprimé en radians.} Ainsi :
\begin{itemize}
\item dire « $t=3$ » signifie que l'angle $(\vec{i},\overrightarrow{\Omega M})$ mesure $3$ radians (environ $171,9^\circ$), et non que l'on a parcouru $3$ unités de longueur sur le cercle ;
\item la longueur de l'arc parcouru entre les paramètres $t_1$ et $t_2$ vaut $R\,|t_2-t_1|$ (formule de la longueur d'un arc : $\ell = R\theta$), et non $|t_2-t_1|$ ;
\item pensez à régler votre calculatrice en mode \textbf{radians} pour évaluer $\cos t$ et $\sin t$.
\end{itemize}
Autre piège : ne pas intervertir $\cos$ et $\sin$. C'est bien $x=a+R\cos t$ et $y=b+R\sin t$ : le cosinus se lit sur l'axe des abscisses, le sinus sur l'axe des ordonnées, comme sur le cercle trigonométrique.
\end{attention}

\courssub{Du paramétrique au cartésien : éliminer le paramètre}

\begin{propriete}[Équivalence entre équation cartésienne et système paramétrique]
Soit $\Omega(a\,;b)$ et $R>0$. Les deux descriptions suivantes définissent le même ensemble de points :
\begin{itemize}
\item l'équation cartésienne : $(x-a)^2+(y-b)^2=R^2$ ;
\item le système paramétrique : $\begin{cases} x=a+R\cos t \\ y=b+R\sin t \end{cases}$, $t\in\mathbb{R}$.
\end{itemize}
\end{propriete}

La démonstration a été menée dans l'activité guidée ci-dessus : le point clé est l'identité $\cos^2 t+\sin^2 t=1$, qui garantit que tout point donné par le système paramétrique vérifie l'équation cartésienne, et l'existence, pour tout point du cercle, d'un angle $t$ associé au vecteur unitaire $\frac{1}{R}\overrightarrow{\Omega M}$.

\begin{methode}[Éliminer le paramètre $t$ pour retrouver l'équation cartésienne]
On part d'un système $\begin{cases} x = a + R\cos t \\ y = b + R\sin t \end{cases}$ (éventuellement donné sous une forme à réorganiser).
\begin{enumerate}
\item \textbf{Isoler} les termes trigonométriques : $x-a = R\cos t$ et $y-b = R\sin t$.
\item \textbf{Élever au carré} les deux égalités : $(x-a)^2=R^2\cos^2 t$ et $(y-b)^2=R^2\sin^2 t$.
\item \textbf{Additionner} membre à membre : $(x-a)^2+(y-b)^2 = R^2(\cos^2 t+\sin^2 t)$.
\item \textbf{Conclure} grâce à $\cos^2 t+\sin^2 t=1$ : $(x-a)^2+(y-b)^2=R^2$. Le paramètre $t$ a disparu : on dit qu'on l'a \emph{éliminé}.
\end{enumerate}
\end{methode}

\begin{exoresolu}[Retrouver centre et rayon à partir d'un paramétrage]
Énoncé. Le plan est muni d'un repère orthonormé. On considère l'ensemble $\mathcal{E}$ des points $M(x\,;y)$ tels que
\[
\begin{cases}
x = -3 + 5\cos t \\
y = 4 + 5\sin t
\end{cases}
\qquad t\in\mathbb{R}.
\]
Déterminer la nature de $\mathcal{E}$, son centre et son rayon, puis donner son équation cartésienne.
\tcblower
Solution détaillée. On isole les termes trigonométriques :
\[
x+3 = 5\cos t \qquad\text{et}\qquad y-4 = 5\sin t.
\]
On élève au carré et on additionne :
\begin{align*}
(x+3)^2+(y-4)^2 &= 25\cos^2 t + 25\sin^2 t \\
&= 25\big(\cos^2 t+\sin^2 t\big) \\
&= 25.
\end{align*}
L'ensemble $\mathcal{E}$ est donc le cercle de centre $\Omega(-3\,;4)$ et de rayon $R=\sqrt{25}=5$, d'équation cartésienne $(x+3)^2+(y-4)^2=25$.
\end{exoresolu}

\courssub{Du cartésien au paramétrique : introduire le paramètre}

La démarche inverse est tout aussi exigible. Partons de l'équation cartésienne $(x-a)^2+(y-b)^2=R^2$.

\begin{methode}[Paramétrer un cercle donné par son équation cartésienne]
\begin{enumerate}
\item \textbf{Identifier} le centre $\Omega(a\,;b)$ et le rayon $R$ dans l'équation $(x-a)^2+(y-b)^2=R^2$ (si l'équation est développée, on la ramène d'abord à la forme canonique par complétion des carrés, comme on l'a vu à la section précédente).
\item \textbf{Réécrire} l'équation sous la forme
\[
\left(\frac{x-a}{R}\right)^2+\left(\frac{y-b}{R}\right)^2=1.
\]
\item \textbf{Reconnaître} une somme de deux carrés égale à $1$ : or les nombres $\cos t$ et $\sin t$ sont précisément caractérisés par $\cos^2 t+\sin^2 t=1$. On pose donc
\[
\frac{x-a}{R}=\cos t \qquad\text{et}\qquad \frac{y-b}{R}=\sin t.
\]
\item \textbf{Conclure} : $\begin{cases} x=a+R\cos t \\ y=b+R\sin t \end{cases}$, $t\in\mathbb{R}$.
\end{enumerate}
\end{methode}

\begin{exoresolu}[Paramétrer un cercle donné par une équation développée]
Énoncé. Soit $\mathcal{C}$ l'ensemble des points $M(x\,;y)$ tels que $x^2+y^2-2x+6y+6=0$. Montrer que $\mathcal{C}$ est un cercle, puis en donner une représentation paramétrique.
\tcblower
Solution détaillée. On complète les carrés :
\begin{align*}
x^2-2x &= (x-1)^2-1,\\
y^2+6y &= (y+3)^2-9.
\end{align*}
L'équation devient :
\[
(x-1)^2-1+(y+3)^2-9+6=0 \iff (x-1)^2+(y+3)^2=4.
\]
$\mathcal{C}$ est le cercle de centre $\Omega(1\,;-3)$ et de rayon $R=2$. Une représentation paramétrique est donc :
\[
\begin{cases}
x = 1 + 2\cos t \\
y = -3 + 2\sin t
\end{cases}
\qquad t\in\mathbb{R}.
\]
\end{exoresolu}

\courssub{Restriction du paramètre : arcs et demi-cercles}

Jusqu'ici, $t$ parcourait $\mathbb{R}$ tout entier : le point $M(t)$ tourne indéfiniment autour de $\Omega$, repassant périodiquement par les mêmes positions (les fonctions cosinus et sinus sont $2\pi$-périodiques). Pour décrire le cercle \emph{une seule fois}, il suffit de prendre $t$ dans un intervalle d'amplitude $2\pi$, par exemple $[0\,;2\pi[$ ou $[-\pi\,;\pi[$.

Mais on peut aussi restreindre davantage $t$ pour ne décrire qu'une \emph{partie} du cercle :

\begin{aretenir}[Paramètre restreint et portion de cercle]
Pour le cercle paramétré par $x=a+R\cos t$, $y=b+R\sin t$ :
\begin{itemize}
\item $t\in[0\,;2\pi[$ : le cercle \textbf{entier}, parcouru une fois dans le sens trigonométrique ;
\item $t\in[0\,;\pi]$ : le \textbf{demi-cercle supérieur} (car $\sin t\geq 0$, donc $y\geq b$) ;
\item $t\in[\pi\,;2\pi]$ : le \textbf{demi-cercle inférieur} ($y\leq b$) ;
\item $t\in\left[-\dfrac{\pi}{2}\,;\dfrac{\pi}{2}\right]$ : le \textbf{demi-cercle de droite} (car $\cos t\geq 0$, donc $x\geq a$) ;
\item $t\in[t_1\,;t_2]$ : l'\textbf{arc de cercle} joignant le point de paramètre $t_1$ au point de paramètre $t_2$, dans le sens trigonométrique.
\end{itemize}
\end{aretenir}

\begin{exemplebox}[Un quart de cercle]
Considérons le cercle de centre $\Omega(1\,;-2)$ et de rayon $3$, paramétré par $x=1+3\cos t$, $y=-2+3\sin t$. Pour $t\in\left[0\,;\dfrac{\pi}{2}\right]$, le point $M(t)$ décrit le quart de cercle joignant $A(4\,;-2)$ (paramètre $t=0$) à $B(1\,;1)$ (paramètre $t=\dfrac{\pi}{2}$), en passant par exemple par le point de paramètre $t=\dfrac{\pi}{4}$ :
\[
x = 1+3\cos\frac{\pi}{4} = 1+\frac{3\sqrt{2}}{2}\approx 3{,}12, \qquad
y = -2+3\sin\frac{\pi}{4} = -2+\frac{3\sqrt{2}}{2}\approx 0{,}12.
\]
Cet arc est situé dans la région $x\geq 1$, $y\geq -2$ : le « quart nord-est » du cercle.
\end{exemplebox}

\begin{popfigure}
\begin{center}
\begin{tikzpicture}
\begin{axis}[
  width=0.85\linewidth,
  axis equal image,
  axis lines=middle,
  xlabel={$x$}, ylabel={$y$},
  xmin=-3.5, xmax=5.5, ymin=-6, ymax=2.5,
  xtick={-2,1,4}, ytick={-5,-2,1},
  grid=both, grid style={popline!40},
  every axis plot/.append style={thick}
]
\addplot[popcurve,domain=0:360,samples=200] ({1+3*cos(x)},{-2+3*sin(x)});
\addplot[only marks,mark=*,mark size=2pt,poporange] coordinates {(4,-2) (1,1) (-2,-2) (1,-5)};
\node[poporange,anchor=south west] at (axis cs:4,-2) {$M(0)$};
\node[poporange,anchor=south west] at (axis cs:1,1) {$M\!\left(\frac{\pi}{2}\right)$};
\node[poporange,anchor=south east] at (axis cs:-2,-2) {$M(\pi)$};
\node[poporange,anchor=north west] at (axis cs:1,-5) {$M\!\left(\frac{3\pi}{2}\right)$};
\addplot[only marks,mark=*,mark size=1.8pt,popblue] coordinates {(1,-2)};
\node[popblue,anchor=south east] at (axis cs:1,-2) {$\Omega$};
\end{axis}
\end{tikzpicture}
\end{center}
\legende{Le cercle $x=1+3\cos t$, $y=-2+3\sin t$ tracé en mode paramétrique, avec les points correspondant aux valeurs remarquables du paramètre.}
\end{popfigure}

\begin{attention}[Un même point, une infinité de paramètres]
À cause de la $2\pi$-périodicité de $\cos$ et $\sin$, le paramètre d'un point donné n'est \textbf{jamais unique} : $M(t)$ et $M(t+2k\pi)$, avec $k\in\mathbb{Z}$, désignent le même point. Par exemple, sur le cercle de centre $\Omega(1\,;-2)$ et de rayon $3$, le point $A(4\,;-2)$ correspond à $t=0$, mais aussi à $t=2\pi$, $t=-4\pi$, etc.

Conséquence pratique : quand on résout une équation du type « trouver $t$ tel que $M(t)=M_0$ », on se place dans un intervalle d'amplitude $2\pi$ (souvent $[0\,;2\pi[$) et on y cherche \emph{toutes} les solutions. Deux situations typiques :
\begin{itemize}
\item l'équation peut avoir \textbf{une seule} solution dans $[0\,;2\pi[$ : par exemple $M_0(1\,;-5)$ impose $\cos t=0$ et $\sin t=-1$, d'où $t=\dfrac{3\pi}{2}$ uniquement ;
\item elle peut en avoir \textbf{deux} si l'on ne fixe qu'une seule coordonnée : par exemple $\cos t=\dfrac12$ donne $t=\dfrac{\pi}{3}$ ou $t=\dfrac{5\pi}{3}$ dans $[0\,;2\pi[$, qui correspondent à deux points \emph{distincts} du cercle (symétriques par rapport à l'axe horizontal passant par $\Omega$).
\end{itemize}
Retenir : pour identifier \emph{un} point du cercle, il faut connaître à la fois $\cos t$ \textbf{et} $\sin t$.
\end{attention}

\courssub{Ouverture : le paramétrage comme modèle de mouvement}

\begin{savaistu}[Vers la Terminale : le mouvement circulaire uniforme]
Le paramétrage du cercle prend tout son sens en physique et en Terminale C/E. Si un mobile parcourt le cercle de centre $\Omega$ et de rayon $R$ à \emph{vitesse angulaire constante} $\omega$ (en rad/s), sa position à l'instant $t$ (en secondes) est :
\[
\begin{cases}
x(t) = a + R\cos(\omega t + \varphi) \\
y(t) = b + R\sin(\omega t + \varphi)
\end{cases}
\]
où $\varphi$ est la position angulaire de départ. C'est le modèle du \cle{mouvement circulaire uniforme} : la grande roue évoquée en introduction, la rotation d'une éolienne à Guider, l'orbite (quasi circulaire) d'un satellite... Le paramètre devient alors le \emph{temps}, et l'angle parcouru est $\theta=\omega t+\varphi$. Retiens l'idée essentielle : \textbf{un paramétrage, c'est un mode d'emploi pour fabriquer la courbe point par point}, ce que font exactement les logiciels de géométrie et les traceurs de courbes.
\end{savaistu}

\begin{exercice}[Pour s'entraîner]
\begin{enumerate}
\item Donner une représentation paramétrique du cercle de centre $\Omega(-1\,;5)$ et de rayon $\sqrt{7}$.
\item L'ensemble des points $M(x\,;y)$ tels que $x=2-4\cos t$, $y=3-4\sin t$, $t\in\mathbb{R}$, est-il un cercle ? Préciser son centre et son rayon. (Attention aux signes !)
\item On considère le cercle $\mathcal{C}$ : $x^2+y^2+4x-2y-11=0$. Déterminer une représentation paramétrique de $\mathcal{C}$, puis les points de $\mathcal{C}$ correspondant à $t=\dfrac{\pi}{6}$ et $t=-\dfrac{\pi}{2}$.
\item Sur le cercle trigonométrique (centre $O$, rayon $1$), quel arc est décrit lorsque $t$ parcourt $\left[\dfrac{\pi}{6}\,;\dfrac{5\pi}{6}\right]$ ? Donner les coordonnées exactes de ses extrémités.
\end{enumerate}
\end{exercice}

\begin{corrige}[Exercice]
\begin{enumerate}
\item $\begin{cases} x=-1+\sqrt{7}\cos t \\ y=5+\sqrt{7}\sin t \end{cases}$, $t\in\mathbb{R}$.
\item On a $x-2=-4\cos t$ et $y-3=-4\sin t$, donc $(x-2)^2+(y-3)^2=16\cos^2t+16\sin^2t=16$. C'est le cercle de centre $\Omega(2\,;3)$ et de rayon $4$. (Les signes « $-$ » ne changent ni le centre ni le rayon : ils correspondent à un sens de parcours horaire et un décalage de paramètre.)
\item Complétion des carrés : $x^2+4x=(x+2)^2-4$ et $y^2-2y=(y-1)^2-1$. L'équation devient $(x+2)^2+(y-1)^2=16$ : cercle de centre $\Omega(-2\,;1)$, rayon $4$. Paramétrage : $x=-2+4\cos t$, $y=1+4\sin t$. Pour $t=\dfrac{\pi}{6}$ : $M\big(-2+2\sqrt{3}\,;\,3\big)$. Pour $t=-\dfrac{\pi}{2}$ : $M(-2\,;-3)$.
\item C'est l'arc joignant $A\left(\dfrac{\sqrt3}{2}\,;\dfrac12\right)$ à $B\left(-\dfrac{\sqrt3}{2}\,;\dfrac12\right)$ en passant par le point $(0\,;1)$ : l'arc supérieur d'angle au centre $\dfrac{2\pi}{3}$.
\end{enumerate}
\end{corrige}

En résumé, le paramétrage d'un cercle transforme une contrainte géométrique (« être à distance $R$ de $\Omega$ ») en une \emph{recette de fabrication} : choisir un angle $t$, calculer $a+R\cos t$ et $b+R\sin t$, et placer le point. Cette double lecture — cartésienne et paramétrique — sera précieuse dans la suite de la leçon, notamment pour l'étude des tangentes, puis en Terminale pour les mouvements et les nombres complexes.

\coursec{Tangente à un cercle en un de ses points}

Dans la section précédente, nous avons appris à décrire un cercle par des équations paramétriques : chaque point du cercle est repéré par un angle $\theta$, ce qui nous permet de «~parcourir~» le cercle. Nous allons maintenant nous intéresser à une droite très particulière liée au cercle : la \cle{tangente} en un de ses points. Tu connais déjà cette notion depuis la classe de Troisième, de manière géométrique ; ici, grâce au produit scalaire et aux vecteurs normaux étudiés au début de cette leçon, nous allons en donner une \emph{équation cartésienne} que tu pourras écrire en quelques lignes. C'est une question classique du Probatoire.

\courssub{Rappel géométrique : la tangente est perpendiculaire au rayon}

\begin{definition}[Tangente à un cercle]
Soit $\mathcal{C}$ un cercle de centre $\Omega$ et de rayon $R$, et $M_0$ un point de $\mathcal{C}$. La \cle{tangente} au cercle $\mathcal{C}$ au point $M_0$ est l'unique droite passant par $M_0$ et \textbf{perpendiculaire au rayon} $[\Omega M_0]$. Le point $M_0$ est appelé \emph{point de contact}.
\end{definition}

L'intuition est simple : la tangente «~frôle~» le cercle en un seul point, sans jamais le traverser. Imagine une roue de vélo posée sur le sol de la cour de récréation : le sol ne touche la roue qu'en un seul point, et le rayon qui aboutit à ce point de contact est vertical, donc perpendiculaire au sol horizontal. Cette propriété géométrique, démontrée au collège, est le point de départ de tout ce qui suit.

\begin{propriete}[Caractérisation de la tangente]
Soit $\mathcal{C}$ le cercle de centre $\Omega$ et de rayon $R$, et $M_0$ un point de $\mathcal{C}$. Une droite $\mathcal{T}$ passant par $M_0$ est la tangente à $\mathcal{C}$ en $M_0$ si et seulement si le vecteur $\overrightarrow{\Omega M_0}$ est un \cle{vecteur normal} de $\mathcal{T}$. Autrement dit, pour tout point $M$ du plan :
\[ M \in \mathcal{T} \iff \overrightarrow{\Omega M_0} \cdot \overrightarrow{M_0 M} = 0. \]
\end{propriete}

En effet, dire que $\mathcal{T}$ est perpendiculaire à la droite $(\Omega M_0)$, c'est dire exactement que pour tout point $M$ de $\mathcal{T}$, les vecteurs $\overrightarrow{\Omega M_0}$ et $\overrightarrow{M_0 M}$ sont orthogonaux, donc que leur produit scalaire est nul. Cette caractérisation vectorielle est la clé : elle transforme une propriété géométrique (perpendicularité) en une équation que l'on peut développer en coordonnées.

\begin{popfigure}
\begin{center}
\begin{tikzpicture}[scale=1.1,>=Stealth]
% Définition des constantes
\def\ang{40}
\def\rad{2}
\def\tlen{1.8}
% Cercle
\draw[popblue,thick] (0,0) circle (\rad);
% Centre
\fill[popdark] (0,0) circle (1.6pt) node[below left=2pt] {$\Omega(a\,;\,b)$};
% Point M0
\coordinate (M0) at (\ang:\rad);
\fill[poporange] (M0) circle (1.8pt) node[above right=1pt] {$M_0(x_0\,;\,y_0)$};
% Rayon
\draw[popblue,thick,->] (0,0) -- (M0) node[midway,below left=1pt] {$\overrightarrow{\Omega M_0}$};
% Tangente : direction perpendiculaire au rayon (angle + 90)
\draw[poporange,thick] ($(M0)+(\ang+90:\tlen)$) -- ($(M0)+(\ang-90:\tlen)$) node[below right=1pt] {$\mathcal{T}$};
% Point M sur la tangente
\coordinate (M) at ($(M0)+(\ang+90:1.1)$);
\fill[popgreen] (M) circle (1.5pt) node[above left=2pt] {$M(x\,;\,y)$};
% Vecteur M0M
\draw[popgreen,thick,->] (M0) -- (M) node[midway,above=2pt,sloped] {$\overrightarrow{M_0 M}$};
% Angle droit (marque standard)
\draw[popdark] ($(M0)+(\ang:0.35)$) -- ($(M0)+(\ang:0.35)+(\ang-90:0.35)$) -- ($(M0)+(\ang-90:0.35)$);
\end{tikzpicture}
\end{center}
\legende{La tangente $\mathcal{T}$ en $M_0$ est perpendiculaire au rayon $[\Omega M_0]$ : le vecteur $\overrightarrow{\Omega M_0}$ est normal à $\mathcal{T}$, et pour tout point $M$ de $\mathcal{T}$, on a $\overrightarrow{\Omega M_0}\cdot\overrightarrow{M_0 M}=0$.}
\end{popfigure}

\courssub{Théorème : équation cartésienne de la tangente en un point du cercle}

\begin{propriete}[Équation de la tangente en un point du cercle]
Le plan est muni d'un repère orthonormé $(O\,;\vec{\imath},\vec{\jmath})$. Soit $\mathcal{C}$ le cercle de centre $\Omega(a\,;\,b)$ et de rayon $R$, et $M_0(x_0\,;\,y_0)$ un point \textbf{appartenant à} $\mathcal{C}$. La tangente à $\mathcal{C}$ en $M_0$ a pour équation :
\[ (x_0 - a)(x - x_0) + (y_0 - b)(y - y_0) = 0. \]
Cette démonstration est formatrice et peut être demandée au Probatoire : elle repose uniquement sur le produit scalaire.
\end{propriete}

\begin{experience}[Démonstration guidée du théorème]
On se place dans un repère orthonormé. Soit $M(x\,;\,y)$ un point quelconque du plan.
\begin{enumerate}
\item \textbf{Coordonnées des vecteurs.} On a $\overrightarrow{\Omega M_0}\begin{pmatrix} x_0 - a \\ y_0 - b \end{pmatrix}$ et $\overrightarrow{M_0 M}\begin{pmatrix} x - x_0 \\ y - y_0 \end{pmatrix}$.
\item \textbf{Traduction de l'appartenance.} D'après la caractérisation précédente :
\[ M \in \mathcal{T} \iff \overrightarrow{\Omega M_0} \cdot \overrightarrow{M_0 M} = 0. \]
\item \textbf{Calcul du produit scalaire.} Dans un repère orthonormé, le produit scalaire de deux vecteurs $\vec{u}\begin{pmatrix} X \\ Y \end{pmatrix}$ et $\vec{v}\begin{pmatrix} X' \\ Y' \end{pmatrix}$ vaut $XX' + YY'$. On obtient donc :
\[ M \in \mathcal{T} \iff (x_0 - a)(x - x_0) + (y_0 - b)(y - y_0) = 0. \]
\item \textbf{Conclusion.} Cette relation est bien une équation cartésienne de droite (elle est de la forme $ux + vy + w = 0$ avec $(u\,;\,v) = (x_0 - a\,;\,y_0 - b) \neq (0\,;\,0)$ car $M_0 \neq \Omega$), et c'est celle de la tangente cherchée.
\end{enumerate}
\end{experience}

\begin{aretenir}[La formule à retenir]
Tangente en $M_0(x_0\,;\,y_0)$ au cercle de centre $\Omega(a\,;\,b)$ :
\[ \boxed{(x_0 - a)(x - x_0) + (y_0 - b)(y - y_0) = 0} \]
On reconnaît le produit scalaire $\overrightarrow{\Omega M_0}\cdot\overrightarrow{M_0 M} = 0$ développé en coordonnées. Plutôt que d'apprendre la formule par cœur, retiens la \emph{démarche} : le vecteur $\overrightarrow{\Omega M_0}$ est normal à la tangente.
\end{aretenir}

\begin{methode}[Déterminer une équation de la tangente en un point du cercle]
Pour déterminer une équation cartésienne de la tangente en $M_0(x_0\,;\,y_0)$ au cercle $\mathcal{C}$ de centre $\Omega(a\,;\,b)$ et de rayon $R$ :
\begin{enumerate}
\item \textbf{Vérifier que $M_0$ appartient au cercle} : calculer $\Omega M_0^2 = (x_0 - a)^2 + (y_0 - b)^2$ et contrôler que $\Omega M_0^2 = R^2$. Cette étape est \textbf{obligatoire} et doit apparaître sur ta copie.
\item \textbf{Écrire la condition d'orthogonalité} : $M(x\,;\,y) \in \mathcal{T} \iff \overrightarrow{\Omega M_0}\cdot\overrightarrow{M_0 M} = 0$.
\item \textbf{Développer} : $(x_0 - a)(x - x_0) + (y_0 - b)(y - y_0) = 0$, puis réduire à une équation de la forme $ux + vy + w = 0$.
\end{enumerate}
\end{methode}

\begin{exemplebox}[Tangente en $M_0(4\,;\,2)$ au cercle de centre $\Omega(1\,;\,-2)$ et de rayon $5$]
\textbf{Étape 1 : vérification.} Calculons :
\[ \Omega M_0^2 = (4 - 1)^2 + (2 - (-2))^2 = 3^2 + 4^2 = 9 + 16 = 25 = R^2. \]
Le point $M_0(4\,;\,2)$ est bien sur le cercle.

\textbf{Étape 2 : orthogonalité.} $\overrightarrow{\Omega M_0}\begin{pmatrix} 3 \\ 4 \end{pmatrix}$ est normal à la tangente.

\textbf{Étape 3 : développement.}
\begin{align*}
M(x\,;\,y) \in \mathcal{T} &\iff 3(x - 4) + 4(y - 2) = 0 \\
&\iff 3x + 4y - 12 - 8 = 0 \\
&\iff 3x + 4y - 20 = 0.
\end{align*}
La tangente en $M_0(4\,;\,2)$ a pour équation $3x + 4y - 20 = 0$.
\end{exemplebox}

\begin{attention}[Piège n°1 : appliquer la formule sans vérifier que $M_0$ est sur le cercle]
L'erreur la plus fréquente au Probatoire consiste à écrire directement $(x_0 - a)(x - x_0) + (y_0 - b)(y - y_0) = 0$ pour un point $M_0$ quelconque. Or si $M_0$ n'appartient \textbf{pas} au cercle, la droite obtenue n'est \emph{pas} tangente au cercle : la formule donne alors une droite passant par $M_0$ et perpendiculaire à $(\Omega M_0)$, mais qui peut être sécante ou extérieure au cercle. \textbf{Sur ta copie, la vérification $\Omega M_0 = R$ est la première ligne de la solution.} Si l'énoncé te donne un point qui n'est pas sur le cercle, c'est que tu es dans la situation de la section suivante (tangentes issues d'un point extérieur), et la méthode sera différente.
\end{attention}

\courssub{Exercice résolu et vérification par la distance}

\begin{exoresolu}[Tangente en $M_0(4\,;\,1)$ au cercle de centre $\Omega(1\,;\,-2)$ et de rayon $3\sqrt{2}$]
Soit $\mathcal{C}$ le cercle de centre $\Omega(1\,;\,-2)$ et de rayon $R = 3\sqrt{2}$, et le point $M_0(4\,;\,1)$.
\begin{enumerate}
\item Vérifier que $M_0$ appartient au cercle $\mathcal{C}$.
\item Déterminer une équation cartésienne de la tangente $\mathcal{T}$ à $\mathcal{C}$ en $M_0$.
\item Vérifier que la distance de $\Omega$ à la droite $\mathcal{T}$ est égale au rayon.
\end{enumerate}
\tcblower
\textbf{1. Appartenance au cercle.}
\[ \Omega M_0^2 = (4 - 1)^2 + (1 - (-2))^2 = 3^2 + 3^2 = 9 + 9 = 18. \]
Or $R^2 = (3\sqrt{2})^2 = 18$. Comme $\Omega M_0^2 = R^2$, le point $M_0(4\,;\,1)$ appartient bien au cercle $\mathcal{C}$.

\textbf{2. Équation de la tangente.} On a $\overrightarrow{\Omega M_0}\begin{pmatrix} 4 - 1 \\ 1 - (-2) \end{pmatrix} = \begin{pmatrix} 3 \\ 3 \end{pmatrix}$. Pour $M(x\,;\,y)$ :
\begin{align*}
M \in \mathcal{T} &\iff \overrightarrow{\Omega M_0}\cdot\overrightarrow{M_0 M} = 0 \\
&\iff 3(x - 4) + 3(y - 1) = 0 \\
&\iff 3x + 3y - 12 - 3 = 0 \\
&\iff 3x + 3y - 15 = 0 \\
&\iff x + y - 5 = 0 \quad \text{(en divisant par } 3\text{)}.
\end{align*}
La tangente $\mathcal{T}$ a pour équation $x + y - 5 = 0$.

\textbf{3. Vérification par la distance.} Utilisons la formule de la distance d'un point à une droite vue dans cette leçon :
\[ d(\Omega\,,\,\mathcal{T}) = \frac{|1 \times 1 + 1 \times (-2) - 5|}{\sqrt{1^2 + 1^2}} = \frac{|1 - 2 - 5|}{\sqrt{2}} = \frac{6}{\sqrt{2}} = \frac{6\sqrt{2}}{2} = 3\sqrt{2} = R. \]
La distance du centre à la tangente est bien égale au rayon : notre résultat est confirmé. Cette vérification est un excellent réflexe de contrôle.
\end{exoresolu}

\begin{propriete}[Lien avec la distance d'un point à une droite]
Une droite $\mathcal{D}$ est tangente au cercle $\mathcal{C}(\Omega\,,\,R)$ si et seulement si $d(\Omega\,,\,\mathcal{D}) = R$. En particulier, pour contrôler une équation de tangente obtenue par le produit scalaire, il suffit de vérifier que la distance du centre à cette droite vaut exactement $R$.
\end{propriete}

Ce critère est cohérent avec l'intuition : si $d(\Omega\,,\,\mathcal{D}) < R$, la droite traverse le cercle (elle est sécante) ; si $d(\Omega\,,\,\mathcal{D}) > R$, elle passe «~à côté~» sans le toucher ; et le cas limite $d(\Omega\,,\,\mathcal{D}) = R$ correspond exactement au contact en un unique point.

\courssub{Cas particuliers : tangentes horizontales et verticales}

Lorsque le rayon $[\Omega M_0]$ est vertical ou horizontal, la tangente est respectivement horizontale ou verticale, et son équation se lit \emph{directement} sur le centre et le rayon, sans calcul de produit scalaire.

\begin{propriete}[Tangentes horizontales et verticales]
Soit $\mathcal{C}$ le cercle de centre $\Omega(a\,;\,b)$ et de rayon $R$.
\begin{itemize}
\item Aux points $(a\,;\,b + R)$ et $(a\,;\,b - R)$ (points «~haut~» et «~bas~» du cercle), les tangentes sont \textbf{horizontales}, d'équations $y = b + R$ et $y = b - R$.
\item Aux points $(a + R\,;\,b)$ et $(a - R\,;\,b)$ (points «~droite~» et «~gauche~» du cercle), les tangentes sont \textbf{verticales}, d'équations $x = a + R$ et $x = a - R$.
\end{itemize}
\end{propriete}

Vérifions la cohérence avec la formule générale. Au point $M_0(a\,;\,b+R)$, on a $\overrightarrow{\Omega M_0}\begin{pmatrix} 0 \\ R \end{pmatrix}$, donc :
\[ 0 \times (x - a) + R\,(y - b - R) = 0 \iff y = b + R. \]
La formule générale redonne bien la lecture directe.

\begin{exemplebox}[Lecture directe]
Soit $\mathcal{C}$ le cercle de centre $\Omega(2\,;\,-1)$ et de rayon $R = 3$.
\begin{itemize}
\item Tangente horizontale en $(2\,;\,2)$ : $y = -1 + 3$, c'est-à-dire $y = 2$.
\item Tangente horizontale en $(2\,;\,-4)$ : $y = -4$.
\item Tangente verticale en $(5\,;\,-1)$ : $x = 5$.
\item Tangente verticale en $(-1\,;\,-1)$ : $x = -1$.
\end{itemize}
\end{exemplebox}

\begin{attention}[Deux pièges de rédaction à éviter]
\begin{itemize}
\item \textbf{Ne pas simplifier trop vite sans raison.} Dans l'équation $3(x - 4) + 3(y - 1) = 0$, on peut diviser par $3$ pour obtenir $x + y - 5 = 0$ : c'est la même droite, et la forme simplifiée est préférable. Mais ne divise jamais par une expression contenant $x$ ou $y$.
\item \textbf{Confondre vecteur normal et vecteur directeur.} Le vecteur $\overrightarrow{\Omega M_0}$ est \emph{normal} à la tangente, pas directeur. Un vecteur directeur de la tangente est $\begin{pmatrix} -(y_0 - b) \\ x_0 - a \end{pmatrix}$. Si l'on te demande un vecteur directeur, ne donne pas $\overrightarrow{\Omega M_0}$ !
\end{itemize}
\end{attention}

\begin{savaistu}[La roue du benskin et la tangente]
À Yaoundé comme à Douala, les motos-taxis — les fameux \emph{benskins} — sillonnent les quartiers. Regarde la roue avant d'une moto posée sur la route : elle ne touche le bitume qu'en \textbf{un seul point}. Mathématiquement, la route est \emph{tangente} à la roue, et le rayon de la roue aboutissant au point de contact est perpendiculaire à la route. C'est exactement le théorème de cette section ! Ce même principe explique pourquoi, sur une route horizontale, l'axe vertical passant par le point de contact porte le poids de la moto de manière optimale : la force de réaction du sol est dirigée le long du rayon, vers le moyeu. Les ingénieurs qui conçoivent les suspensions utilisent précisément cette géométrie.
\end{savaistu}

\begin{exercice}[Pour s'entraîner]
Soit $\mathcal{C}$ le cercle de centre $\Omega(-1\,;\,2)$ et de rayon $R = 5$.
\begin{enumerate}
\item Vérifier que les points $A(2\,;\,6)$ et $B(3\,;\,-1)$ appartiennent à $\mathcal{C}$.
\item Déterminer une équation cartésienne de la tangente à $\mathcal{C}$ en $A$, puis de la tangente à $\mathcal{C}$ en $B$.
\item Vérifier, pour la tangente en $A$, que $d(\Omega\,,\,\mathcal{T}_A) = R$.
\item Donner, sans calcul, les équations des tangentes horizontales et verticales à $\mathcal{C}$.
\end{enumerate}
\end{exercice}

\begin{corrige}[Exercice : tangentes au cercle de centre $\Omega(-1\,;\,2)$]
\textbf{1.} $\Omega A^2 = (2-(-1))^2 + (6-2)^2 = 9 + 16 = 25 = R^2$, donc $A \in \mathcal{C}$. De même $\Omega B^2 = (3-(-1))^2 + (-1-2)^2 = 16 + 9 = 25 = R^2$, donc $B \in \mathcal{C}$.

\textbf{2.} Tangente en $A$ : $\overrightarrow{\Omega A}\begin{pmatrix} 3 \\ 4 \end{pmatrix}$, donc $3(x - 2) + 4(y - 6) = 0$, soit $3x + 4y - 30 = 0$.

Tangente en $B$ : $\overrightarrow{\Omega B}\begin{pmatrix} 4 \\ -3 \end{pmatrix}$, donc $4(x - 3) - 3(y + 1) = 0$, soit $4x - 3y - 15 = 0$.

\textbf{3.} $d(\Omega\,,\,\mathcal{T}_A) = \dfrac{|3(-1) + 4(2) - 30|}{\sqrt{9 + 16}} = \dfrac{|-3 + 8 - 30|}{5} = \dfrac{25}{5} = 5 = R$. La vérification est concluante.

\textbf{4.} Tangentes horizontales : $y = 2 + 5 = 7$ et $y = 2 - 5 = -3$. Tangentes verticales : $x = -1 + 5 = 4$ et $x = -1 - 5 = -6$.
\end{corrige}

\begin{aretenir}[L'essentiel de la section]
\begin{itemize}
\item La tangente en $M_0$ au cercle $\mathcal{C}(\Omega\,,\,R)$ est la droite passant par $M_0$ et perpendiculaire au rayon $[\Omega M_0]$ : le vecteur $\overrightarrow{\Omega M_0}$ est un \cle{vecteur normal} de la tangente.
\item Équation : $(x_0 - a)(x - x_0) + (y_0 - b)(y - y_0) = 0$, obtenue en développant $\overrightarrow{\Omega M_0}\cdot\overrightarrow{M_0 M} = 0$.
\item \textbf{Toujours vérifier d'abord} que $M_0$ appartient au cercle ($\Omega M_0 = R$).
\item Contrôle possible : $d(\Omega\,,\,\mathcal{T}) = R$.
\item Tangentes horizontales $y = b \pm R$ et verticales $x = a \pm R$ : lecture directe.
\end{itemize}
\end{aretenir}

Nous savons désormais déterminer la tangente en un point \emph{du} cercle. Mais que se passe-t-il lorsque le point donné est \emph{extérieur} au cercle, comme un projecteur qui éclaire un ballon depuis l'extérieur ? Il existe alors \textbf{deux} tangentes passant par ce point : c'est l'objet de la section suivante.

\coursec{Tangentes à un cercle passant par un point extérieur}

Dans la section précédente, nous avons appris à écrire l'équation de la tangente à un cercle \emph{en un point du cercle} : il suffisait d'exploiter la perpendicularité entre le rayon et la tangente. Mais une question naturelle se pose alors, et elle est très fréquente au Probatoire : \textbf{si le point $A$ n'est pas sur le cercle, peut-on encore mener des tangentes au cercle passant par $A$ ? Combien ? Comment trouver leurs équations ?} Toute cette section est consacrée à cette question, dans le cas le plus riche : celui où $A$ est \emph{extérieur} au cercle.

\courssub{Combien de tangentes peut-on mener d'un point donné ?}

Commençons par une expérience de pensée. Place une pièce de 100 FCFA sur la table (notre cercle), et une gomme à côté (notre point $A$). Avec une règle posée contre la gomme, cherche les positions où la règle \emph{effleure} la pièce sans la couper. Tu constates qu'il y a exactement \textbf{deux} telles positions : une qui touche la pièce « par le haut », une « par le bas ». Si maintenant la gomme est posée \emph{sur} le bord de la pièce, une seule position convient. Et si la gomme était à l'intérieur du disque ? Aucune droite passant par elle ne pourrait effleurer le cercle sans le traverser.

\begin{propriete}[Nombre de tangentes issues d'un point (admise)]
Soit $(\mathcal{C})$ un cercle de centre $\Omega$ et de rayon $R$, et $A$ un point du plan.
\begin{itemize}
\item Si $A$ est \textbf{extérieur} au cercle ($\Omega A > R$), il passe par $A$ \cle{exactement deux tangentes} au cercle.
\item Si $A$ est \textbf{sur} le cercle ($\Omega A = R$), il passe par $A$ \cle{exactement une tangente} (étudiée à la section précédente).
\item Si $A$ est \textbf{intérieur} au cercle ($\Omega A < R$), il n'existe \cle{aucune tangente} au cercle passant par $A$.
\end{itemize}
De plus, dans le premier cas, si $T_1$ et $T_2$ sont les points de contact, alors les deux segments tangents ont la même longueur :
\[ AT_1 = AT_2 = \sqrt{\Omega A^2 - R^2}. \]
\end{propriete}

\begin{experience}[Pourquoi $AT_1 = AT_2 = \sqrt{\Omega A^2 - R^2}$ ?]
Justifions géométriquement la formule des longueurs égales.
\begin{enumerate}
\item Soit $T$ un point de contact d'une tangente issue de $A$. Par la propriété fondamentale de la tangente, le rayon $[\Omega T]$ est perpendiculaire à la tangente $(AT)$ : le triangle $\Omega T A$ est donc \textbf{rectangle en $T$}.
\item Appliquons le théorème de Pythagore dans ce triangle :
\[ \Omega A^2 = \Omega T^2 + AT^2 = R^2 + AT^2. \]
\item On en tire $AT^2 = \Omega A^2 - R^2$, d'où $AT = \sqrt{\Omega A^2 - R^2}$ (cette quantité est bien définie car $A$ extérieur signifie $\Omega A > R$, donc $\Omega A^2 - R^2 > 0$).
\item Ce calcul ne dépend pas du point de contact choisi : $AT_1 = AT_2$, et la figure est symétrique par rapport à la droite $(\Omega A)$.
\end{enumerate}
\end{experience}

\begin{popfigure}
\begin{center}
\begin{tikzpicture}[scale=1.1]
\coordinate (O) at (0,0);
\coordinate (A) at (5,0);
\coordinate (T1) at (0.8,1.833);
\coordinate (T2) at (0.8,-1.833);
% cercle
\draw[popblue, thick] (O) circle (2);
% tangentes
\draw[poporange, thick] (A) -- (T1);
\draw[poporange, thick] (A) -- (T2);
% rayons
\draw[popgreen, thick] (O) -- (T1);
\draw[popgreen, thick] (O) -- (T2);
% axe de symétrie
\draw[popmuted, dashed] (-1.2,0) -- (6,0);
% angles droits en T1 et T2
\coordinate (P1) at ($(T1)!0.16!(O)$);
\coordinate (Q1) at ($(T1)!0.16!(A)$);
\draw[popdark] (P1) -- ($(P1)+(Q1)-(T1)$) -- (Q1);
\coordinate (P2) at ($(T2)!0.16!(O)$);
\coordinate (Q2) at ($(T2)!0.16!(A)$);
\draw[popdark] (P2) -- ($(P2)+(Q2)-(T2)$) -- (Q2);
% points
\fill[popblue] (O) circle (1.5pt) node[below left] {$\Omega$};
\fill[popink] (A) circle (1.5pt) node[right] {$A$};
\fill[poporange] (T1) circle (1.5pt) node[above left] {$T_1$};
\fill[poporange] (T2) circle (1.5pt) node[below left] {$T_2$};
\node[popgreen] at (0.15,1.15) {\small $R$};
\node[poporange] at (3.1,1.25) {\small $AT_1$};
\node[poporange] at (3.1,-1.25) {\small $AT_2$};
\end{tikzpicture}
\end{center}
\legende{D'un point extérieur $A$ partent exactement deux tangentes ; les rayons de contact sont perpendiculaires aux tangentes et $AT_1 = AT_2$.}
\end{popfigure}

\courssub{Méthode 1 : la méthode algébrique (condition $d(\Omega, D) = R$)}

L'idée est lumineuse : une droite $(D)$ passant par $A$ est tangente au cercle si et seulement si elle \emph{effleure} le cercle, c'est-à-dire si et seulement si la distance du centre $\Omega$ à $(D)$ est exactement égale au rayon $R$. Nous savons calculer cette distance depuis la section 2 : tout est donc en place.

\begin{methode}[Trouver les tangentes à un cercle issues d'un point extérieur]
Soit $(\mathcal{C})$ de centre $\Omega$ et de rayon $R$, et $A(x_A\,;\,y_A)$ un point extérieur.
\begin{enumerate}
\item \textbf{Vérifier} que $A$ est extérieur : calculer $\Omega A$ et constater que $\Omega A > R$.
\item \textbf{Écrire} une droite $(D)$ passant par $A$ et de pente $m$ :
\[ (D)\ :\ y - y_A = m(x - x_A), \]
puis la mettre sous forme cartésienne $mx - y + (y_A - m x_A) = 0$.
\item \textbf{Traduire la tangence} : écrire $d(\Omega, D) = R$, ce qui donne une équation en $m$ (avec une valeur absolue).
\item \textbf{Résoudre} : en élevant au carré, on obtient une équation du second degré en $m$, qui fournit en général deux pentes $m_1$ et $m_2$.
\item \textbf{Ne pas oublier} de tester séparément la \cle{droite verticale} $x = x_A$, qui n'est pas de la forme $y = mx + p$.
\item \textbf{Conclure} en donnant les équations cartésiennes des tangentes, et vérifier.
\end{enumerate}
\end{methode}

\begin{attention}[Le piège classique : la droite verticale]
L'équation $y - y_A = m(x - x_A)$ décrit \textbf{toutes} les droites passant par $A$... \emph{sauf une} : la droite verticale $x = x_A$, qui n'a pas de pente. Il arrive qu'elle soit tangente au cercle ! Si tu résous l'équation en $m$ et que tu ne trouves qu'\textbf{une seule solution} alors que $A$ est extérieur, c'est presque toujours le signe que la deuxième tangente est la droite verticale. Teste-la systématiquement : calcule $d(\Omega, x = x_A) = |x_\Omega - x_A|$ et compare à $R$.
\end{attention}

\begin{exoresolu}[Tangentes au cercle $x^2 + y^2 = 5$ issues de $A(5\,;\,0)$]
Le plan est muni d'un repère orthonormé. Soit $(\mathcal{C})$ le cercle d'équation $x^2 + y^2 = 5$ et le point $A(5\,;\,0)$. Déterminer les équations des tangentes à $(\mathcal{C})$ passant par $A$.
\tcblower
\textbf{Solution détaillée.}

\textbf{Étape 1 : position de $A$.} Le cercle a pour centre $\Omega(0\,;\,0)$ et rayon $R = \sqrt{5}$. On a $\Omega A = 5 > \sqrt{5}$ : $A$ est extérieur, il y a donc exactement deux tangentes.

\textbf{Étape 2 : droites de pente $m$ passant par $A$.}
\[ (D)\ :\ y - 0 = m(x - 5) \quad\Longleftrightarrow\quad mx - y - 5m = 0. \]

\textbf{Étape 3 : condition de tangence.}
\[ d(\Omega, D) = \frac{|m \cdot 0 - 0 - 5m|}{\sqrt{m^2 + 1}} = \frac{|5m|}{\sqrt{m^2+1}} = \sqrt{5}. \]

\textbf{Étape 4 : résolution.} En élevant au carré (les deux membres sont positifs, l'équivalence est conservée) :
\begin{align*}
\frac{25m^2}{m^2 + 1} &= 5 \\
25m^2 &= 5m^2 + 5 \\
20m^2 &= 5 \\
m^2 &= \frac{1}{4} \quad\Longrightarrow\quad m = \frac{1}{2} \ \text{ou}\ m = -\frac{1}{2}.
\end{align*}

\textbf{Étape 5 : la droite verticale.} La droite $x = 5$ passe par $A$ ; sa distance à $\Omega$ vaut $|5 - 0| = 5 \neq \sqrt{5}$ : elle n'est \emph{pas} tangente. Tout est cohérent : nous avons bien nos deux tangentes.

\textbf{Étape 6 : conclusion.} Avec $m = \dfrac{1}{2}$ : $y = \dfrac{1}{2}(x-5)$, soit $x - 2y - 5 = 0$. Avec $m = -\dfrac{1}{2}$ : $y = -\dfrac{1}{2}(x-5)$, soit $x + 2y - 5 = 0$.

\textbf{Vérification.} Pour la première : $d(\Omega, D) = \dfrac{|-5|}{\sqrt{1+4}} = \dfrac{5}{\sqrt{5}} = \sqrt{5} = R$. Parfait. Les points de contact sont les projetés orthogonaux de $\Omega$ sur ces droites : $T_1(1\,;\,-2)$ et $T_2(1\,;\,2)$, qui vérifient bien $1^2 + (\pm 2)^2 = 5$. Enfin, $AT_1 = AT_2 = \sqrt{\Omega A^2 - R^2} = \sqrt{25 - 5} = \sqrt{20} = 2\sqrt{5}$, conformément à la propriété.
\end{exoresolu}

\courssub{Méthode 2 : la méthode géométrique (cercle de diamètre $[\Omega A]$)}

Voici un regard différent, plus géométrique, qui éclaire la structure de la figure et qui est parfois demandé comme complément. Cherchons directement les \emph{points de contact} $T$ au lieu des droites.

Un point $T$ est un point de contact d'une tangente issue de $A$ si et seulement si :
\begin{itemize}
\item $T$ appartient au cercle $(\mathcal{C})$ : $\Omega T = R$ ;
\item la tangente en $T$ passe par $A$, c'est-à-dire $\overrightarrow{\Omega T} \perp \overrightarrow{AT}$.
\end{itemize}
Or dire que l'angle $\widehat{\Omega T A}$ est droit, c'est dire, par la réciproque du théorème de l'angle inscrit (un triangle rectangle est inscriptible dans le cercle de diamètre son hypoténuse), que $T$ appartient au \cle{cercle de diamètre $[\Omega A]$}.

\begin{propriete}[{Points de contact et cercle de diamètre $[\Omega A]$}]
Soit $(\mathcal{C})$ un cercle de centre $\Omega$ et $A$ un point extérieur. Les points de contact $T_1$ et $T_2$ des tangentes issues de $A$ sont les points d'intersection de $(\mathcal{C})$ avec le \cle{cercle de diamètre $[\Omega A]$}.
\end{propriete}

\begin{popfigure}
\begin{center}
\begin{tikzpicture}[scale=1.1]
\coordinate (O) at (0,0);
\coordinate (A) at (5,0);
\coordinate (I) at (2.5,0);
\coordinate (T1) at (1,2);
\coordinate (T2) at (1,-2);
% cercle de diamètre [OA]
\draw[poppurple, thick, dashed] (I) circle (2.5);
% cercle initial
\draw[popblue, thick] (O) circle (2.236);
% tangentes
\draw[poporange, thick] (A) -- (T1);
\draw[poporange, thick] (A) -- (T2);
% rayons
\draw[popgreen, thick] (O) -- (T1);
\draw[popgreen, thick] (O) -- (T2);
% points
\fill[popblue] (O) circle (1.5pt) node[below left] {$\Omega$};
\fill[popink] (A) circle (1.5pt) node[right] {$A$};
\fill[poppurple] (I) circle (1.5pt) node[below] {\small $I$};
\fill[poporange] (T1) circle (1.5pt) node[above] {$T_1$};
\fill[poporange] (T2) circle (1.5pt) node[below] {$T_2$};
\node[poppurple] at (4.1,1.9) {\small diamètre $[\Omega A]$};
\node[popblue] at (-1.55,-1.7) {\small $(\mathcal{C})$};
\end{tikzpicture}
\end{center}
\legende{Le cercle de diamètre $[\Omega A]$ (pointillés violets, centre $I$ milieu de $[\Omega A]$) coupe $(\mathcal{C})$ exactement aux points de contact $T_1$ et $T_2$.}
\end{popfigure}

\begin{exemplebox}[Retrouver les points de contact par la méthode géométrique]
Reprenons $(\mathcal{C}) : x^2 + y^2 = 5$ et $A(5\,;\,0)$. Le milieu de $[\Omega A]$ est $I\left(\dfrac{5}{2}\,;\,0\right)$ et $\Omega A = 5$, donc le cercle de diamètre $[\Omega A]$ a pour rayon $\dfrac{5}{2}$ et pour équation :
\[ \left(x - \tfrac{5}{2}\right)^2 + y^2 = \tfrac{25}{4} \quad\Longleftrightarrow\quad x^2 + y^2 - 5x = 0. \]
En soustrayant membre à membre avec $x^2 + y^2 = 5$, il reste $5 - 5x = 0$, soit $x = 1$. En reportant dans l'équation de $(\mathcal{C})$ : $y^2 = 4$, donc $y = \pm 2$. On retrouve $T_1(1\,;\,2)$ et $T_2(1\,;\,-2)$, et les tangentes sont les droites $(AT_1)$ et $(AT_2)$. Élégant, non ?
\end{exemplebox}

\courssub{Synthèse : quelle méthode choisir ?}

\begin{aretenir}[L'essentiel de la section]
\begin{itemize}
\item D'un point \textbf{extérieur} $A$ partent \cle{exactement deux tangentes} au cercle ; les segments tangents vérifient $AT_1 = AT_2 = \sqrt{\Omega A^2 - R^2}$.
\item \textbf{Méthode algébrique} (la plus opérationnelle au Probatoire) : poser $(D) : y - y_A = m(x - x_A)$, écrire $d(\Omega, D) = R$, résoudre l'équation du second degré en $m$, \textbf{puis tester la droite verticale} $x = x_A$.
\item \textbf{Méthode géométrique} : les points de contact sont les intersections de $(\mathcal{C})$ avec le cercle de diamètre $[\Omega A]$.
\item Toujours \textbf{vérifier} : la distance de $\Omega$ à chaque droite trouvée doit valoir exactement $R$.
\end{itemize}
Choix pratique : si l'énoncé demande les \emph{équations des tangentes}, la méthode algébrique est directe ; s'il demande les \emph{points de contact} ou une construction, la méthode géométrique est idéale. Les deux se contrôlent mutuellement.
\end{aretenir}

\begin{exercice}[À toi de jouer]
Soit $(\mathcal{C})$ le cercle d'équation $x^2 + y^2 = 1$ et $A(0\,;\,3)$.
\begin{enumerate}
\item Justifier qu'il existe deux tangentes à $(\mathcal{C})$ passant par $A$.
\item Déterminer leurs équations par la méthode algébrique.
\item La droite verticale passant par $A$ est-elle tangente ?
\end{enumerate}
\end{exercice}

\begin{corrige}[Exercice 1]
\begin{enumerate}
\item Le centre est $\Omega(0\,;\,0)$ et $R = 1$. On a $\Omega A = 3 > R$ : $A$ est extérieur, donc il existe exactement deux tangentes.
\item $(D) : y = mx + 3$, soit $mx - y + 3 = 0$. Condition de tangence : $\dfrac{|3|}{\sqrt{m^2+1}} = 1$, d'où $m^2 + 1 = 9$, soit $m^2 = 8$ et $m = \pm 2\sqrt{2}$. Tangentes : $2\sqrt{2}\,x - y + 3 = 0$ et $2\sqrt{2}\,x + y - 3 = 0$.
\item La verticale $x = 0$ est à distance $0$ de $\Omega$ : elle passe par le centre et coupe le cercle en deux points ; elle n'est pas tangente.
\end{enumerate}
\end{corrige}

\begin{savaistu}[Des tangentes dans les stades et sur les routes]
Quand un projecteur du stade d'Olembe éclaire la pelouse, les rayons qui « rasent » un obstacle circulaire (un poteau, un ballon posé au sol vu de loin) suivent exactement les \textbf{tangentes} au cercle issues de la source lumineuse : la zone d'ombre est délimitée par elles. En balistique, pour qu'un projectile lancé depuis $A$ frôle un obstacle circulaire sans le toucher, sa trajectoire initiale doit suivre une tangente au cercle. Les ingénieurs en éclairage et en télécommunications (pour le dégagement des faisceaux hertziens de CAMTEL autour des collines) utilisent quotidiennement ces calculs de tangentes issues d'un point extérieur !
\end{savaistu}

\newpage

\coursec{Activité d'intégration}

\coursec{Géométrie analytique du plan : Configurations et transformations élémentaires — Activité d'intégration}

\courssub{Objectifs de l'activité}

À l'issue de cette activité d'intégration, tu seras capable de :
\begin{itemize}
\item mobiliser simultanément plusieurs savoir-faire de la leçon dans une situation authentique d'aménagement du territoire ;
\item écrire l'\cle{équation normale} d'une droite et l'utiliser pour lire directement une distance à l'origine ;
\item calculer la \cle{distance d'un point à une droite} et interpréter le résultat pour vérifier la conformité d'un projet à une norme technique ;
\item déterminer les équations des \cle{tangentes à un cercle} passant par un point extérieur ;
\item donner les \cle{équations paramétriques} d'un cercle pour modéliser un parcours circulaire ;
\item traduire des grandeurs géométriques en grandeurs physiques (conversion d'unités) et rédiger une conclusion argumentée.
\end{itemize}

\courssub{Situation}

Dans le village de \textbf{Guidiguis}, dans la région de l'Extrême-Nord du Cameroun, le comité de développement local a fait réaliser un \textbf{forage} destiné à l'approvisionnement en eau potable des habitants. L'ingénieur du génie rural a établi un plan cadastral du site, rapporté à un repère orthonormé $(O\,;\vec{\imath},\vec{\jmath})$ du plan, dans lequel \textbf{l'unité de longueur vaut 10 mètres} sur le terrain.

Sur ce plan :
\begin{itemize}
\item le pourtour circulaire du forage est modélisé par le cercle $\mathscr{C}$ de centre $\Omega(4\,;3)$ et de rayon $R = 2$ (soit 20 mètres en réalité) ;
\item une route rectiligne projetée, destinée à relier le village à la route nationale, est modélisée par la droite $(\mathscr{R})$ d'équation cartésienne :
\[ 3x + 4y - 34 = 0 \; ; \]
\item l'origine $O$ du repère correspond à la case du chef du village, point de référence du plan cadastral ;
\item un technicien de la commune, placé au point $A(8\,;7)$, doit installer deux clôtures rectilignes de protection, chacune \textbf{tangente} au pourtour du forage ;
\item enfin, la coopérative agricole souhaite programmer un système d'\textbf{arrosage circulaire} qui parcourt exactement le pourtour du forage : il faut donc en fournir une description paramétrique.
\end{itemize}

La réglementation en vigueur exige que toute route passe à \textbf{au moins 15 mètres} (soit $1{,}5$ unité du plan) du centre d'un forage, afin de protéger la nappe et les installations des vibrations et des pollutions.

Le comité de développement te consulte, toi l'élève de Première, en tant qu'expert mathématique : la route projetée est-elle conforme ? Quelles clôtures le technicien doit-il installer ? Comment programmer l'arrosage ?

\courssub{Consignes}

\begin{enumerate}
\item \textbf{Équation normale de la route.}
\begin{enumerate}
\item Justifier que le vecteur $\vec{n}(3\,;4)$ est un vecteur normal à la droite $(\mathscr{R})$, puis calculer sa norme.
\item Écrire l'équation normale de $(\mathscr{R})$.
\item En déduire, sans nouveau calcul, la distance de l'origine $O$ (la case du chef) à la route. Convertir le résultat en mètres.
\end{enumerate}
\item \textbf{Conformité de la route.}
\begin{enumerate}
\item Calculer la distance du centre $\Omega(4\,;3)$ du forage à la droite $(\mathscr{R})$.
\item Convertir cette distance en mètres, la comparer à la norme réglementaire de 15 mètres, et conclure par une phrase argumentée destinée au comité de développement.
\end{enumerate}
\item \textbf{Installation des clôtures.}
\begin{enumerate}
\item Vérifier que le point $A(8\,;7)$ est extérieur au cercle $\mathscr{C}$ (on pourra calculer $\Omega A$ et comparer au rayon).
\item On cherche les droites passant par $A$ sous la forme $y = m(x-8) + 7$. Justifier cette écriture, puis traduire la condition de tangence par une équation d'inconnue $m$.
\item Résoudre cette équation et donner les équations cartésiennes des deux clôtures.
\end{enumerate}
\item \textbf{Programmation de l'arrosage.} Donner un système d'équations paramétriques du cercle $\mathscr{C}$, en précisant le paramètre et son intervalle de variation pour un tour complet.
\end{enumerate}

\courssub{Ressources à mobiliser}

\begin{aretenir}[Boîte à outils de la leçon]
\begin{itemize}
\item \textbf{Vecteur normal :} si $(d) : ax + by + c = 0$, alors $\vec{n}(a\,;b)$ est un vecteur normal à $(d)$.
\item \textbf{Équation normale :} en divisant par $\|\vec{n}\| = \sqrt{a^2+b^2}$, on obtient
\[ \frac{a}{\sqrt{a^2+b^2}}\,x + \frac{b}{\sqrt{a^2+b^2}}\,y + \frac{c}{\sqrt{a^2+b^2}} = 0, \]
dans laquelle la distance de $O$ à $(d)$ se lit directement : $d(O,(d)) = \dfrac{|c|}{\sqrt{a^2+b^2}}$.
\item \textbf{Distance d'un point à une droite :} pour $M_0(x_0\,;y_0)$ et $(d) : ax+by+c=0$,
\[ d(M_0,(d)) = \frac{|ax_0 + by_0 + c|}{\sqrt{a^2+b^2}}. \]
\item \textbf{Tangence :} une droite $(d)$ est tangente au cercle de centre $\Omega$ et de rayon $R$ si et seulement si $d(\Omega,(d)) = R$.
\item \textbf{Point extérieur :} $A$ est extérieur au cercle si et seulement si $\Omega A > R$ ; il passe alors exactement deux tangentes par $A$.
\item \textbf{Équations paramétriques du cercle} de centre $\Omega(a\,;b)$ et de rayon $R$ :
\[ \begin{cases} x = a + R\cos t \\ y = b + R\sin t \end{cases}, \quad t \in [0\,;2\pi]. \]
\end{itemize}
\end{aretenir}

\courssub{Démarche et corrigé commenté}

\begin{exoresolu}[Aménagement autour du forage de Guidiguis]
\textbf{Question 1.} Écrire l'équation normale de la route $(\mathscr{R}) : 3x+4y-34=0$ et en déduire sa distance à l'origine.

\textbf{Question 2.} La route respecte-t-elle la norme des 15 mètres par rapport au centre $\Omega(4\,;3)$ du forage ?

\textbf{Question 3.} Déterminer les deux clôtures tangentes au forage passant par $A(8\,;7)$.

\textbf{Question 4.} Donner les équations paramétriques du pourtour du forage.
\tcblower
\textbf{1.a)} La droite $(\mathscr{R})$ a pour équation $3x+4y-34=0$ : un vecteur normal est donc $\vec{n}(3\,;4)$. Sa norme est :
\[ \|\vec{n}\| = \sqrt{3^2+4^2} = \sqrt{25} = 5. \]

\textbf{1.b)} On divise tous les coefficients par $5$ ; l'équation normale de $(\mathscr{R})$ est :
\[ \frac{3}{5}x + \frac{4}{5}y - \frac{34}{5} = 0. \]

\textbf{1.c)} Dans une équation normale, le terme constant (en valeur absolue) est exactement la distance de l'origine à la droite :
\[ d(O,(\mathscr{R})) = \frac{|-34|}{5} = \frac{34}{5} = 6{,}8 \text{ unités}. \]
Comme une unité vaut 10 mètres, la case du chef se trouve à $\SI{68}{m}$ de la route projetée.

\textbf{2.a)} Appliquons la formule de la distance d'un point à une droite avec $\Omega(4\,;3)$ :
\[ d(\Omega,(\mathscr{R})) = \frac{|3\times 4 + 4\times 3 - 34|}{\sqrt{3^2+4^2}} = \frac{|12+12-34|}{5} = \frac{|-10|}{5} = 2 \text{ unités}. \]

\textbf{2.b)} Deux unités correspondent à $2 \times 10 = \SI{20}{m}$. Or $20 > 15$ : la distance est \textbf{supérieure} à la norme réglementaire. On rédige la conclusion pour le comité : \emph{« La route projetée passe à 20 mètres du centre du forage ; elle est donc conforme à la réglementation qui exige au moins 15 mètres. Le projet peut être maintenu en l'état. »}

Remarque au passage : $d(\Omega,(\mathscr{R})) = 2 = R$ signifie aussi que la route est \textbf{tangente} au cercle $\mathscr{C}$ : elle effleure exactement le pourtour du forage en un unique point. Le comité pourrait d'ailleurs juger prudent de reculer légèrement le tracé pour créer une marge de sécurité.

\textbf{3.a)} Calculons la distance de $A$ au centre :
\[ \Omega A = \sqrt{(8-4)^2 + (7-3)^2} = \sqrt{16+16} = \sqrt{32} = 4\sqrt{2} \approx 5{,}66. \]
Comme $\Omega A = 4\sqrt{2} > 2 = R$, le point $A$ est \textbf{extérieur} au cercle : il existe bien exactement deux tangentes passant par $A$.

\textbf{3.b)} Toute droite non verticale passant par $A(8\,;7)$ a une équation de la forme $y - 7 = m(x-8)$, c'est-à-dire :
\[ mx - y + 7 - 8m = 0. \]
(La droite verticale $x = 8$ donne $d(\Omega, x=8) = |8-4| = 4 \neq 2$ : elle n'est pas tangente, l'hypothèse « non verticale » ne fait donc perdre aucune solution.)

La condition de tangence s'écrit $d(\Omega, (d)) = R$, soit :
\[ \frac{|4m - 3 + 7 - 8m|}{\sqrt{m^2+1}} = 2
\quad\Longleftrightarrow\quad
\frac{|4 - 4m|}{\sqrt{m^2+1}} = 2. \]

\textbf{3.c)} Élevons au carré (les deux membres sont positifs) :
\begin{align*}
(4-4m)^2 &= 4(m^2+1) \\
16 - 32m + 16m^2 &= 4m^2 + 4 \\
12m^2 - 32m + 12 &= 0 \\
3m^2 - 8m + 3 &= 0.
\end{align*}
Le discriminant est $\Delta = 64 - 36 = 28 = (2\sqrt{7})^2$, d'où :
\[ m_1 = \frac{8 - 2\sqrt{7}}{6} = \frac{4-\sqrt{7}}{3}, \qquad m_2 = \frac{4+\sqrt{7}}{3}. \]
Les deux clôtures ont donc pour équations :
\[ (\mathscr{C}_1) : \ y = \frac{4-\sqrt{7}}{3}(x-8) + 7, \qquad (\mathscr{C}_2) : \ y = \frac{4+\sqrt{7}}{3}(x-8) + 7. \]
Numériquement, $m_1 \approx 0{,}45$ et $m_2 \approx 2{,}22$ : une clôture « douce » et une clôture « raide », ce qui correspond bien à la géométrie de la situation.

\textbf{4.} Le cercle de centre $\Omega(4\,;3)$ et de rayon $2$ admet pour équations paramétriques :
\[ \begin{cases} x = 4 + 2\cos t \\ y = 3 + 2\sin t \end{cases}, \quad t \in [0\,;2\pi], \]
où $t$ est l'angle (en radians) repérant la position du jet d'arrosage autour du centre. Lorsque $t$ décrit $[0\,;2\pi]$, le point $(x\,;y)$ effectue exactement un tour complet du pourtour : le programme d'arrosage est prêt.
\end{exoresolu}

\begin{attention}[Les pièges de cette activité]
\begin{itemize}
\item \textbf{Oublier la valeur absolue} dans la formule de distance : $3\times 4 + 4\times 3 - 34 = -10$, et une distance ne peut pas être négative.
\item \textbf{Confondre unités du plan et mètres} : tous les résultats doivent être multipliés par 10 avant d'être comparés aux normes exprimées en mètres.
\item \textbf{Diviser par $\sqrt{a^2+b^2}$ et non par $a^2+b^2$} dans l'équation normale et la distance.
\item \textbf{Oublier de vérifier que $A$ est extérieur} avant de chercher les tangentes : si $A$ était intérieur, l'équation en $m$ n'aurait aucune solution.
\item \textbf{Perdre la tangente verticale} : lorsqu'on écrit $y = m(x-x_A)+y_A$, il faut toujours tester séparément la droite verticale $x = x_A$.
\end{itemize}
\end{attention}

\courssub{Bilan de l'activité}

Cette activité t'a permis de mettre en évidence plusieurs idées essentielles de la géométrie analytique :
\begin{itemize}
\item l'\textbf{équation normale} n'est pas qu'un exercice de style : elle transforme une équation de droite en un instrument de mesure, où la distance à l'origine se lit immédiatement ;
\item la \textbf{formule de la distance d'un point à une droite} est un outil de décision concrète : elle a permis de trancher, par un calcul objectif, la question de la conformité d'un projet routier à une norme de sécurité ;
\item la condition de tangence $d(\Omega,(d)) = R$ convertit un problème géométrique (construire des tangentes) en une \textbf{équation algébrique} que l'on sait résoudre : c'est là tout l'esprit de la géométrie analytique inventée par Descartes ;
\item les \textbf{équations paramétriques} décrivent un cercle non plus comme un lieu statique, mais comme un \textbf{mouvement} commandé par un paramètre, ce qui est exactement le langage des machines programmables (arrosage automatique, robots, trajectoires) ;
\item enfin, tu as pris l'habitude professionnelle de \textbf{convertir, comparer et conclure par écrit} : en mathématiques appliquées comme au Probatoire, un calcul sans phrase de conclusion argumentée reste un travail inachevé.
\end{itemize}

\begin{savaistu}[Les forages et la géométrie au Cameroun]
L'accès à l'eau potable est un enjeu majeur dans les régions du Nord et de l'Extrême-Nord du Cameroun, où des centaines de forages sont réalisés chaque année avec l'appui de partenaires techniques. Les ingénieurs du génie rural utilisent quotidiennement des plans cotés dans des repères orthonormés : distances de sécurité autour des points d'eau, tracés de canalisations perpendiculaires aux routes, zones de protection circulaires... Autant de situations où les outils de cette leçon — vecteur normal, distance point-droite, tangentes à un cercle — sont des instruments de travail réels, et non de simples exercices scolaires.
\end{savaistu}

\newpage

\coursec{Méthodes et exercices résolus}

\coursec{Géométrie analytique du plan}

\courssub{Donner un vecteur normal d'une droite}

\begin{methode}[Trouver un vecteur normal à une droite]
Pour déterminer un vecteur normal à une droite $(D)$, applique les réflexes suivants :
\begin{enumerate}
\item Si $(D)$ est donnée par une équation cartésienne $ax+by+c=0$, alors un vecteur normal est $\vec{n}(a\,;\,b)$. C'est immédiat : les coefficients de $x$ et $y$ sont les coordonnées de $\vec{n}$.
\item Si $(D)$ est donnée par une équation réduite $y=mx+p$, réécris-la sous la forme $mx-y+p=0$ : un vecteur normal est $\vec{n}(m\,;\,-1)$.
\item Si $(D)$ est donnée par un point $A$ et un vecteur directeur $\vec{u}(\alpha\,;\,\beta)$, alors tout vecteur $\vec{n}$ vérifiant $\vec{u}\cdot\vec{n}=0$ convient, par exemple $\vec{n}(-\beta\,;\,\alpha)$ ou $\vec{n}(\beta\,;\,-\alpha)$.
\item Si $(D)$ passe par deux points $A$ et $B$, calcule $\overrightarrow{AB}$ (vecteur directeur), puis applique le point 3.
\end{enumerate}
\textbf{Réflexe de vérification :} un vecteur normal n'est jamais unique : tous les vecteurs colinéaires à $\vec{n}$ non nuls conviennent. Vérifie toujours que $\vec{n}\neq\vec{0}$ et que $\vec{u}\cdot\vec{n}=0$.
\end{methode}

\begin{popfigure}
\begin{tikzpicture}[scale=0.8, >=Stealth]
\draw[popline, thick] (-1, -1) -- (5, 3) node[right] {$(D)$};
\draw[popblue, thick, ->] (2, 1) -- (3.2, 0.4) node[below right] {$\vec{n}(3\,;\,-4)$};
\draw[poporange, thick, ->] (2, 1) -- (3.6, 2.2) node[above right] {$\vec{u}(4\,;\,3)$};
\fill (2,1) circle (2pt) node[below left] {$A$};
\draw (2.5, 1.2) node[popdark, rotate=30] {$\perp$};
\end{tikzpicture}
\legende{Un vecteur normal $\vec{n}$ est orthogonal à tout vecteur directeur $\vec{u}$ de la droite.}
\end{popfigure}

\begin{exoresolu}[Vecteurs normaux dans différentes situations]
Dans le plan muni d'un repère orthonormé $(O\,;\vec{i},\vec{j})$, on considère les droites :
$(D_1) : 3x-4y+5=0$ ; \quad $(D_2) : y=2x-7$ ; \quad $(D_3)$ passant par $A(1\,;\,2)$ et dirigée par $\vec{u}(5\,;\,-3)$ ; \quad $(D_4)$ passant par $M(0\,;\,4)$ et $N(6\,;\,1)$.

Déterminer un vecteur normal à chacune de ces quatre droites.
\tcblower
\textbf{Droite $(D_1)$.} L'équation $3x-4y+5=0$ est de la forme $ax+by+c=0$ avec $a=3$ et $b=-4$. Un vecteur normal est donc :
\[\vec{n_1}(3\,;\,-4).\]

\textbf{Droite $(D_2)$.} On réécrit l'équation réduite sous forme cartésienne :
\[y=2x-7 \iff 2x-y-7=0.\]
On lit $a=2$, $b=-1$, d'où un vecteur normal :
\[\vec{n_2}(2\,;\,-1).\]

\textbf{Droite $(D_3)$.} Le vecteur $\vec{u}(5\,;\,-3)$ dirige $(D_3)$. On cherche $\vec{n_3}(a\,;\,b)$ tel que $\vec{u}\cdot\vec{n_3}=0$, c'est-à-dire $5a-3b=0$. En échangeant les coordonnées et en changeant l'un des signes, on obtient :
\[\vec{n_3}(3\,;\,5) \quad\text{car}\quad 5\times 3+(-3)\times 5=15-15=0.\]

\textbf{Droite $(D_4)$.} Un vecteur directeur est $\overrightarrow{MN}(6-0\,;\,1-4)$, soit $\overrightarrow{MN}(6\,;\,-3)$. On peut simplifier par $3$ : $\vec{u'}(2\,;\,-1)$ dirige aussi $(D_4)$. Un vecteur normal est alors :
\[\vec{n_4}(1\,;\,2) \quad\text{car}\quad 2\times 1+(-1)\times 2=0.\]

\textbf{Conclusion :} $\vec{n_1}(3\,;\,-4)$, $\vec{n_2}(2\,;\,-1)$, $\vec{n_3}(3\,;\,5)$ et $\vec{n_4}(1\,;\,2)$ sont des vecteurs normaux respectifs de $(D_1)$, $(D_2)$, $(D_3)$ et $(D_4)$.
\end{exoresolu}

\courssub{Écrire une équation normale d'une droite}

\begin{methode}[Équation normale d'une droite]
On appelle \cle{équation normale} d'une droite une équation cartésienne de la forme :
\[(\cos\theta)\,x+(\sin\theta)\,y-p=0 \quad \text{ou} \quad x\cos\theta+y\sin\theta=p\]
où $\vec{n}(\cos\theta\,;\,\sin\theta)$ est un vecteur normal \textbf{unitaire} ($\|\vec{n}\|=1$) et $p$ est la distance de l'origine $O$ à la droite ($p\geq 0$).

Pour obtenir l'équation normale d'une droite $(D) : ax+by+c=0$ :
\begin{enumerate}
\item Calcule $\sqrt{a^2+b^2}$ (norme du vecteur normal $\vec{n}(a\,;\,b)$).
\item Divise tous les coefficients de l'équation par $\sqrt{a^2+b^2}$.
\item Si le terme constant obtenu (côté gauche) est strictement positif, multiplie toute l'équation par $-1$ afin que le second membre $p$ soit positif ou nul.
\item Conclus : l'équation s'écrit $x\cos\theta+y\sin\theta=p$ avec $p=\dfrac{|c|}{\sqrt{a^2+b^2}}$.
\end{enumerate}
\textbf{À retenir :} dans une équation normale, les coefficients de $x$ et $y$ vérifient $\cos^2\theta+\sin^2\theta=1$, et $p=d(O\,;\,(D))$.
\end{methode}

\begin{exoresolu}[Mise sous forme normale]
Le plan est muni d'un repère orthonormé $(O\,;\vec{i},\vec{j})$. On considère la droite $(D) : 3x+4y-10=0$.
\begin{enumerate}
\item Déterminer une équation normale de $(D)$.
\item En déduire la distance de l'origine $O$ à la droite $(D)$.
\end{enumerate}
\tcblower
\textbf{1.} Un vecteur normal de $(D)$ est $\vec{n}(3\,;\,4)$. Sa norme est :
\[\|\vec{n}\|=\sqrt{3^2+4^2}=\sqrt{9+16}=\sqrt{25}=5.\]
On divise tous les coefficients de l'équation par $5$ :
\[\frac{3}{5}x+\frac{4}{5}y-2=0 \quad\iff\quad \frac{3}{5}x+\frac{4}{5}y=2.\]
Le second membre vaut $2>0$ : l'équation est bien sous forme normale. On a $\cos\theta=\dfrac{3}{5}$ et $\sin\theta=\dfrac{4}{5}$, et on vérifie :
\[\left(\frac{3}{5}\right)^2+\left(\frac{4}{5}\right)^2=\frac{9}{25}+\frac{16}{25}=\frac{25}{25}=1.\]
Une équation normale de $(D)$ est donc :
\[\frac{3}{5}x+\frac{4}{5}y-2=0 \quad \left(\text{ou } \frac{3}{5}x+\frac{4}{5}y=2\right).\]

\textbf{2.} Dans l'équation normale $x\cos\theta+y\sin\theta=p$, le nombre $p$ est la distance de $O$ à $(D)$. Ici $p=2$, donc :
\[d\big(O\,;\,(D)\big)=2.\]

\textbf{Conclusion :} l'équation normale de $(D)$ est $\dfrac{3}{5}x+\dfrac{4}{5}y=2$, et l'origine est à la distance $2$ de $(D)$.
\end{exoresolu}

\courssub{Calculer la distance d'un point à une droite}

\begin{popfigure}
\begin{tikzpicture}[scale=0.9, >=Stealth]
\draw[popline, thick] (-1, -0.5) -- (5, 3.5) node[right] {$(D) : ax+by+c=0$};
\fill[popblue] (1, 2.5) circle (2.5pt) node[above left] {$A(x_A,y_A)$};
\fill[poporange] (2.5, 1.5) circle (2.5pt) node[below right] {$H$};
\draw[popblue, thick, dashed] (1, 2.5) -- (2.5, 1.5);
\draw (2.2, 1.8) node[popdark] {$\perp$};
\draw[popmuted, <->] (1, 2.5) -- node[left] {$d(A,(D))$} (2.5, 1.5);
\draw[->, popdark] (0,0) -- (1, 2.5) node[midway, left] {$\overrightarrow{OA}$};
\end{tikzpicture}
\legende{La distance d'un point $A$ à une droite $(D)$ est la longueur du segment perpendiculaire $AH$.}
\end{popfigure}

\begin{methode}[Distance d'un point à une droite]
Dans un repère orthonormé, soient $A(x_A\,;\,y_A)$ un point et $(D) : ax+by+c=0$ une droite. La distance de $A$ à $(D)$ est :
\[d\big(A\,;\,(D)\big)=\frac{|ax_A+by_A+c|}{\sqrt{a^2+b^2}}.\]
Démarche :
\begin{enumerate}
\item Vérifie que l'équation de $(D)$ est bien sous la forme $ax+by+c=0$ (tout dans le membre de gauche, $0$ à droite). Sinon, transpose d'abord.
\item Identifie $a$, $b$, $c$ et les coordonnées de $A$.
\item Calcule le numérateur $|ax_A+by_A+c|$ puis le dénominateur $\sqrt{a^2+b^2}$.
\item Simplifie le résultat (dénominateur rationnel si possible).
\end{enumerate}
\textbf{Réflexes :} si $A\in(D)$, le numérateur est nul et la distance vaut $0$ ; la distance est toujours un réel \textbf{positif ou nul} ; elle s'exprime en unités de longueur du repère.
\end{methode}

\begin{exoresolu}[Distance d'un point à une droite — contexte]
Pour électrifier un quartier de Bafoussam, on modélise le tracé d'une ligne électrique par la droite $(D) : 4x-3y+10=0$ dans un repère orthonormé où l'unité est le hectomètre. La maison de la famille Kamdem est représentée par le point $K(5\,;\,2)$.
\begin{enumerate}
\item Calculer la distance de la maison $K$ à la ligne $(D)$.
\item La réglementation exige une distance minimale de $1$ hectomètre. La norme est-elle respectée ?
\end{enumerate}
\tcblower
\textbf{1.} L'équation $4x-3y+10=0$ est bien sous la forme $ax+by+c=0$ avec $a=4$, $b=-3$, $c=10$. On applique la formule avec $x_K=5$ et $y_K=2$ :
\begin{align*}
d\big(K\,;\,(D)\big) &= \frac{|4\times 5+(-3)\times 2+10|}{\sqrt{4^2+(-3)^2}}\\
&= \frac{|20-6+10|}{\sqrt{16+9}}\\
&= \frac{|24|}{\sqrt{25}}\\
&= \frac{24}{5}=4{,}8.
\end{align*}
La distance de la maison $K$ à la droite $(D)$ est donc $4{,}8$ unités, soit $4{,}8$ hectomètres, c'est-à-dire $480$ mètres.

\textbf{2.} On compare : $4{,}8>1$. La distance réelle ($480$ m) est largement supérieure au minimum réglementaire ($100$ m).

\textbf{Conclusion :} $d\big(K\,;\,(D)\big)=\dfrac{24}{5}=4{,}8$ hectomètres ; la norme de sécurité est respectée.
\end{exoresolu}

\courssub{Déterminer le centre, le rayon et une équation cartésienne d'un cercle}

\begin{methode}[Centre, rayon, équation cartésienne d'un cercle]
\textbf{De l'équation vers les éléments.} Toute équation de la forme $x^2+y^2+ax+by+c=0$ définit un cercle (ou un point, ou l'ensemble vide). Pour le savoir :
\begin{enumerate}
\item Regroupe les termes en $x$ et en $y$, puis complète les carrés :
\[x^2+ax=\left(x+\frac{a}{2}\right)^2-\frac{a^2}{4} \qquad\text{et}\qquad y^2+by=\left(y+\frac{b}{2}\right)^2-\frac{b^2}{4}.\]
\item Écris l'équation sous la forme $(x-x_\Omega)^2+(y-y_\Omega)^2=R'$ avec $R'=\dfrac{a^2}{4}+\dfrac{b^2}{4}-c$.
\item Conclus selon le signe de $R'$ :
\begin{itemize}
\item si $R'>0$ : cercle de centre $\Omega\left(-\dfrac{a}{2}\,;\,-\dfrac{b}{2}\right)$ et de rayon $R=\sqrt{R'}$ ;
\item si $R'=0$ : ensemble réduit au point $\Omega$ ;
\item si $R'<0$ : ensemble vide.
\end{itemize}
\end{enumerate}
\textbf{Des éléments vers l'équation.} Le cercle de centre $\Omega(x_\Omega\,;\,y_\Omega)$ et de rayon $R$ a pour équation cartésienne :
\[(x-x_\Omega)^2+(y-y_\Omega)^2=R^2 \quad\iff\quad x^2+y^2-2x_\Omega x-2y_\Omega y+x_\Omega^2+y_\Omega^2-R^2=0.\]
\textbf{Réflexe :} un point $M(x\,;\,y)$ appartient au cercle si et seulement si $\Omega M=R$, c'est-à-dire $(x-x_\Omega)^2+(y-y_\Omega)^2=R^2$.
\end{methode}

\begin{exoresolu}[Nature d'un ensemble et éléments caractéristiques]
Le plan est muni d'un repère orthonormé $(O\,;\vec{i},\vec{j})$. On considère l'ensemble $(\mathcal{C})$ des points $M(x\,;\,y)$ tels que :
\[x^2+y^2-4x+6y-12=0.\]
\begin{enumerate}
\item Démontrer que $(\mathcal{C})$ est un cercle dont on précisera le centre $\Omega$ et le rayon $R$.
\item Le point $A(5\,;\,-1)$ appartient-il à $(\mathcal{C})$ ?
\end{enumerate}
\tcblower
\textbf{1.} On complète les carrés. Pour les termes en $x$ : $x^2-4x=(x-2)^2-4$. Pour les termes en $y$ : $y^2+6y=(y+3)^2-9$. L'équation devient :
\begin{align*}
(x-2)^2-4+(y+3)^2-9-12 &= 0\\
(x-2)^2+(y+3)^2 &= 4+9+12\\
(x-2)^2+(y+3)^2 &= 25.
\end{align*}
On reconnaît la forme $(x-x_\Omega)^2+(y-y_\Omega)^2=R^2$ avec $R^2=25>0$.

\textbf{Conclusion :} $(\mathcal{C})$ est le cercle de centre $\Omega(2\,;\,-3)$ et de rayon $R=\sqrt{25}=5$.

\textbf{2.} Le point $A(5\,;\,-1)$ appartient à $(\mathcal{C})$ si et seulement si $\Omega A^2=25$. Calculons :
\[\Omega A^2=(5-2)^2+(-1-(-3))^2=3^2+2^2=9+4=13.\]
Or $13\neq 25$, donc $\Omega A=\sqrt{13}\neq 5$.

\textbf{Conclusion :} le point $A$ n'appartient pas au cercle $(\mathcal{C})$ ; comme $\Omega A=\sqrt{13}<5$, $A$ est même intérieur au disque.
\end{exoresolu}

\courssub{Déterminer des équations paramétriques d'un cercle}

\begin{popfigure}
\begin{tikzpicture}[scale=0.8, >=Stealth]
\draw[popline, thick] (0,0) circle (2cm);
\fill (0,0) circle (2pt) node[below left] {$\Omega(x_\Omega,y_\Omega)$};
\fill (1.6, 1.2) circle (2pt) node[above right] {$M(x,y)$};
\draw[popblue, thick, ->] (0,0) -- (1.6, 1.2) node[midway, above left] {$R$};
\draw[poporange, thick, ->] (0.5, 0) arc (0:36.87:0.5) node[midway, right] {$t$};
\draw[popmuted, dashed] (1.6, 0) -- (1.6, 1.2);
\draw[popmuted, dashed] (0, 1.2) -- (1.6, 1.2);
\node at (1.6, -0.3) {$x_\Omega+R\cos t$};
\node at (-0.3, 1.2) {$y_\Omega+R\sin t$};
\end{tikzpicture}
\legende{Repérage paramétrique d'un point $M$ sur le cercle : $x=x_\Omega+R\cos t$, $y=y_\Omega+R\sin t$.}
\end{popfigure}

\begin{methode}[Équations paramétriques d'un cercle]
Le cercle $(\mathcal{C})$ de centre $\Omega(x_\Omega\,;\,y_\Omega)$ et de rayon $R$ admet pour \cle{représentation paramétrique} :
\[\begin{cases} x=x_\Omega+R\cos t \\ y=y_\Omega+R\sin t \end{cases}, \quad t\in\mathbb{R}\ \text{(ou } t\in[0\,;\,2\pi[\text{)}.\]
Démarche :
\begin{enumerate}
\item Identifie le centre $\Omega$ et le rayon $R$ (à partir de l'équation cartésienne, complète les carrés si nécessaire).
\item Reporte dans le système ci-dessus.
\item Pour vérifier qu'un point $M_0(x_0\,;\,y_0)$ est sur le cercle et trouver son paramètre $t$, résous $\cos t=\dfrac{x_0-x_\Omega}{R}$ et $\sin t=\dfrac{y_0-y_\Omega}{R}$.
\end{enumerate}
\textbf{Réflexe de vérification :} en éliminant $t$ grâce à $\cos^2 t+\sin^2 t=1$, on doit retrouver l'équation cartésienne $(x-x_\Omega)^2+(y-y_\Omega)^2=R^2$.
\end{methode}

\begin{exoresolu}[Paramétrage d'un cercle]
Dans un repère orthonormé $(O\,;\vec{i},\vec{j})$, on considère le cercle $(\mathcal{C})$ d'équation :
\[x^2+y^2+2x-8y+8=0.\]
\begin{enumerate}
\item Déterminer le centre et le rayon de $(\mathcal{C})$, puis donner une représentation paramétrique de $(\mathcal{C})$.
\item Déterminer une valeur du paramètre $t$ correspondant au point $B(1\,;\,4)$.
\end{enumerate}
\tcblower
\textbf{1.} Complétons les carrés : $x^2+2x=(x+1)^2-1$ et $y^2-8y=(y-4)^2-16$. L'équation devient :
\[(x+1)^2-1+(y-4)^2-16+8=0 \iff (x+1)^2+(y-4)^2=9.\]
Donc $(\mathcal{C})$ est le cercle de centre $\Omega(-1\,;\,4)$ et de rayon $R=\sqrt{9}=3$. Une représentation paramétrique est :
\[\begin{cases} x=-1+3\cos t \\ y=4+3\sin t \end{cases}, \quad t\in[0\,;\,2\pi[.\]

\textbf{2.} Vérifions d'abord que $B\in(\mathcal{C})$ :
\[\Omega B^2=(1-(-1))^2+(4-4)^2=2^2+0=4\neq 9.\]
Le point $B$ ne vérifie pas $\Omega B=3$ : \textbf{$B$ n'appartient pas au cercle $(\mathcal{C})$}, il n'existe donc aucun paramètre $t$ lui correspondant.

Prenons plutôt le point $B'(2\,;\,4)$ : $\Omega B'^2=(2+1)^2+0=9$, donc $B'\in(\mathcal{C})$. On résout :
\[\cos t=\frac{2-(-1)}{3}=1 \quad\text{et}\quad \sin t=\frac{4-4}{3}=0.\]
L'unique solution dans $[0\,;\,2\pi[$ est $t=0$.

\textbf{Conclusion :} $(\mathcal{C})$ a pour centre $\Omega(-1\,;\,4)$, pour rayon $3$, et pour équations paramétriques $x=-1+3\cos t$, $y=4+3\sin t$. Le point $B(1\,;\,4)$ n'est pas sur le cercle ; le point $B'(2\,;\,4)$ correspond au paramètre $t=0$.
\end{exoresolu}

\courssub{Déterminer une équation cartésienne d'une tangente en un point d'un cercle}

\begin{popfigure}
\begin{tikzpicture}[scale=0.8, >=Stealth]
\draw[popline, thick] (0,0) circle (2cm);
\fill (0,0) circle (2pt) node[below left] {$\Omega$};
\fill (1.6, 1.2) circle (2pt) node[above right] {$A$};
\draw[popblue, thick, ->] (0,0) -- (1.6, 1.2) node[midway, above left] {$\overrightarrow{\Omega A}$};
\draw[poporange, thick] (-1, 2.8) -- (3.5, -0.8) node[right] {$(T)$};
\draw (1.3, 1.5) node[popdark] {$\perp$};
\draw[popmuted, dashed] (0,0) -- (1.6, 1.2);
\end{tikzpicture}
\legende{La tangente $(T)$ en $A$ est perpendiculaire au rayon $\Omega A$ ; $\overrightarrow{\Omega A}$ est un vecteur normal de $(T)$.}
\end{popfigure}

\begin{methode}[Tangente à un cercle en un de ses points]
Soient $(\mathcal{C})$ un cercle de centre $\Omega$ et $A(x_A\,;\,y_A)$ un point \textbf{du cercle}. La tangente $(T)$ à $(\mathcal{C})$ en $A$ est la droite passant par $A$ et \textbf{perpendiculaire au rayon} $(\Omega A)$ : le vecteur $\overrightarrow{\Omega A}$ est donc un \textbf{vecteur normal} de $(T)$.
\begin{enumerate}
\item Vérifie d'abord que $A\in(\mathcal{C})$, c'est-à-dire $\Omega A=R$. Sans cette condition, la méthode ne s'applique pas.
\item Calcule les coordonnées de $\overrightarrow{\Omega A}(x_A-x_\Omega\,;\,y_A-y_\Omega)$.
\item Écris que $M(x\,;\,y)\in(T) \iff \overrightarrow{\Omega A}\cdot\overrightarrow{AM}=0$ :
\[(x_A-x_\Omega)(x-x_A)+(y_A-y_\Omega)(y-y_A)=0.\]
\item Développe et réduis pour obtenir une équation cartésienne $ax+by+c=0$.
\end{enumerate}
\textbf{Formule de duplication (admise) :} pour le cercle $(x-x_\Omega)^2+(y-y_\Omega)^2=R^2$, la tangente en $A$ a pour équation :
\[(x_A-x_\Omega)(x-x_\Omega)+(y_A-y_\Omega)(y-y_\Omega)=R^2.\]
\end{methode}

\begin{exoresolu}[Tangente en un point d'un cercle]
Le plan est muni d'un repère orthonormé $(O\,;\vec{i},\vec{j})$. Soit $(\mathcal{C})$ le cercle d'équation :
\[x^2+y^2-2x+4y-20=0.\]
\begin{enumerate}
\item Déterminer le centre $\Omega$ et le rayon $R$ de $(\mathcal{C})$.
\item Vérifier que le point $A(4\,;\,2)$ appartient à $(\mathcal{C})$.
\item Déterminer une équation cartésienne de la tangente $(T)$ à $(\mathcal{C})$ en $A$.
\end{enumerate}
\tcblower
\textbf{1.} Complétons les carrés : $x^2-2x=(x-1)^2-1$ et $y^2+4y=(y+2)^2-4$. L'équation devient :
\[(x-1)^2+(y+2)^2=1+4+20=25.\]
Donc $\Omega(1\,;\,-2)$ et $R=5$.

\textbf{2.} Calculons $\Omega A^2$ :
\[\Omega A^2=(4-1)^2+(2-(-2))^2=3^2+4^2=9+16=25=R^2.\]
Donc $\Omega A=5=R$ : le point $A$ appartient bien au cercle $(\mathcal{C})$.

\textbf{3.} Le vecteur $\overrightarrow{\Omega A}(4-1\,;\,2-(-2))$, soit $\overrightarrow{\Omega A}(3\,;\,4)$, est normal à la tangente $(T)$. Un point $M(x\,;\,y)$ appartient à $(T)$ si et seulement si :
\[\overrightarrow{\Omega A}\cdot\overrightarrow{AM}=0 \iff 3(x-4)+4(y-2)=0.\]
Développons :
\[3x-12+4y-8=0 \iff 3x+4y-20=0.\]

\textbf{Vérification (facultative mais recommandée) :} la distance de $\Omega$ à $(T)$ doit valoir $R$ :
\[d\big(\Omega\,;\,(T)\big)=\frac{|3\times 1+4\times(-2)-20|}{\sqrt{3^2+4^2}}=\frac{|3-8-20|}{5}=\frac{25}{5}=5=R.\]
La tangente est bien à la distance $R$ du centre : le résultat est cohérent.

\textbf{Conclusion :} la tangente à $(\mathcal{C})$ en $A(4\,;\,2)$ a pour équation $3x+4y-20=0$.
\end{exoresolu}

\courssub{Déterminer les tangentes à un cercle passant par un point extérieur}

\begin{popfigure}
\begin{tikzpicture}[scale=0.8, >=Stealth]
\draw[popline, thick] (0,0) circle (1.5cm);
\fill (0,0) circle (2pt) node[below left] {$\Omega$};
\fill (4, 0) circle (2pt) node[below right] {$P$};
\draw[poporange, thick] (4,0) -- (0.56, 1.4) node[above left] {$T_1$};
\draw[poporange, thick] (4,0) -- (0.56, -1.4) node[below left] {$T_2$};
\draw[popmuted, dashed] (0,0) -- (0.56, 1.4);
\draw[popmuted, dashed] (0,0) -- (0.56, -1.4);
\draw (0.3, 0.7) node[popdark] {$\perp$};
\draw (0.3, -0.7) node[popdark] {$\perp$};
\draw[<->, popdark] (0, -1.8) -- (4, -1.8) node[midway, below] {$\Omega P > R$};
\end{tikzpicture}
\legende{Depuis un point $P$ extérieur au cercle ($\Omega P > R$), on peut tracer exactement deux tangentes.}
\end{popfigure}

\begin{methode}[Tangentes à un cercle issues d'un point extérieur]
Soient $(\mathcal{C})$ un cercle de centre $\Omega$ et de rayon $R$, et $P$ un point \textbf{extérieur} au cercle ($\Omega P>R$). Il existe exactement \textbf{deux} tangentes à $(\mathcal{C})$ passant par $P$.
\begin{enumerate}
\item Vérifie que $P$ est extérieur : calcule $\Omega P$ et compare à $R$.
\item Écris l'équation d'une droite $(\Delta)$ passant par $P$ : soit sous la forme $a(x-x_P)+b(y-y_P)=0$ avec $(a,b)\neq(0,0)$, soit en traitant séparément le cas d'une droite verticale $x=x_P$.
\item Traduis la tangence : $(\Delta)$ est tangente à $(\mathcal{C})$ si et seulement si
\[d\big(\Omega\,;\,(\Delta)\big)=R.\]
\item Résous l'équation obtenue (en général en $a$ et $b$, ou en la pente $m$) ; tu dois trouver \textbf{deux solutions}.
\item Écris les équations des deux tangentes et conclus.
\end{enumerate}
\textbf{Réflexes :} si tu ne trouves qu'une solution, c'est qu'une tangente verticale (pente infinie) t'a échappé — teste $x=x_P$. Si $\Omega P=R$, $P$ est sur le cercle (une seule tangente) ; si $\Omega P<R$, il n'y en a aucune.
\end{methode}

\begin{exoresolu}[Tangentes issues d'un point extérieur — type Probatoire]
Dans le plan muni d'un repère orthonormé $(O\,;\vec{i},\vec{j})$, on considère le cercle $(\mathcal{C})$ de centre $\Omega(1\,;\,2)$ et de rayon $R=2$, et le point $P(5\,;\,2)$.
\begin{enumerate}
\item Montrer que $P$ est extérieur au cercle $(\mathcal{C})$.
\item Déterminer les équations des tangentes à $(\mathcal{C})$ passant par $P$.
\end{enumerate}
\tcblower
\textbf{1.} Calculons la distance $\Omega P$ :
\[\Omega P=\sqrt{(5-1)^2+(2-2)^2}=\sqrt{16}=4.\]
Comme $\Omega P=4>2=R$, le point $P$ est extérieur au cercle : il existe exactement deux tangentes à $(\mathcal{C})$ passant par $P$.

\textbf{2.} \emph{Première famille : droites non verticales.} Une droite $(\Delta)$ passant par $P(5\,;\,2)$ et de pente $m$ a pour équation :
\[y-2=m(x-5) \iff mx-y+2-5m=0.\]
$(\Delta)$ est tangente à $(\mathcal{C})$ si et seulement si $d\big(\Omega\,;\,(\Delta)\big)=R$, c'est-à-dire :
\[\frac{|m\times 1-2+2-5m|}{\sqrt{m^2+1}}=2 \iff \frac{|-4m|}{\sqrt{m^2+1}}=2 \iff \frac{4|m|}{\sqrt{m^2+1}}=2.\]
En simplifiant par $2$ puis en élevant au carré (les deux membres sont positifs) :
\[\frac{4m^2}{m^2+1}=1 \iff 4m^2=m^2+1 \iff 3m^2=1 \iff m^2=\frac{1}{3}.\]
D'où deux solutions : $m=\dfrac{1}{\sqrt{3}}=\dfrac{\sqrt{3}}{3}$ ou $m=-\dfrac{\sqrt{3}}{3}$.

\emph{Cas de la droite verticale} $x=5$ : la distance de $\Omega(1\,;\,2)$ à cette droite est $|5-1|=4\neq 2$ : elle n'est pas tangente. Aucune solution n'a été perdue.

Les deux tangentes ont donc pour équations :
\[(T_1) : y-2=\frac{\sqrt{3}}{3}(x-5) \quad\text{et}\quad (T_2) : y-2=-\frac{\sqrt{3}}{3}(x-5).\]

\textbf{Vérification pour $(T_1)$} écrite $\dfrac{\sqrt{3}}{3}x-y+2-\dfrac{5\sqrt{3}}{3}=0$ :
\[d\big(\Omega\,;\,(T_1)\big)=\frac{\left|\dfrac{\sqrt{3}}{3}-2+2-\dfrac{5\sqrt{3}}{3}\right|}{\sqrt{\dfrac{1}{3}+1}}=\frac{\left|-\dfrac{4\sqrt{3}}{3}\right|}{\sqrt{\dfrac{4}{3}}}=\frac{\dfrac{4\sqrt{3}}{3}}{\dfrac{2}{\sqrt{3}}}=\frac{4\sqrt{3}}{3}\times\frac{\sqrt{3}}{2}=\frac{4\times 3}{6}=2=R.\]
Le résultat est confirmé.

\textbf{Conclusion :} les deux tangentes à $(\mathcal{C})$ issues de $P$ sont :
\[(T_1) : \sqrt{3}\,x-3y+6-5\sqrt{3}=0 \qquad (T_2) : \sqrt{3}\,x+3y-6-5\sqrt{3}=0.\]
\end{exoresolu}

\begin{attention}[Les 5 erreurs qui coûtent des points au Probatoire]
\begin{enumerate}
\item \textbf{Oublier la valeur absolue dans la formule de la distance.} La distance $d\big(A\,;\,(D)\big)=\dfrac{|ax_A+by_A+c|}{\sqrt{a^2+b^2}}$ est toujours positive. Écrire $d=-\dfrac{24}{5}$ ou « oublier » le signe sans valeur absolue est une erreur de cours éliminatoire.
\item \textbf{Appliquer la formule de distance à une équation mal préparée.} Si la droite est donnée par $y=2x-7$, il faut d'abord écrire $2x-y-7=0$. Utiliser $a=2$, $b=0$, $c=-7$ (en oubliant le $-y$) fausse tout le calcul.
\item \textbf{Confondre vecteur normal et vecteur directeur.} Pour $ax+by+c=0$, le vecteur $(a\,;\,b)$ est \emph{normal}, le vecteur $(-b\,;\,a)$ est \emph{directeur}. Dans la tangente en un point $A$ d'un cercle, c'est $\overrightarrow{\Omega A}$ (le rayon) qui est le vecteur \emph{normal} de la tangente, pas son vecteur directeur.
\item \textbf{Se tromper de signe dans le centre du cercle.} L'équation $(x-2)^2+(y+3)^2=25$ donne le centre $\Omega(2\,;\,-3)$ et non $\Omega(-2\,;\,3)$ : on lit $x_\Omega$ dans $(x-x_\Omega)^2$. De même, ne pas oublier d'écrire $R=\sqrt{25}=5$ et non $R=25$.
\item \textbf{Oublier de vérifier les conditions d'application et de conclure.} Avant de chercher une tangente en $A$, vérifie que $A$ est sur le cercle ($\Omega A=R$) ; avant de chercher les tangentes issues de $P$, vérifie que $P$ est extérieur ($\Omega P>R$). Et n'oublie pas de tester la droite verticale $x=x_P$ : si tu ne trouves qu'une pente $m$, la deuxième tangente est probablement verticale. Termine toujours par une phrase de conclusion.
\end{enumerate}
\end{attention}

\newpage

\coursec{Exercices}

\courssub{Je vérifie mes connaissances}

\begin{exercice}[QCM : vecteur normal]
Le plan est muni d'un repère orthonormé $(O\,;\vec{i},\vec{j})$. Pour chaque question, une seule réponse est exacte ; indique-la en justifiant brièvement.
\begin{enumerate}
\item Un vecteur normal à la droite $(D) : 3x - 4y + 5 = 0$ est :
\begin{enumerate}
\item $\vec{u}\begin{pmatrix} 3 \\ -4 \end{pmatrix}$ \qquad
\item $\vec{u}\begin{pmatrix} 4 \\ 3 \end{pmatrix}$ \qquad
\item $\vec{u}\begin{pmatrix} -4 \\ 3 \end{pmatrix}$
\end{enumerate}
\item La distance du point $A(1\,;1)$ à la droite $(D) : x + y - 4 = 0$ est égale à :
\begin{enumerate}
\item $\sqrt{2}$ \qquad \item $2\sqrt{2}$ \qquad \item $\dfrac{2}{\sqrt{2}}$
\end{enumerate}
\item Le cercle d'équation $x^2 + y^2 - 2x + 4y - 4 = 0$ a pour centre et rayon :
\begin{enumerate}
\item $\Omega(1\,;-2)$, $R = 3$ \qquad \item $\Omega(-1\,;2)$, $R = 3$ \qquad \item $\Omega(1\,;-2)$, $R = 2$
\end{enumerate}
\item Une équation de la tangente au cercle de centre $O$ et de rayon $5$ au point $M_0(3\,;4)$ est :
\begin{enumerate}
\item $3x + 4y - 25 = 0$ \qquad \item $4x - 3y = 0$ \qquad \item $3x + 4y + 25 = 0$
\end{enumerate}
\end{enumerate}
\end{exercice}

\begin{exercice}[Vrai ou faux justifié]
Le plan est muni d'un repère orthonormé. Pour chaque affirmation, dire si elle est vraie ou fausse en justifiant la réponse.
\begin{enumerate}
\item Si $\vec{n}\begin{pmatrix} a \\ b \end{pmatrix}$ est un vecteur normal à $(D)$, alors $\vec{u}\begin{pmatrix} -b \\ a \end{pmatrix}$ est un vecteur directeur de $(D)$.
\item L'équation $x^2 + y^2 + 2x - 2y + 5 = 0$ est l'équation d'un cercle de rayon $1$.
\item La distance d'un point $A$ à une droite $(D)$ est la distance de $A$ à un point quelconque de $(D)$.
\item Par un point extérieur à un cercle, il passe exactement deux tangentes à ce cercle.
\end{enumerate}
\end{exercice}

\begin{exercice}[Lecture directe sur les équations]
\begin{enumerate}
\item Donner un vecteur normal et un vecteur directeur de chacune des droites :
\[ (D_1) : 2x + 5y - 7 = 0 \qquad ; \qquad (D_2) : x = 3 \qquad ; \qquad (D_3) : y = -2x + 1. \]
\item Pour chaque équation, dire si elle définit un cercle ; si oui, préciser le centre et le rayon :
\[ x^2 + y^2 - 4x + 6y + 9 = 0 \qquad ; \qquad x^2 + y^2 + 2x + 6 = 0 \qquad ; \qquad x^2 + y^2 - 10x + 25 = 0. \]
\end{enumerate}
\end{exercice}

\begin{exercice}[Calculs directs de distances]
Le plan est muni d'un repère orthonormé.
\begin{enumerate}
\item Calculer la distance du point $A(2\,;-1)$ à la droite $(D) : 5x - 12y + 26 = 0$.
\item Calculer la distance du point $B(-3\,;4)$ à la droite $(\Delta) : 3x + 4y - 22 = 0$.
\item Déterminer la distance entre les droites parallèles $(D) : 2x - y + 1 = 0$ et $(D') : 2x - y - 4 = 0$.
\end{enumerate}
\end{exercice}

\courssub{Je m'entraîne}

\begin{exercice}[Équation normale d'une droite]
Le plan est muni d'un repère orthonormé $(O\,;\vec{i},\vec{j})$. On considère la droite $(D) : 5x - 12y + 26 = 0$ et le point $A(2\,;-1)$.
\begin{enumerate}
\item Donner un vecteur normal $\vec{n}$ à $(D)$ et vérifier que $\|\vec{n}\| = 13$.
\item En déduire une équation normale de $(D)$, c'est-à-dire une équation de la forme $\cos\theta\, x + \sin\theta\, y + c = 0$.
\item Calculer la distance de $A$ à $(D)$ en utilisant l'équation normale, puis vérifier le résultat avec la formule du cours.
\item Déterminer une équation de la droite $(D')$ passant par $A$ et perpendiculaire à $(D)$.
\end{enumerate}
\end{exercice}

\begin{exercice}[Aire d'un triangle]
Le plan est muni d'un repère orthonormé. On considère les points $A(1\,;2)$, $B(5\,;0)$ et $C(3\,;6)$.
\begin{enumerate}
\item Déterminer une équation cartésienne de la droite $(BC)$.
\item Calculer la distance du point $A$ à la droite $(BC)$.
\item Calculer la longueur $BC$.
\item En déduire l'aire du triangle $ABC$.
\end{enumerate}
\end{exercice}

\begin{exercice}[Cercle et droite : position relative]
Le plan est muni d'un repère orthonormé.
\begin{enumerate}
\item Montrer que l'équation $x^2 + y^2 - 6x + 2y + 6 = 0$ est celle d'un cercle $(\mathscr{C})$ dont on précisera le centre $\Omega$ et le rayon $R$.
\item Donner des équations paramétriques de $(\mathscr{C})$.
\item On considère la droite $(D) : x + y - 1 = 0$. Calculer la distance de $\Omega$ à $(D)$ et en déduire la position relative de $(D)$ et $(\mathscr{C})$.
\item Déterminer les coordonnées des points d'intersection éventuels de $(D)$ et $(\mathscr{C})$.
\end{enumerate}
\end{exercice}

\begin{exercice}[Tangente en un point d'un cercle]
Le plan est muni d'un repère orthonormé. On considère le cercle $(\mathscr{C})$ d'équation $(x - 1)^2 + (y + 2)^2 = 25$ et le point $M_0(4\,;2)$.
\begin{enumerate}
\item Vérifier que $M_0$ appartient à $(\mathscr{C})$.
\item Préciser le centre $\Omega$ de $(\mathscr{C})$ et déterminer les coordonnées du vecteur $\overrightarrow{\Omega M_0}$.
\item En utilisant le fait que la tangente en $M_0$ est perpendiculaire au rayon $[\Omega M_0]$, déterminer une équation cartésienne de la tangente $(T)$ à $(\mathscr{C})$ en $M_0$.
\item La droite $(T)$ coupe-t-elle l'axe des ordonnées ? Si oui, en quel point ?
\end{enumerate}
\end{exercice}

\courssub{Je me prépare au Probatoire}

\begin{exercice}[Type Probatoire : tangentes issues d'un point extérieur]
Le plan est muni d'un repère orthonormé $(O\,;\vec{i},\vec{j})$, unité graphique : \SI{1}{cm}. On considère le cercle $(\mathscr{C})$ d'équation $x^2 + y^2 = 9$ et le point $A(5\,;5)$.

\textbf{Partie A : position du point}
\begin{enumerate}
\item Préciser le centre et le rayon de $(\mathscr{C})$.
\item Calculer la distance $OA$ et en déduire que $A$ est extérieur à $(\mathscr{C})$.
\end{enumerate}

\textbf{Partie B : recherche des tangentes}

Soit $(\Delta)$ une droite passant par $A$, de coefficient directeur $m$, d'équation $y - 5 = m(x - 5)$.
\begin{enumerate}
\item Montrer que $(\Delta)$ est tangente à $(\mathscr{C})$ si et seulement si :
\[ \dfrac{|5 - 5m|}{\sqrt{1 + m^2}} = 3. \]
\item En déduire que $m$ est solution de l'équation $16m^2 - 50m + 16 = 0$, puis déterminer les deux valeurs de $m$.
\item Écrire les équations cartésiennes des deux tangentes à $(\mathscr{C})$ issues de $A$.
\end{enumerate}

\textbf{Partie C : longueur des segments tangents}

On note $T_1$ et $T_2$ les points de contact de ces tangentes avec $(\mathscr{C})$.
\begin{enumerate}
\item En utilisant le triangle $OAT_1$ rectangle en $T_1$, calculer la longueur $AT_1$.
\item Que peut-on dire des longueurs $AT_1$ et $AT_2$ ? Justifier.
\end{enumerate}
\end{exercice}

\begin{exercice}[Type Probatoire : droites à distance donnée d'un point]
Le plan est muni d'un repère orthonormé $(O\,;\vec{i},\vec{j})$. On considère les points $A(0\,;5)$ et $B(3\,;1)$.

\textbf{Partie A : recherche des droites}

On cherche les droites $(D)$ passant par $A$ et situées à la distance $2$ du point $B$.
\begin{enumerate}
\item Vérifier que la droite d'équation $x = 0$ ne convient pas.
\item Soit $(D)$ une droite passant par $A$, d'équation $y = mx + 5$. Montrer que $d(B\,;D) = \dfrac{|3m + 4|}{\sqrt{m^2 + 1}}$.
\item Résoudre l'équation $d(B\,;D) = 2$ et en déduire les équations des droites cherchées.
\end{enumerate}

\textbf{Partie B : interprétation géométrique}

On considère le cercle $(\mathscr{C})$ de centre $B$ et de rayon $2$.
\begin{enumerate}
\item Écrire une équation cartésienne de $(\mathscr{C})$.
\item Calculer la distance $AB$ et justifier que $A$ est extérieur à $(\mathscr{C})$.
\item Expliquer pourquoi les droites trouvées dans la partie A sont exactement les tangentes à $(\mathscr{C})$ passant par $A$.
\item Donner des équations paramétriques de $(\mathscr{C})$.
\end{enumerate}
\end{exercice}

\begin{exercice}[Type Probatoire : problème de synthèse en contexte]
Une entreprise de télécommunications veut installer une antenne relais dans une plaine. Dans un repère orthonormé $(O\,;\vec{i},\vec{j})$ où l'unité représente \SI{1}{km}, trois villages sont modélisés par les points $P(0\,;0)$, $Q(8\,;0)$ et $R(2\,;6)$. L'antenne doit être placée en un point $\Omega$ à égale distance des trois villages, et sa zone de couverture est le disque de centre $\Omega$ passant par $P$.

\textbf{Partie A : localisation de l'antenne}
\begin{enumerate}
\item Déterminer une équation de la médiatrice $(\Delta_1)$ du segment $[PQ]$.
\item Déterminer une équation de la médiatrice $(\Delta_2)$ du segment $[PR]$.
\item En déduire les coordonnées du point $\Omega$, centre du cercle circonscrit au triangle $PQR$.
\end{enumerate}

\textbf{Partie B : zone de couverture}
\begin{enumerate}
\item Calculer le rayon $R$ du cercle $(\mathscr{C})$ de centre $\Omega$ passant par $P$, puis écrire une équation cartésienne de $(\mathscr{C})$.
\item Donner des équations paramétriques de $(\mathscr{C})$.
\item Une route rectiligne est modélisée par la droite $(D) : x + y - 10 = 0$. Calculer la distance de $\Omega$ à $(D)$ et en déduire si la route traverse la zone de couverture.
\end{enumerate}

\textbf{Partie C : câble tangent}

On veut poser un câble en ligne droite, tangent au cercle $(\mathscr{C})$ au point $M_0(0\,;0)$.
\begin{enumerate}
\item Déterminer une équation cartésienne de la tangente $(T)$ à $(\mathscr{C})$ en $M_0$.
\item Le point $S(10\,;0)$, extrémité prévue du câble, appartient-il à $(T)$ ?
\end{enumerate}
\end{exercice}

\newpage

\coursec{Corrigés des exercices}

\begin{corrige}[Exercice 1 — QCM : vecteur normal]
\begin{enumerate}
\item \textbf{Réponse a.} Pour une droite d'équation $ax+by+c=0$, le vecteur $\vec{n}\begin{pmatrix} a \\ b \end{pmatrix}$ est normal. Ici $a=3$, $b=-4$, donc $\vec{n}\begin{pmatrix} 3 \\ -4 \end{pmatrix}$ convient. Vérification par produit scalaire : un vecteur directeur de $(D)$ est $\vec{u}\begin{pmatrix} 4 \\ 3 \end{pmatrix}$ et $3\times 4 + (-4)\times 3 = 0$.
\item \textbf{Réponse c.} $d(A\,;D)=\dfrac{|1+1-4|}{\sqrt{1^2+1^2}}=\dfrac{2}{\sqrt{2}}$. (Remarque : $\dfrac{2}{\sqrt{2}}=\sqrt{2}$, donc la valeur numérique coïncide avec la réponse a ; la réponse c est la forme issue directement de la formule.)
\item \textbf{Réponse a.} On a $a=-2$, $b=4$, $c=-4$ : centre $\Omega\left(-\dfrac{a}{2}\,;-\dfrac{b}{2}\right)=(1\,;-2)$ et $R^2=1^2+(-2)^2-(-4)=1+4+4=9$, donc $R=3$.
\item \textbf{Réponse a.} La tangente en $M_0(3\,;4)$ admet $\overrightarrow{OM_0}\begin{pmatrix} 3 \\ 4 \end{pmatrix}$ comme vecteur normal : $3(x-3)+4(y-4)=0$, soit $3x+4y-25=0$. On vérifie que $M_0$ appartient bien au cercle : $3^2+4^2=25=5^2$.
\end{enumerate}
\end{corrige}

\begin{corrige}[Exercice 2 — Vrai ou faux justifié]
\begin{enumerate}
\item \textbf{Vrai.} $\vec{n}\cdot\vec{u}=a\times(-b)+b\times a=0$, donc $\vec{u}$ est orthogonal à $\vec{n}$ : c'est un vecteur directeur de $(D)$.
\item \textbf{Faux.} Centre $(-1\,;1)$ et $R^2=(-1)^2+1^2-5=-3<0$ : l'équation ne représente aucun point, et a fortiori pas un cercle.
\item \textbf{Faux.} La distance de $A$ à $(D)$ est la distance de $A$ à son \emph{projeté orthogonal} $H$ sur $(D)$, c'est-à-dire le \emph{minimum} des distances $AM$ pour $M\in(D)$, et non la distance à un point quelconque.
\item \textbf{Vrai.} Si $A$ est extérieur au cercle de centre $\Omega$, les tangentes issues de $A$ correspondent aux points de contact $T$ tels que le triangle $\Omega AT$ soit rectangle en $T$ : il y en a exactement deux, symétriques par rapport à la droite $(\Omega A)$.
\end{enumerate}
\end{corrige}

\begin{corrige}[Exercice 3 — Lecture directe sur les équations]
\begin{enumerate}
\item $(D_1):2x+5y-7=0$ : vecteur normal $\vec{n}_1\begin{pmatrix} 2 \\ 5 \end{pmatrix}$, vecteur directeur $\vec{u}_1\begin{pmatrix} -5 \\ 2 \end{pmatrix}$.

$(D_2):x=3$, c'est-à-dire $x-3=0$ : vecteur normal $\vec{n}_2\begin{pmatrix} 1 \\ 0 \end{pmatrix}$ (droite verticale), vecteur directeur $\vec{u}_2\begin{pmatrix} 0 \\ 1 \end{pmatrix}$.

$(D_3):y=-2x+1$, c'est-à-dire $2x+y-1=0$ : vecteur normal $\vec{n}_3\begin{pmatrix} 2 \\ 1 \end{pmatrix}$, vecteur directeur $\vec{u}_3\begin{pmatrix} 1 \\ -2 \end{pmatrix}$.
\item $x^2+y^2-4x+6y+9=0$ : centre $(2\,;-3)$, $R^2=4+9-9=4$, donc $R=2$ : \textbf{cercle} de centre $(2\,;-3)$ et de rayon $2$.

$x^2+y^2+2x+6=0$ : centre $(-1\,;0)$, $R^2=1+0-6=-5<0$ : \textbf{pas un cercle} (ensemble vide).

$x^2+y^2-10x+25=0$ : centre $(5\,;0)$, $R^2=25+0-25=0$ : $R=0$, l'équation se réduit à $(x-5)^2+y^2=0$, vérifiée uniquement par le point $(5\,;0)$ : \textbf{ce n'est pas un cercle} (cercle de rayon nul, réduit à un point).
\end{enumerate}
\end{corrige}

\begin{corrige}[Exercice 4 — Calculs directs de distances]
\begin{enumerate}
\item $d(A\,;D)=\dfrac{|5\times 2-12\times(-1)+26|}{\sqrt{5^2+(-12)^2}}=\dfrac{|10+12+26|}{\sqrt{25+144}}=\dfrac{48}{13}$.
\item $d(B\,;\Delta)=\dfrac{|3\times(-3)+4\times 4-22|}{\sqrt{9+16}}=\dfrac{|-9+16-22|}{5}=\dfrac{15}{5}=3$.
\item Les deux droites sont parallèles (mêmes coefficients $2$ et $-1$). On choisit un point de $(D)$, par exemple $M(0\,;1)$ (car $2\times 0-1+1=0$). Alors :
\[ d(D\,;D')=d(M\,;D')=\dfrac{|2\times 0-1-4|}{\sqrt{4+1}}=\dfrac{5}{\sqrt{5}}=\sqrt{5}. \]
\end{enumerate}
\end{corrige}

\begin{corrige}[Exercice 5 — Équation normale d'une droite]
\begin{enumerate}
\item $\vec{n}\begin{pmatrix} 5 \\ -12 \end{pmatrix}$ est normal à $(D)$ et $\|\vec{n}\|=\sqrt{25+144}=\sqrt{169}=13$.
\item On divise l'équation par $\|\vec{n}\|=13$ : $\dfrac{5}{13}x-\dfrac{12}{13}y+2=0$. Comme $\left(\dfrac{5}{13}\right)^2+\left(\dfrac{12}{13}\right)^2=1$, il existe $\theta$ tel que $\cos\theta=\dfrac{5}{13}$ et $\sin\theta=-\dfrac{12}{13}$ : c'est bien une équation normale.
\item Avec l'équation normale : $d(A\,;D)=\left|\dfrac{5}{13}\times 2-\dfrac{12}{13}\times(-1)+2\right|=\left|\dfrac{10+12+26}{13}\right|=\dfrac{48}{13}$. La formule du cours donne le même résultat (exercice 4, question 1) : $d(A\,;D)=\dfrac{48}{13}\approx 3{,}69$.
\item $(D')$ est perpendiculaire à $(D)$, donc $\vec{n}$ est un \emph{vecteur directeur} de $(D')$ : un vecteur normal de $(D')$ est $\begin{pmatrix} 12 \\ 5 \end{pmatrix}$. D'où $(D'):12x+5y+c=0$. Comme $A(2\,;-1)\in(D')$ : $24-5+c=0$, soit $c=-19$. Finalement :
\[ (D'):12x+5y-19=0. \]
\end{enumerate}
\end{corrige}

\begin{corrige}[Exercice 6 — Aire d'un triangle]
\begin{enumerate}
\item $\overrightarrow{BC}\begin{pmatrix} 3-5 \\ 6-0 \end{pmatrix}=\begin{pmatrix} -2 \\ 6 \end{pmatrix}$, que l'on simplifie en $\begin{pmatrix} 1 \\ -3 \end{pmatrix}$. Un vecteur normal est $\begin{pmatrix} 3 \\ 1 \end{pmatrix}$, d'où $(BC):3x+y+c=0$. Avec $B(5\,;0)$ : $15+c=0$, donc $c=-15$ :
\[ (BC):3x+y-15=0. \]
Vérification avec $C$ : $9+6-15=0$. \checkmark
\item $d(A\,;BC)=\dfrac{|3\times 1+2-15|}{\sqrt{9+1}}=\dfrac{10}{\sqrt{10}}=\sqrt{10}$.
\item $BC=\sqrt{(-2)^2+6^2}=\sqrt{4+36}=\sqrt{40}=2\sqrt{10}$.
\item $\mathscr{A}=\dfrac{1}{2}\times BC\times d(A\,;BC)=\dfrac{1}{2}\times 2\sqrt{10}\times\sqrt{10}=\dfrac{1}{2}\times 20=\boxed{10}$ unités d'aire.
\end{enumerate}
\end{corrige}

\begin{corrige}[Exercice 7 — Cercle et droite : position relative]
\begin{enumerate}
\item Équation de la forme $x^2+y^2+ax+by+c=0$ avec $a=-6$, $b=2$, $c=6$. Centre $\Omega\left(3\,;-1\right)$ et $R^2=3^2+(-1)^2-6=9+1-6=4$, donc $R=2$. Forme canonique : $(x-3)^2+(y+1)^2=4$. C'est bien un cercle.
\item Équations paramétriques :
\[ \begin{cases} x=3+2\cos t \\ y=-1+2\sin t \end{cases} \quad t\in\mathbb{R}. \]
\item $d(\Omega\,;D)=\dfrac{|3+(-1)-1|}{\sqrt{1+1}}=\dfrac{1}{\sqrt{2}}=\dfrac{\sqrt{2}}{2}\approx 0{,}71$. Comme $d(\Omega\,;D)<R=2$, la droite $(D)$ est \textbf{sécante} au cercle : deux points d'intersection.
\item Sur $(D)$ : $y=1-x$. On substitue dans l'équation du cercle :
\begin{align*}
x^2+(1-x)^2-6x+2(1-x)+6&=0\\
2x^2-10x+9&=0.
\end{align*}
Discriminant : $\Delta=100-72=28>0$, d'où $x=\dfrac{10\pm 2\sqrt{7}}{4}=\dfrac{5\pm\sqrt{7}}{2}$ et $y=1-x=\dfrac{-3\mp\sqrt{7}}{2}$. Les deux points d'intersection sont :
\[ I_1\left(\dfrac{5+\sqrt{7}}{2}\,;\dfrac{-3-\sqrt{7}}{2}\right) \qquad I_2\left(\dfrac{5-\sqrt{7}}{2}\,;\dfrac{-3+\sqrt{7}}{2}\right). \]
\end{enumerate}
\end{corrige}

\begin{corrige}[Exercice 8 — Tangente en un point d'un cercle]
\begin{enumerate}
\item $(4-1)^2+(2+2)^2=9+16=25$ : $M_0$ appartient bien à $(\mathscr{C})$.
\item Centre $\Omega(1\,;-2)$, rayon $5$. $\overrightarrow{\Omega M_0}\begin{pmatrix} 4-1 \\ 2-(-2) \end{pmatrix}=\begin{pmatrix} 3 \\ 4 \end{pmatrix}$.
\item La tangente $(T)$ en $M_0$ est perpendiculaire à $(\Omega M_0)$, donc $\overrightarrow{\Omega M_0}$ est un vecteur normal de $(T)$. Pour $M(x\,;y)\in(T)$ : $\overrightarrow{\Omega M_0}\cdot\overrightarrow{M_0M}=0$, soit :
\[ 3(x-4)+4(y-2)=0 \iff (T):3x+4y-20=0. \]
\item Pour $x=0$ : $4y-20=0$, donc $y=5$. $(T)$ coupe l'axe des ordonnées au point $(0\,;5)$.
\end{enumerate}
\end{corrige}

\begin{corrige}[Exercice 9 — Type Probatoire : tangentes issues d'un point extérieur]
\textbf{Partie A}
\begin{enumerate}
\item $(\mathscr{C}):x^2+y^2=9$ : centre $O(0\,;0)$, rayon $R=3$.
\item $OA=\sqrt{5^2+5^2}=\sqrt{50}=5\sqrt{2}\approx 7{,}07$. Comme $OA>R$ ($5\sqrt{2}>3$ car $50>9$), le point $A$ est extérieur à $(\mathscr{C})$.
\end{enumerate}

\textbf{Partie B}
\begin{enumerate}
\item $(\Delta):y-5=m(x-5)$ s'écrit $mx-y+5-5m=0$. $(\Delta)$ est tangente à $(\mathscr{C})$ si et seulement si la distance du centre $O$ à $(\Delta)$ est égale au rayon :
\[ d(O\,;\Delta)=\dfrac{|m\times 0-0+5-5m|}{\sqrt{m^2+1}}=\dfrac{|5-5m|}{\sqrt{1+m^2}}=3. \]
\item En élevant au carré (les deux membres sont positifs) :
\begin{align*}
(5-5m)^2=9(1+m^2) &\iff 25-50m+25m^2=9+9m^2\\
&\iff 16m^2-50m+16=0.
\end{align*}
Discriminant : $\Delta=(-50)^2-4\times 16\times 16=2500-1024=1476=36\times 41$, donc $\sqrt{\Delta}=6\sqrt{41}$. D'où :
\[ m_1=\dfrac{50-6\sqrt{41}}{32}=\dfrac{25-3\sqrt{41}}{16} \qquad m_2=\dfrac{25+3\sqrt{41}}{16}. \]
(Numériquement : $m_1\approx 0{,}36$ et $m_2\approx 2{,}76$.)
\item Les deux tangentes ont pour équations :
\[ (T_1):y-5=\dfrac{25-3\sqrt{41}}{16}(x-5) \qquad (T_2):y-5=\dfrac{25+3\sqrt{41}}{16}(x-5), \]
soit sous forme cartésienne $m_i x-y+5-5m_i=0$ pour $i\in\{1,2\}$.
\end{enumerate}

\textbf{Partie C}
\begin{enumerate}
\item Le triangle $OAT_1$ est rectangle en $T_1$ (la tangente est perpendiculaire au rayon au point de contact). Par le théorème de Pythagore :
\[ AT_1=\sqrt{OA^2-OT_1^2}=\sqrt{50-9}=\sqrt{41}\ \text{cm}\approx 6{,}4\ \text{cm}. \]
\item $AT_1=AT_2=\sqrt{41}$ : les triangles $OAT_1$ et $OAT_2$ sont rectangles, ont la même hypoténuse $OA$ et des côtés $OT_1=OT_2=R$ ; ils sont donc isométriques. Les segments tangents issus d'un même point extérieur ont la même longueur.
\end{enumerate}

\textbf{Barème indicatif (sur 10)} : Partie A : 2 pts (centre-rayon 1 pt ; $OA$ et conclusion 1 pt) — Partie B : 5 pts (mise en équation de la distance 1,5 pt ; équation du second degré 1,5 pt ; résolution 1 pt ; équations des tangentes 1 pt) — Partie C : 3 pts (Pythagore 1,5 pt ; égalité justifiée 1,5 pt).

\textbf{Points de vigilance} : citer la condition de tangence $d(O\,;\Delta)=R$ ; justifier l'élévation au carré (membres positifs) ; ne pas oublier de vérifier que $A$ est extérieur \emph{avant} d'affirmer l'existence de deux tangentes ; préciser que le triangle est rectangle \emph{en $T_1$} (propriété fondamentale de la tangente).
\end{corrige}

\begin{corrige}[Exercice 10 — Type Probatoire : droites à distance donnée d'un point]
\textbf{Partie A}
\begin{enumerate}
\item La droite $x=0$ est la droite verticale passant par $A$. $d(B\,;x=0)=|3|=3\neq 2$ : elle ne convient pas. Cela justifie que l'on cherche les autres droites sous la forme $y=mx+5$.
\item $(D):y=mx+5$ s'écrit $mx-y+5=0$. Alors :
\[ d(B\,;D)=\dfrac{|m\times 3-1+5|}{\sqrt{m^2+1}}=\dfrac{|3m+4|}{\sqrt{m^2+1}}. \]
\item $d(B\,;D)=2\iff (3m+4)^2=4(m^2+1)\iff 9m^2+24m+16=4m^2+4\iff 5m^2+24m+12=0$.

Discriminant : $\Delta=24^2-4\times 5\times 12=576-240=336=16\times 21$, donc $\sqrt{\Delta}=4\sqrt{21}$. D'où :
\[ m_1=\dfrac{-24-4\sqrt{21}}{10}=\dfrac{-12-2\sqrt{21}}{5} \qquad m_2=\dfrac{-12+2\sqrt{21}}{5}. \]
Les droites cherchées sont :
\[ (D_1):y=\dfrac{-12-2\sqrt{21}}{5}\,x+5 \qquad (D_2):y=\dfrac{-12+2\sqrt{21}}{5}\,x+5. \]
\end{enumerate}

\textbf{Partie B}
\begin{enumerate}
\item $(\mathscr{C}):(x-3)^2+(y-1)^2=4$.
\item $AB=\sqrt{(3-0)^2+(1-5)^2}=\sqrt{9+16}=\sqrt{25}=5$. Comme $AB=5>R=2$, $A$ est extérieur à $(\mathscr{C})$.
\item Une droite est tangente à $(\mathscr{C})$ si et seulement si la distance de son centre $B$ à la droite est égale au rayon $2$. Les droites de la partie A passent par $A$ et vérifient exactement $d(B\,;D)=2$ : ce sont donc les tangentes à $(\mathscr{C})$ issues de $A$. Comme $A$ est extérieur, on retrouve bien qu'il en existe exactement deux.
\item Équations paramétriques :
\[ \begin{cases} x=3+2\cos t \\ y=1+2\sin t \end{cases} \quad t\in\mathbb{R}. \]
\end{enumerate}

\textbf{Barème indicatif (sur 10)} : Partie A : 5 pts (cas $x=0$ : 1 pt ; formule de distance : 1,5 pt ; équation et résolution : 1,5 pt ; équations des droites : 1 pt) — Partie B : 5 pts (équation du cercle : 1 pt ; $AB$ et position : 1,5 pt ; justification tangence : 1,5 pt ; équations paramétriques : 1 pt).

\textbf{Points de vigilance} : traiter \emph{séparément} le cas de la droite verticale $x=0$ (elle n'a pas de coefficient directeur : l'oublier est une erreur classique) ; la condition « distance égale au rayon » doit être explicitée comme critère de tangence ; penser à justifier que $A$ est extérieur pour conclure à l'existence des deux tangentes.
\end{corrige}

\begin{corrige}[Exercice 11 — Type Probatoire : problème de synthèse en contexte]
\textbf{Partie A}
\begin{enumerate}
\item $P(0\,;0)$ et $Q(8\,;0)$ : le milieu de $[PQ]$ est $(4\,;0)$ et $[PQ]$ est horizontal, donc sa médiatrice est verticale :
\[ (\Delta_1):x=4. \]
\item Le milieu de $[PR]$ est $I(1\,;3)$ et $\overrightarrow{PR}\begin{pmatrix} 2 \\ 6 \end{pmatrix}$ est un vecteur normal de $(\Delta_2)$. Pour $M(x\,;y)\in(\Delta_2)$ : $\overrightarrow{PR}\cdot\overrightarrow{IM}=0$, soit $2(x-1)+6(y-3)=0$, d'où :
\[ (\Delta_2):2x+6y-20=0 \iff x+3y-10=0. \]
\item $\Omega$ est l'intersection de $(\Delta_1)$ et $(\Delta_2)$ : $x=4$ et $4+3y-10=0$ donnent $y=2$. Donc $\Omega(4\,;2)$. L'antenne sera installée à \SI{4}{km} à l'est et \SI{2}{km} au nord de $P$.
\end{enumerate}

\textbf{Partie B}
\begin{enumerate}
\item $R=\Omega P=\sqrt{(4-0)^2+(2-0)^2}=\sqrt{16+4}=\sqrt{20}=2\sqrt{5}\approx 4{,}47$ km. Équation cartésienne :
\[ (\mathscr{C}):(x-4)^2+(y-2)^2=20. \]
Vérification : $Q$ : $(8-4)^2+(0-2)^2=16+4=20$ \checkmark ; $R$ : $(2-4)^2+(6-2)^2=4+16=20$ \checkmark.
\item Équations paramétriques :
\[ \begin{cases} x=4+2\sqrt{5}\cos t \\ y=2+2\sqrt{5}\sin t \end{cases} \quad t\in\mathbb{R}. \]
\item $d(\Omega\,;D)=\dfrac{|4+2-10|}{\sqrt{1+1}}=\dfrac{4}{\sqrt{2}}=2\sqrt{2}\approx 2{,}83$ km. Comme $d(\Omega\,;D)=2\sqrt{2}<R=2\sqrt{5}$ (car $8<20$), la droite $(D)$ est sécante au cercle : \textbf{la route traverse la zone de couverture}.
\end{enumerate}

\textbf{Partie C}
\begin{enumerate}
\item $M_0(0\,;0)=P$ appartient bien à $(\mathscr{C})$. $\overrightarrow{\Omega M_0}\begin{pmatrix} 0-4 \\ 0-2 \end{pmatrix}=\begin{pmatrix} -4 \\ -2 \end{pmatrix}$ est normal à la tangente $(T)$. Pour $M(x\,;y)\in(T)$ : $\overrightarrow{\Omega M_0}\cdot\overrightarrow{M_0M}=0$, soit $-4x-2y=0$, d'où :
\[ (T):2x+y=0. \]
\item Pour $S(10\,;0)$ : $2\times 10+0=20\neq 0$, donc $S$ n'appartient pas à $(T)$. Le câble tangent en $M_0$ ne peut pas aboutir en $S$ ; il faudra revoir le tracé ou le point d'ancrage.
\end{enumerate}

\textbf{Barème indicatif (sur 10)} : Partie A : 3 pts (médiatrices : 1 pt chacune ; $\Omega$ : 1 pt) — Partie B : 4 pts (rayon et équation : 1,5 pt ; paramétriques : 1 pt ; distance et conclusion : 1,5 pt) — Partie C : 3 pts (tangente : 2 pts ; test de $S$ : 1 pt).

\textbf{Points de vigilance} : pour la médiatrice, exiger la caractérisation $MP=MQ$ ou l'orthogonalité avec le milieu ; pour la position droite-cercle, comparer $d(\Omega\,;D)$ et $R$ en justifiant l'inégalité (comparer les carrés si besoin) ; pour la tangente, vérifier d'abord que $M_0\in(\mathscr{C})$ et citer la propriété « la tangente est perpendiculaire au rayon au point de contact ».
\end{corrige}

\newpage

\coursec{Fiche bilan}

\begin{popfigure}
\begin{tikzpicture}[mindmap, grow cyclic, every node/.style={concept, text=white, font=\small\bfseries},
root concept/.append style={concept color=popdark, font=\bfseries},
level 1/.append style={level distance=3.4cm, sibling angle=51.4},
level 2/.append style={level distance=2.1cm, sibling angle=45, font=\scriptsize}]
\node[root concept]{Géométrie analytique du plan}
  child[concept color=popblue]{ node{Vecteur normal}
    child{ node{$\vec{n}(a;b)$ perpendiculaire à $(D)$} }
    child{ node{$ax+by+c=0$} }
  }
  child[concept color=poporange]{ node{Équation normale}
    child{ node{$\dfrac{ax+by+c}{\sqrt{a^2+b^2}}=0$} }
    child{ node{Coefficients unitaires} }
  }
  child[concept color=popgreen]{ node{Distance point-droite}
    child{ node{$d(M,D)=\dfrac{|ax_0+by_0+c|}{\sqrt{a^2+b^2}}$} }
  }
  child[concept color=poppurple]{ node{Cercle : équations}
    child{ node{$(x-a)^2+(y-b)^2=R^2$} }
    child{ node{$x^2+y^2+\alpha x+\beta y+\gamma=0$} }
  }
  child[concept color=poppink]{ node{Équations paramétriques}
    child{ node{$x=a+R\cos t$} }
    child{ node{$y=b+R\sin t$} }
  }
  child[concept color=popgold]{ node{Tangente en $M_0\in\mathcal{C}$}
    child{ node{$\overrightarrow{\Omega M_0}\cdot\overrightarrow{M_0M}=0$} }
  }
  child[concept color=popblue!70]{ node{Tangentes par un point extérieur}
    child{ node{$d(\Omega,\Delta)=R$} }
    child{ node{Deux tangentes} }
  };
\end{tikzpicture}
\legende{Carte mentale de la leçon : les sept idées forces de la géométrie analytique du plan.}
\end{popfigure}

\begin{bilan}
\textbf{Définitions clés}
\begin{itemize}
\item Un \cle{vecteur normal} à une droite $(D)$ est tout vecteur non nul orthogonal à la direction de $(D)$. Si $(D):ax+by+c=0$, alors $\vec{n}(a;b)$ est normal à $(D)$ et $\vec{u}(-b;a)$ dirige $(D)$.
\item Une \cle{équation normale} de $(D)$ est une équation de la forme $\dfrac{ax+by+c}{\sqrt{a^2+b^2}}=0$ : les coefficients de $x$ et $y$ forment un vecteur normal \emph{unitaire}.
\item Le \cle{cercle} de centre $\Omega(a;b)$ et de rayon $R>0$ est l'ensemble des points $M$ tels que $\Omega M=R$.
\end{itemize}

\textbf{Formules à connaître par c\oe ur}
\begin{center}
\begin{tabularx}{\linewidth}{|l|X|}\hline
\rowcolor{popblueL} \textbf{Notion} & \textbf{Formule} \\ \hline
Distance point-droite & $d(M_0,D)=\dfrac{|ax_0+by_0+c|}{\sqrt{a^2+b^2}}$ \\ \hline
Équation cartésienne du cercle & $(x-a)^2+(y-b)^2=R^2$ \\ \hline
Forme développée & $x^2+y^2-2ax-2by+a^2+b^2-R^2=0$ \\ \hline
Centre et rayon & $\Omega\left(-\dfrac{\alpha}{2};-\dfrac{\beta}{2}\right)$, \quad $R=\sqrt{\dfrac{\alpha^2+\beta^2}{4}-\gamma}$ (si positif) \\ \hline
Équations paramétriques & $\begin{cases} x=a+R\cos t \\ y=b+R\sin t \end{cases}$, $t\in\mathbb{R}$ \\ \hline
Tangente en $M_0$ & $\overrightarrow{\Omega M_0}\cdot\overrightarrow{M_0M}=0$ \\ \hline
Tangente par point extérieur & droites $\Delta$ par $M_0$ vérifiant $d(\Omega,\Delta)=R$ \\ \hline
\end{tabularx}
\end{center}

\textbf{Méthodes express}
\begin{itemize}
\item \textbf{Trouver le centre et le rayon} : regrouper les termes en $x$ et en $y$, compléter les carrés, identifier avec $(x-a)^2+(y-b)^2=R^2$.
\item \textbf{Tangente en $M_0(x_0;y_0)$ du cercle} : écrire $(x_0-a)(x-x_0)+(y_0-b)(y-y_0)=0$, puis développer.
\item \textbf{Tangentes par un point extérieur $M_0$} : écrire les droites passant par $M_0$ (pente $m$ : $y-y_0=m(x-x_0)$, sans oublier la droite verticale $x=x_0$), puis résoudre $d(\Omega,\Delta)=R$. On obtient en général \emph{deux} tangentes.
\item \textbf{Position d'un point par rapport au cercle} : comparer $\Omega M$ à $R$ (intérieur, sur le cercle, extérieur).
\end{itemize}
\end{bilan}

\begin{aretenir}[Checklist avant le Probatoire]
\begin{itemize}
\item[$\square$] Je sais lire un vecteur normal $\vec{n}(a;b)$ et un vecteur directeur $\vec{u}(-b;a)$ sur l'équation $ax+by+c=0$.
\item[$\square$] Je sais écrire l'équation normale d'une droite en divisant par $\sqrt{a^2+b^2}$.
\item[$\square$] Je connais par c\oe ur la formule $d(M_0,D)=\dfrac{|ax_0+by_0+c|}{\sqrt{a^2+b^2}}$ et je n'oublie pas la valeur absolue.
\item[$\square$] Je sais compléter les carrés pour passer de $x^2+y^2+\alpha x+\beta y+\gamma=0$ au centre et au rayon.
\item[$\square$] Je vérifie toujours que la quantité sous la racine est strictement positive avant de conclure à un cercle.
\item[$\square$] Je sais donner les équations paramétriques $x=a+R\cos t$, $y=b+R\sin t$ d'un cercle.
\item[$\square$] Je sais écrire la tangente en un point du cercle grâce à l'orthogonalité $\overrightarrow{\Omega M_0}\cdot\overrightarrow{M_0M}=0$.
\item[$\square$] Je vérifie d'abord que $M_0$ appartient bien au cercle avant d'appliquer la formule de la tangente en $M_0$.
\item[$\square$] Pour un point extérieur, je cherche les droites $\Delta$ telles que $d(\Omega,\Delta)=R$ et je pense à la droite verticale.
\item[$\square$] Je sais qu'un point extérieur donne deux tangentes, un point du cercle une seule, un point intérieur aucune.
\item[$\square$] Je rédige proprement : je cite le repère, je justifie l'orthogonalité, je conclus par une équation simplifiée.
\end{itemize}
\end{aretenir}

\findecours

\end{document}
